\documentclass[11pt,letterpaper]{article} \usepackage[utf8]{inputenc}
\usepackage[margin=1in]{geometry} \usepackage{amsmath} \usepackage{amssymb}
\usepackage{amsthm} \usepackage{mathtools} \usepackage{microtype}
\clubpenalty=10000
\widowpenalty=10000
\usepackage{graphicx} \usepackage{tikz} \usepackage{booktabs} \usepackage{array}
\usepackage{float} \usepackage{caption} \usepackage{subcaption}
\usepackage{xcolor} \usepackage{algorithm} \usepackage{algpseudocode}
\usepackage{aliascnt} \usepackage[numbers,sort&compress]{natbib}
\usepackage{hyperref} \usepackage[nameinlink,capitalise,noabbrev]{cleveref}
\usepackage{todonotes}
\makeatletter
\def\theHALG@line{\thealgorithm.\arabic{ALG@line}}
\makeatother
\hypersetup{hidelinks,pdftitle={NP-Hardness and Additive Inapproximability of Deterministic Communication Complexity}}
\usetikzlibrary{arrows.meta,positioning,calc,matrix,fit,backgrounds}
\theoremstyle{plain} \newtheorem{theorem}{Theorem}[section]
\newaliascnt{lemma}{theorem} \newtheorem{lemma}[lemma]{Lemma}
\aliascntresetthe{lemma} \newaliascnt{proposition}{theorem}
\newtheorem{proposition}[proposition]{Proposition}
\aliascntresetthe{proposition} \newaliascnt{corollary}{theorem}
\newtheorem{corollary}[corollary]{Corollary} \aliascntresetthe{corollary}
\newaliascnt{claim}{theorem} \newtheorem{claim}[claim]{Claim}
\aliascntresetthe{claim} \newaliascnt{observation}{theorem}
\newtheorem{observation}[observation]{Observation}
\aliascntresetthe{observation} \newaliascnt{fact}{theorem}
\newtheorem{fact}[fact]{Fact} \aliascntresetthe{fact}
\newaliascnt{conjecture}{theorem}
\newtheorem{conjecture}[conjecture]{Conjecture} \aliascntresetthe{conjecture}
\theoremstyle{definition} \newaliascnt{definition}{theorem}
\newtheorem{definition}[definition]{Definition} \aliascntresetthe{definition}
\newaliascnt{assumption}{theorem}
\newtheorem{assumption}[assumption]{Assumption} \aliascntresetthe{assumption}
\newaliascnt{example}{theorem} \newtheorem{example}[example]{Example}
\aliascntresetthe{example} \theoremstyle{remark} \newaliascnt{remark}{theorem}
\newtheorem{remark}[remark]{Remark} \aliascntresetthe{remark}

%% Setup carried over from notes.tex
\usetikzlibrary{arrows.meta,positioning,calc,matrix,fit,decorations.pathreplacing}
\definecolor{aliceblue}{RGB}{32,96,151} \definecolor{boborange}{RGB}{183,96,27}
\definecolor{answergreen}{RGB}{34,111,81} \definecolor{linkblue}{RGB}{32,71,116}
\crefname{lemma}{Lemma}{Lemmas}
\crefname{proposition}{Proposition}{Propositions}
\crefname{corollary}{Corollary}{Corollaries}
\crefname{definition}{Definition}{Definitions}
\crefname{example}{Example}{Examples} \crefname{remark}{Remark}{Remarks}

%% Macros carried over from notes.tex
\newcommand{\D}{\operatorname{D}}
\newcommand{\Int}[2]{\mathsf{I}_{#1}(#2)}
\newcommand{\RP}[3]{\mathcal{P}(#1;#2,#3)}
\newcommand{\card}[1]{\lvert#1\rvert}
\newcommand{\ceil}[1]{\left\lceil#1\right\rceil}
\newcommand{\floor}[1]{\left\lfloor#1\right\rfloor}
\newcommand{\ind}{\mathbf{1}}
\newcommand{\NP}{\mathsf{NP}}
\newcommand{\BPP}{\mathsf{BPP}}
\newcommand{\SAT}{\mathsf{SAT}}
\DeclareMathOperator{\rank}{rank} \crefname{claim}{Claim}{Claims}
\crefname{observation}{Observation}{Observations} \crefname{fact}{Fact}{Facts}
\crefname{conjecture}{Conjecture}{Conjectures}
\crefname{assumption}{Assumption}{Assumptions}

\begin{document}
\title{NP-Hardness and Additive Inapproximability of Deterministic Communication Complexity}
\author{Serge Gaspers\thanks{University of New South Wales} \and
Tao Zixu He\thanks{University of Massachusetts Amherst} \and
Simon Mackenzie\footnotemark[1]}
\date{}
\hypersetup{pageanchor=false}
\begin{titlepage}
\maketitle
\thispagestyle{empty}

\begin{abstract}
Given the full truth table of a Boolean function whose input is split
between two players, how hard is it to determine the minimum number of
bits they must exchange to compute it? We prove that computing
deterministic communication complexity is NP-hard. This exact NP-hardness result, for the standard model was obtained independently and concurrently
with Hirahara, Ilango and Loff. This paper includes that result and
develops its construction to prove NP-hardness of approximation within
every fixed additive error. We also establish exact and fixed-error
NP-hardness in Yao's original model, where the players alternate after
every bit and only the receiver of the last bit needs to determine the
answer. Yao asked in 1979 how hard it is to compute communication complexity,
and our result answers his question under his original conventions.

In the standard model, we encode a graph-coloring instance so that a
non-colorable graph forces one extra bit of communication, and interlacing
makes this lower bound hold against every protocol. Each message of a protocol restricts the possible inputs to a subtable. We prove lower bounds that still hold on such subtables, so the output of one round can serve as the input of the next. To keep the table polynomial
for every fixed gap, we replace exponentially large sets of column tuples
by short explicit weighted lists. Under strict alternation, simulating
multi-bit messages would cost more than the gap, so we give a separate construction for this model.

For the standard model, we also show that accurate estimators would give
faster algorithms for sparse $3$-SAT. A
polynomial-time randomized estimator of communication complexity with additive error
$c\log T/\log\log T$, where $T$ is the number of table entries and
$0\le c<1/8$, yields a bounded-error $3$-SAT algorithm on $N$ variables
and $O(N)$ clauses in time $2^{O(N^\alpha)}$ for every fixed
$8c<\alpha<1$. A deterministic polynomial-time estimator with error
$c_*\log\log T$, for an effective absolute constant $c_*>0$, yields a
deterministic algorithm in time $2^{O((\log N)^2)}$ on such formulas.
The communication protocols being measured remain deterministic.
\end{abstract}

\end{titlepage}
\setcounter{page}{1}
\hypersetup{pageanchor=true}

\section{Introduction}
\label{sec:intro}

Alice and Bob know a Boolean matrix $M$. Alice receives a row index and
Bob receives a column index. They must determine the entry at their
intersection, communicating as few bits as possible. Their local
computation is free, and they must answer correctly for every pair of
indices. The deterministic communication complexity $\D(M)$ is the
smallest worst-case number of bits used by any such protocol.

We study how hard it is to estimate $\D(M)$ when the algorithm receives
the whole table as input. The table is total, meaning every entry is specified, and we write $T$ for its number
of entries. An estimator runs in time measured in $T$ and outputs a number,
not a protocol. A randomized estimator must meet its error guarantee with
probability at least $2/3$ on every table, while the protocols whose cost it
estimates remain deterministic. Yao posed the problem of computing $\D(M)$
exactly from the table when he introduced communication complexity in
1979~\cite{yao1979distributed}.

We prove that computing deterministic communication complexity is
NP-hard, and that this remains true when any fixed additive error is
allowed. The exact NP-hardness result for the standard model was first
presented in earlier versions of this
work~\cite{gaspers-he-mackenzie-npcc}, independently and concurrently
with Hirahara, Ilango and Loff~\cite{hirahara2025communication}.

We also prove exact and fixed-error NP-hardness in Yao's original
model: the players alternate after every bit, and only the receiver
of the last bit needs to determine the answer. The earlier
exact-hardness results did not cover these speaking and output
conventions. Our alternating construction therefore resolves Yao's
question in the model in which it was posed.

\paragraph{The challenge of approximation.}
In the standard model, the exact-hardness reduction maps a $3$-CNF formula
to a table whose optimal cost is one bit larger if the formula is
unsatisfiable than if it is satisfiable. An estimate within additive error
$a$ distinguishes the two cases only if their costs differ by more than
$2a$. A natural way to enlarge the gap is to combine several such tables
and hope that their one-bit gaps add up. However, a protocol for the combined problem may share
communication among them, so their costs need not add, and the lower bound
must hold against every such protocol.
Combining problems also enlarges the truth table. Since the reduction
writes out every entry, the table must have size polynomial in the size of
the formula for each fixed gap. Given $k$, our reductions compute a table
and a number $B$ such that the optimal cost is $\le B$ if the formula is
satisfiable and $\ge B+k$ otherwise. We call this an \emph{additive gap}
of $k$ bits. We do not establish hardness of approximation within any
fixed multiplicative factor above one.

\paragraph{Two rules for speaking.}
In an \emph{unrestricted} protocol, the bits sent so far determine who
speaks next, and the same player may send several consecutive bits. In a
\emph{strictly alternating} protocol, the players take turns after every
bit. Where an unrestricted protocol lets one player send two consecutive
bits, an alternating protocol must let the other player speak in between,
so the two optimum costs can differ.

\paragraph{Interlacing.}
Interlacing combines $q$ copies of a matrix $M$ through a pointer. In the $q$-fold interlace $\Int{q}{M}$, Alice
receives a row $x$ of $M$ together with an index $i\in\{1,\ldots,q\}$, and
Bob receives a list $(y_1,\ldots,y_q)$ of $q$ columns of $M$, possibly
with repetitions. The required output is $M(x,y_i)$.
Alice's index acts as a one-step pointer into Bob's input, as in pointer
chasing~\cite{nisan1993rounds}. \Cref{fig:lookup} shows the case $q=2$
for a $2\times2$ matrix.

There is a natural protocol: Alice sends $i$ using $\log q$ bits, and the
players then run an optimal protocol for $M$ on $x$ and $y_i$, for a total
of $\D(M)+\log q$ bits. But is communicating the index an unavoidable extra
cost? Perhaps Bob could speak first, or the players could combine
information about the index and the answer, saving bits.

The interlacing framework rules out such savings when $M$ is robustly
hard, that is, when every large subtable of $M$ is nearly as hard as $M$
itself. Then any protocol that does not start by revealing Alice's index
leaves a subproblem that is still too hard for the bits that remain, so
the natural protocol is optimal~\cite{mackenzie2025refuting,gaspers-he-mackenzie-npcc}.

The interlacing method for lower bounds was developed by Mackenzie and
Saffidine~\cite{mackenzie2025refuting} to refute the exact direct-sum
conjecture, and underlies our original NP-hardness
reduction~\cite{gaspers-he-mackenzie-npcc}.
If $M$ has $n$ columns, Bob has $n^q$ possible lists. To keep its table
polynomial, that reduction already replaced them by a polynomial-size
family of lists.
Repeating the construction needs more: the table produced by one round must
again have lower bounds on its large subtables, so that the next round can
use it as the base matrix $M$. An intermediate version obtained
additive inapproximability under ETH using subexponential-time
constructions~\cite{gaspers2026interlacingeth}.

\paragraph{Polynomial-size gap amplification.}

Our deterministic construction keeps explicit input sets of polynomial
size at every amplification step, so every fixed gap is constructed in
deterministic polynomial time. Our randomized construction keeps all inputs
during the amplification and represents them implicitly. Only at the end
does it sample rows and columns of the final table. A rank-based count shows
that the sample preserves the gap with high probability
(\cref{thm:compression}). In both constructions, the
lower bounds for unrestricted protocols come from interlacing. Because the
two constructions control the table size differently, they give different
tradeoffs between the approximation error and the running time of the
$3$-SAT algorithm that an estimator would imply.

Strict alternation needs a separate construction. The short protocols in our
unrestricted construction send multi-bit messages. Making them alternate
adds turns, and these extra bits can exceed the gap. We therefore
construct a different table, prove its lower bound using the fixed order of
speaking, and again keep only a small explicit set of Bob's inputs.

\paragraph{Conventions for the statements.}
By default, the conversation alone must determine the answer, which we call
\emph{public output}. We write $D_A^{\mathrm{alt}}(M)$ for the minimum
cost with public output and strict alternation, with Alice starting.
Yao's original convention instead allows the receiver of the last bit
to use their own input as well as the conversation to determine the
answer. We denote its optimum by $Y_{\mathrm{free}}(M)$, where the
protocol may choose its starting player independently of the inputs.
In both conventions, the bits exchanged must determine when the
conversation has ended.

\paragraph{From approximation to satisfiability algorithms.}
For errors growing with the table size, we state unconditional
algorithmic implications: an estimator for communication complexity
would yield a $3$-SAT algorithm with an explicitly bounded running time.
These tradeoffs do not assume ETH or any other running-time lower bound
for satisfiability. An inapproximability conclusion follows upon assuming
that the resulting $3$-SAT algorithm does not exist. The fixed-error
NP-hardness results need no such assumption.

\subsection{Results}
\label{sec:results}

The following table summarizes the fixed-error hardness results and two
examples of our general algorithmic implication. The full tradeoff is
stated in \cref{thm:time-tradeoff}.

\begin{center}
\small
\begin{tabular}{@{}p{.18\linewidth}p{.27\linewidth}p{.25\linewidth}p{.19\linewidth}@{}}
\toprule
Protocol model & Additive error & Consequence for $3$-SAT & Estimator type\\
\midrule
Unrestricted & Every fixed error & Polynomial time & Deterministic or randomized\\
\addlinespace
Strict alternation & Every fixed error & Polynomial time & Deterministic or randomized\\
\addlinespace
Unrestricted & $c_*\log\log T$ & $2^{O((\log N)^2)}$ & Deterministic\\
\addlinespace
Unrestricted & $c\log T/\log\log T$, $c<1/8$ & $2^{O(N^\alpha)}$, every $8c<\alpha<1$ & Bounded-error randomized\\
\bottomrule
\end{tabular}
\end{center}

The last two rows concern formulas with $N$ variables and at most $c_sN$
clauses, for any fixed density bound $c_s$. Each row describes the algorithm
that a polynomial-time estimator would imply; randomized estimators give
bounded-error algorithms. The first two rows are NP-hardness results under
deterministic reductions.

In each reduction below,
$\Phi$ is a $3$-CNF formula. The YES (satisfiable) case is $\Phi\in\SAT$,
and the NO case is $\Phi\notin\SAT$. The reduction computes the threshold
$B$ from the formula and the requested gap, without deciding
satisfiability.

\begin{theorem}[Fixed gaps for unrestricted protocols]
\label{intro:deterministic}
For every fixed integer $k\ge1$, a deterministic polynomial-time algorithm
maps a $3$-CNF formula $\Phi$ to an explicit total Boolean table $M$ and an
integer $B$ such that
\[
\Phi\in\SAT\ \Longrightarrow\ \D(M)\le B,\qquad
\Phi\notin\SAT\ \Longrightarrow\ \D(M)\ge B+k.
\]
The polynomial may depend on $k$. The output may be square with
power-of-two side length.
\end{theorem}

A different construction gives the same fixed-gap hardness under
strict alternation.

\begin{theorem}[Fixed gaps under strict alternation]
\label{yao:main}
For every fixed integer $k\ge1$, a deterministic polynomial-time algorithm
maps a $3$-CNF formula $\Phi$ to an explicit total Boolean table $M$ and an
integer $B$ such that
\[
\Phi\in\SAT\ \Longrightarrow\ D_A^{\mathrm{alt}}(M)\le B,\qquad
\Phi\notin\SAT\ \Longrightarrow\ D_A^{\mathrm{alt}}(M)\ge B+k.
\]
The polynomial bounds on output size and running time may depend on $k$.
\end{theorem}

For additive error $a$, choose $k>2a$ and compare the estimate with
$B+k/2$. Thus, for every fixed error, estimating
either cost in the preceding two theorems in deterministic polynomial time
implies $\mathsf P=\NP$, and bounded-error randomized estimation implies
$\NP\subseteq\BPP$. Changing the starting player costs at most one bit, and so does
announcing the answer known to the receiver (\cref{yao:comparison}). By increasing the fixed
gap, we therefore also obtain hardness under Yao's output convention.

\begin{corollary}[Yao's output convention]\label{yao:bpp}
For every fixed nonnegative additive error, a deterministic polynomial-time
estimator of $Y_{\mathrm{free}}(M)$ on all explicit total Boolean tables
implies $\mathsf P=\NP$, and a bounded-error randomized one implies
$\NP\subseteq\BPP$. The same conclusions hold for estimating the public-output
alternating cost with Alice first, Bob first, or a starter chosen
independently of the inputs.
\end{corollary}

Taking error zero in \cref{yao:bpp} answers Yao's original question by
proving NP-hardness under his speaking and output rules. We do not establish
a growing-error bound for the alternating model.

For unrestricted protocols, choosing the gap as a function of the formula
size gives an explicit tradeoff between approximation error and the
running time of the resulting $3$-SAT algorithm
(\cref{thm:time-tradeoff}). Here $N$ counts variables in the formula and
$T$ counts entries in the communication table. The following two
consequences illustrate the randomized and deterministic regimes; the
constant $1/8$ uses the sharper size bound for the randomized
construction.

\begin{corollary}[Randomized estimation and $3$-SAT running time]
\label{thm:quantitative}
Fix $0\le c<1/8$. Suppose a bounded-error randomized polynomial-time
estimator approximates $\D(M)$ within additive error
$c\log T/\log\log T$ on every sufficiently large explicit total Boolean
table with $T$ entries. Then, for every fixed $8c<\alpha<1$ and every
fixed $c_s$, sparse $3$-SAT on $N$ variables and at most $c_sN$ clauses has
a bounded-error randomized algorithm with worst-case running time
$2^{O(N^\alpha)}$. It suffices for the estimator to work on square tables
whose side length is a power of two.
\end{corollary}

For the smaller error scale below, the deterministic construction gives
a quasipolynomial-time algorithm for sparse $3$-SAT.

\begin{corollary}[Deterministic estimation and $3$-SAT running time]
\label{det:time}
There is an effective absolute constant $c_*>0$ with the following property.
If a deterministic polynomial-time estimator approximates $\D(M)$ within
additive error $c_*\log\log T$ on every sufficiently large explicit total
Boolean table, then, for every fixed $c_s$, $3$-SAT on $N$ variables and
at most $c_sN$ clauses has a deterministic algorithm running in time
$2^{O((\log N)^2)}$. It suffices for the estimator to work on square tables
whose side length is a power of two.
\end{corollary}

These consequences quantify the improvement in satisfiability algorithms
that the corresponding estimators would yield.
\Cref{sec:time-tradeoff} gives the general tradeoff and explains how to
choose the gap for other running-time bounds.

The proofs of \cref{intro:deterministic,det:time} appear in
\cref{sec:deterministic-endpoints}. We prove
\cref{thm:time-tradeoff,thm:quantitative} in \cref{sec:time-tradeoff}, and
\cref{yao:main,yao:bpp} together at the end of
\cref{sec:alternating}. \Cref{sec:overview} explains the ideas behind these proofs.

\subsection{Formal verification and proof dependencies}
\label{sec:formal-verification}

A companion Lean development provides machine-checked proofs for the
constructions and hardness statements covered by its statement
correspondence~\cite{interlacing2026lean}. Its formal statements use
explicit interfaces for the cited source reductions and record their
hypotheses.
The public proof inspector~\cite{interlacing2026inspector} links manuscript
statements to the corresponding Lean declarations and exposes their
definitions and hypotheses. The repository provides the source, build
instructions, and verification records.

Communication lower bounds, explicit constructions, and bounds on output
size and running time are distinct proof obligations. The scope of the
formalization is determined by the statements proved in the companion
development and the external interfaces they assume. The fixed-error
NP-hardness results and the implications from estimation to $3$-SAT
algorithms stated above do not assume ETH. A running-time hypothesis is
needed only to turn a growing-error implication into an
inapproximability conclusion.

The earlier direct-sum construction of Mackenzie and
Saffidine~\cite{mackenzie2025refuting} has separately been formalized in
Lean by Moakhar et al~\cite{moakhar2026beyond}; its statements and proofs are available in
a public inspector~\cite{beyondlibrary2026inspector}.

ChatGPT and Codex were used in proof development and checking;
the \hyperref[sec:ai-methodology]{AI methodology statement}
describes their roles.

\clearpage
\section{Technical overview}
\label{sec:overview}

Our reduction turns a coloring obstruction into one extra bit of
communication. A proper coloring gives a protocol of a known cost $C-1$.
Conversely, interlacing forces every protocol of that cost to commit to
a partition of Alice's inputs before Bob reveals which constraint is
being tested. In the NO case, the extracted assignment violates a
constraint, forcing an additional bit. The difficulty is to prove this
against arbitrary interaction, and then to repeat the construction
without losing either the lower bound or polynomial output size.

We explain three steps: the one-bit gap, the lower bounds under input
deletion that make the gap repeatable, and the explicit short input lists
that make every fixed number of rounds a polynomial-time reduction.
Strict alternation requires a different encoding, described last.

\subsection{Constructing the source matrix}

Our source reduction maps a formula to a graph of maximum degree
$\Delta=O(1)$ that is three-colorable if the formula is satisfiable (the
YES case), and otherwise leaves at least a constant fraction $\gamma$ of
its edges monochromatic under every three-coloring (the NO case,
\cref{src:gap-colouring}). This NO condition survives small deletions.
Since $|E|\ge|V|/2$, deleting fewer than $\gamma|V|/(4\Delta)$
vertices destroys fewer than $\gamma|E|/2$ edges, so at least
$\gamma|E|/2$ monochromatic edges remain.

The construction takes a base communication matrix $R$ of cost $d=\D(R)$
and adds rows encoding the graph to interlaced copies of $R$. The first
round starts from a simple matrix of known cost (\cref{cl:seed}). We
allow a general base matrix so that the construction can be repeated, with
the output of one round as the base matrix of the next. We also need lower bounds
that survive deletion of some inputs. We measure the retained rows and columns
using fixed positive rational weights, summing to one on each side. These act
as fractional counts: they let us control how much each input family contributes,
independently of its number of labels. Later, adjusting weights also lets us
shorten tuple lists while preserving the counts needed by the proof.
The weights do not change communication cost: every protocol must still compute
every entry correctly. The deterministic construction never samples from them,
and its final output is an ordinary Boolean table with the weights discarded.

A proper three-coloring places the vertices in
three groups so that the endpoints of every edge lie in different groups. We call
these groups \emph{bins}. The interlacing construction uses a power-of-two
number $s\ge4$ of bins, so the extra bins must be filled without giving
the vertices extra colors. We treat each graph vertex as an \emph{item} to be
placed in a bin, and add $s-3$ auxiliary items. At an edge, its two
endpoints and the auxiliary items are \emph{active}, and all other items
are \emph{inactive}. We require that each bin contain at most one active
item at every edge. The auxiliary items are active at every edge, so no two
of them share a bin. Since every vertex of our source graph
lies on an edge, no vertex can share a bin with an auxiliary item, and the
vertices share the remaining three bins, exactly as in three-coloring.

Each item $i$ gives a family of Alice's possible inputs: the pairs
$(i,u)$ with $u$ a row of $R$. A coloring partitions these inputs by the
bin of their item. Naming the bin need not identify Alice's item and
does not reveal her row of $R$.

The full construction uses $4s$ labeled groups of items, each containing a
copy of the graph's items. For each edge in each group, a binary vector over all items marks the
active items of that group with 1 and every other item with 0. We index these
vectors by $\alpha$. In the NO case, every assignment of items to bins has a
constant fraction of indices at which one bin contains two active items and
a third item that is inactive there, even after deleting a small fraction of
the items (\cref{src:item-witness}). We say there is a \emph{collision}
at such an index. The inactive third item comes from another group. There
are $4s$ groups but only $s$ bins, and at most $s$ groups can have a bin
containing items of no other group. Hence for most groups the bin holding an
active pair also holds an item of another group, which is inactive at that index.
The reduction never computes an
assignment. Instead, the lower-bound proof extracts one from any proposed
short protocol.

The matrix has two kinds of Alice inputs. An \emph{item row} consists
of an item $i$ and a row $u$ of the base matrix $R$. The other inputs,
called \emph{template rows}, are grouped into $s$ bins and are obtained
from $R$ by nested interlaces. Their purpose is to force every protocol
of cost at most $C-1$ to begin with $\log s$ consecutive bits from Alice.
These bits determine a bin before Bob reveals any index.

Bob's columns have two families. The first supplies a separate index
and list of base columns for each bin, together with shared base-column
data; it enforces the initial $\log s$ bits. The second uses the same
index $\alpha$ and the same list of base columns in every bin. At this
index, the template rows use $s-2$ positions of the list, active items
use distinct additional positions, and inactive items use one further
common position. Thus two active items and one inactive item in the
same bin expose $s+1$ positions.

\paragraph{Defining the table without a coloring.}
A template row carries its bin label $p$ and a tuple $z$. For each
index $\alpha$, the tuple specifies a pair $z_\alpha=(r,u)$: a position
$r\in\{1,\ldots,s-2\}$ and a base row $u$. Bob's first-family column
supplies $(\alpha_p,\gamma_p)$ for each bin, where $\gamma_p$ is a list
of $2s-1$ base columns. The entry is $R(u,\gamma_{p,r})$, with
$(r,u)=z_{\alpha_p}$. Thus each template row uses only its own bin's
data. Letting these data vary independently gives the outer interlace
that forces Alice's initial messages.

The tuple $z$ is important. If a template row carried just one pair
$(r,u)$, its entry after the bin was chosen would not depend on Bob's
index $\alpha_p$. The template subproblem would no longer require that
index to be communicated. Instead, $z$ supplies a separate pair for
each index: Bob selects which part of Alice's input matters, and then
Alice's position $r$ selects which part of Bob's list matters. The
players are left with one instance of $R$. These successive selections
are the nested interlaces, with transposition exchanging the players'
roles between levels.

The outer bin choice adds another interlace, since Bob has independent
data for every bin. We must prove that any sufficiently short protocol
starts by resolving that choice; specifying a short protocol that does
so would not suffice. Once the prefix is forced, the item rows fall
into the same parts as the template rows, and a bin assignment can be
extracted from them.

An item row $(i,u)$ has no bin label: the reduction cannot use a
coloring to assign it one. Bob therefore also holds a shared list
$\xi=(\xi_1,\ldots,\xi_{s+1})$. For each index $\alpha$, give the active
items distinct positions $\pi_\alpha(i)$ in a fixed public order.
If all Bob's indices agree on $\alpha$ and item $i$ is active there,
the entry is $R(u,\xi_{\pi_\alpha(i)})$; otherwise it is
$R(u,\xi_{s+1})$. These rules use only the graph and the two inputs.
In the coloring protocol, the bin and index identify an active item,
and Bob can check agreement of his indices privately to select the
appropriate base column.

The second family sets every bin's index to the same $\alpha$ and every
list to the same $\gamma$, with
$\xi=(\gamma_{s-1},\ldots,\gamma_{2s-1})$.
The template and item rows now read positions of one common list.
A single protocol must handle both families, and before Bob speaks,
Alice's messages depend only on her row. Consequently, the partition
forced by the first family is also the partition tested by the second.

The second family has separate labels and a fixed positive share of
the column weight. In the full-product construction, its columns
already occur in the first family but may have very small total
weight there. The separate weight prevents a small deletion from
removing all collision tests. \Cref{def:source-matrix} gives the full
entry rules and weights.

\paragraph{A coloring gives a protocol.}
Fix the bin assignment supplied by a proper three-coloring. Alice names
her bin $p$, using $\log s$ bits, and Bob sends its index $\alpha_p$,
using $\log q_2$ bits, where the number $q_2$ of indices is a power of two.
Alice then sends one of
\[
\underbrace{s-2}_{\text{template positions}}
\quad+\quad
\underbrace{1}_{\text{active item}}
\quad+\quad
\underbrace{1}_{\text{inactive item}}
=s
\]
possible messages. An active item is uniquely determined by its bin and
the index, since a proper coloring puts at most one active item in each
bin. Bob can therefore identify the selected base column from his input,
and the players finish by running a protocol for $R$.

The two Alice messages each cost $\log s$ bits, giving total cost
$C-1:=d+2\log s+\log q_2$. For an arbitrary graph, Alice can instead name
her active item's position explicitly. This uses at most $2s-1$ message
values and costs one additional bit. Thus the table always has a
cost-$C$ protocol, and a coloring saves one bit
(\cref{src:upper}).

\subsection{Why the reduction rules out every short protocol}

We prove that the NO-case lower bound survives deletion of a small fraction
of the row and column weight, so that it can be reused in later rounds.
Suppose a protocol of cost $C-1$, with any pattern of interaction, computes
the remaining matrix. From such a protocol we
extract a single bin assignment by grouping Alice's rows according to the
first $\log s$ bits, which are sent before Bob can reveal any index. For a NO graph, this fixed assignment must contain
collisions.

\paragraph{A short protocol fixes a bin assignment.}
Group the template rows by their bin $p$. In the first column family,
Bob supplies a separate index and list of base columns for every bin,
and every combination occurs. A template row in bin $p$ depends only
on the data for that bin. On template rows the matrix is therefore an
$s$-fold interlace whose blocks are the bins.

Within one bin, the template rows already require
$d+\log q_2+\log s$ bits, by two applications of the interlacing lower
bound (\cref{src:profiles}). Of the $C-1$ available bits, only $\log s$
remain to tell the bins apart. Suppose Bob sends a bit while $u$ bins are
still possible. Every Alice row possible before his bit remains possible.
The interlacing lower bound for $u$ bins still applies to the
restriction keeping the larger half of his columns, and one bit fewer now
remains, so the remaining bits do not suffice. Similarly, if one of
Alice's bits sent more than half of the $u$ bins to one side, that side
would need more bits than remain. The
separation theorem (\cref{cl:separation}) therefore forces the first
$\log s$ bits to be Alice's and to split the bins evenly. Grouping her rows
by these $\log s$ bits gives $s$ parts, one for each bin. Each part contains almost all the template-row weight of one bin
and very little of the weight of the others.

We now discard the few items that lost more than half of the weight of
their rows to the initial deletion, and
assign each remaining item to a part containing at least a $1/(2s)$
fraction of the weight of its rows $(i,u)$. Such a part exists by
averaging. The rows of one item need not stay in one part,
since different base rows $u$ may follow different paths. However, the
assignment is fixed before Bob has spoken, so it cannot depend on
his edge. For a NO graph, \cref{src:item-witness} guarantees many indices
at which one assigned bin contains two active items and an inactive item.
Thus any proposed short protocol determines an assignment with such triples,
which a proper coloring would avoid.

\paragraph{A collision gives an extra position.}
Restrict Bob to the second column family at a collision index $\alpha$.
All bins now use the same index and the same list of base columns.
Within the collision's bin, the template rows and the three items read
\[
\underbrace{s-2}_{\text{template positions}}
\quad+\quad
\underbrace{2}_{\text{distinct active positions}}
\quad+\quad
\underbrace{1}_{\text{inactive position}}
=s+1
\]
distinct positions. The inactive item is essential: two active items
alone would give only $(s-2)+2=s$ positions. Since $s$ is a power of two,
naming one of $s+1$ positions needs one more bit than naming one of $s$.
To obtain the communication lower bound, we must show that interaction
cannot bypass this obstruction.

\paragraph{Why collisions force more communication.}
The extra position explains the gap, but counting positions alone does not
prove a lower bound against an interactive protocol. Bob's replies restrict which lists of base columns remain possible, and a lower
bound for the original column set need not survive this restriction.
We restrict to the second family, in which every bin has the same
index $\alpha$. After Alice's $\log s$ bits, Bob knows $\alpha$ and may send bits
that narrow down the possible values of $\alpha$ before Alice speaks again,
so that each branch of the protocol covers only some indices. If
collisions occurred at only a few indices, Bob could isolate them on a
branch with a larger budget. We therefore need collisions at a constant
fraction of all indices.

We track the retained list inputs through Bob's messages and stop when
Alice next speaks. Each usable index initially contributes one tracked
set; a collision contributes an additional mark, recording its extra
bit. An unmarked set follows a child retaining at least half its weight.
A marked set either retains its mark on one child or splits into two
unmarked sets. In the latter case its contribution changes from two to
one plus one. Thus the total count is preserved.

A stopping node reached after $r$ Bob bits can accommodate a count of at
most $q_2 2^{-r}$ within its remaining communication budget. The stopping
nodes are prefix-free, so Kraft's inequality bounds their total capacity
by $q_2$. Initially, however, the count exceeds $q_2$: the collision
indices outnumber the indices lost to the initial input deletions.
This contradiction gives the NO lower bound
(\cref{lem:collision-tracking,src:near-full}). Together with the coloring
protocol, it yields the one-bit gap against every protocol, whatever
its pattern of interaction.

\subsection{Making the one-bit gap repeatable}

A one-bit gap does not automatically yield a larger gap. Using the output
as the next base matrix requires lower bounds on the restrictions that
remain after messages, and a hard matrix can have easy restrictions.
For example, restricting Bob's tuples in $R(u,y_i)$ to those with all
coordinates equal removes the need to learn $i$.

We therefore maintain a stronger invariant in the NO case. Using the
fixed positive rational input weights, for a threshold $\rho$ fixed
across amplification rounds, require
\[
\begin{array}{c|c|c}
\text{retained row weight}&\text{retained column weight}
   &\text{communication lower bound}\\ \hline
\ge\rho&\ge3/5&d\\
\ge\rho&\ge3/10&d-1 .
\end{array}
\]
The threshold $\rho$ may depend on the source size. We call these the
\emph{deletion bounds} (\cref{def:profile}). The second column threshold
is half the first: after Bob sends one bit, the heavier branch retains
at least half the column weight, and the allowed lower bound drops by
one. The YES upper bounds do not require this invariant.

The graph construction supplies its extra bit only on restrictions
retaining almost all the weight on both sides. The deletion bounds
needed for another round allow most of the rows and a constant
fraction of the columns to disappear. Two further interlaces close
this gap.

\paragraph{Preserving the extra bit after input deletion.}
Transpose the source matrix to $N$, so its two column families become
row families, and form $\Int{w}{N}$. Alice holds a copy index and a row
of that copy; Bob holds one column for each copy. Consider a restriction
retaining almost all row weight but only a small specified fraction of
the tuple weight. Suppose it has a protocol of cost $C+\log w-1$, one
bit less than naming a copy and running the cost-$C$ protocol for $N$.

Pad shorter paths with fixed Alice bits, and stop each path just before
Bob's first bit or just after Alice has sent $\log w$ bits, whichever
comes first.
The first-family rows force almost all their weight in most copies to
reach distinct stopping nodes, one assigned to each such copy. These
nodes are reached after $\log w$ Alice bits, leaving $C-1$ bits.
Bob has not spoken, so his tuple set is still the initial restriction.
Projection estimates show that almost every copy still has almost
all its columns available. If its assigned node also kept almost all
its rows, the source lower bound would require $C$ further bits.
The node must therefore miss a substantial amount of its copy's rows.

Almost all first-family weight is already at the assigned nodes, and
little weight was deleted initially. Summing over the assigned copies
therefore forces substantial surviving second-family weight to lie
elsewhere. Since only Alice has spoken, every surviving row reaches
exactly one stopping node. It must reach either another copy's
assigned node or a node with no assigned copy.

Neither group can receive enough weight. An assigned node already
contains the hard first-family rows of its own copy; receiving too
much second-family weight from other copies would force its cost
above the remaining budget. The unassigned nodes have only a small
total capacity. The combined receiving bound is smaller than the
weight that must be displaced, giving a contradiction
(\cref{res:alignment}). Transposing the resulting matrix exchanges
the two retained-weight conditions, yielding $W$ with its full lower
bound when a small fraction of rows and almost all columns remain.

\paragraph{Recovering the bit in the deletion bounds.}
The matrix $W$ also satisfies the desired deletion bounds for constant
retained column weight, but one bit below its full cost. A final
interlace combines this guarantee with the full-cost bound just proved.
For full products, a tuple set of retained weight $v$ cannot have small
projections at too many coordinates. If its projection onto coordinate
$i$ has original weight $p_i$, containment in the product of its
projections gives
\[
v\le\prod_i p_i,
\qquad\text{hence}\qquad
\sum_i\log(1/p_i)\le\log(1/v).
\]
Among sufficiently many usable copies, at least one therefore retains
almost all its columns, allowing the full-cost bound for $W$ to apply.

Finding one such copy alone would recover only the cost of $W$; we
must also account for the bits needed to identify a copy. An induction
over messages does this. At an Alice bit, sufficiently many row blocks
remain on one branch. At a Bob bit, keeping the heavier tuple set
increases its logarithmic loss by at most one. Choosing enough copies
at the outset accommodates these losses. The deletion bounds control
the intermediate restrictions, and the full-cost bound closes the
remaining one-bit deficit (\cref{cl:bridge}). For the explicit lists,
the conditioning argument described next supplies the required
projection control instead of the full-product inequality
(\cref{det:restore}).

Consequently, one complete round has a known overhead $h$ and satisfies
the following guarantee. In the YES case,
\[
\D(R^+)\le\D(R)+h.
\]
In the NO case, if $\D(R)=d\ge3$ and $R$ satisfies the deletion bounds
at value $d$, then
\[
\D(R^+)=d+h+1,
\]
and $R^+$ satisfies the same deletion bounds at value $d+h+1$
(\cref{thm:amplifier}). This restoration is what makes iteration
possible. Starting from a matrix of known cost $d_0$ satisfying the
deletion bounds gives
\[
\text{YES: }\D(R_k)\le d_0+kh,
\qquad
\text{NO: }\D(R_k)=d_0+kh+k.
\]
The threshold $d_0+kh$ is computable without deciding whether the graph
is colorable.

\subsection{Keeping the deterministic construction polynomial}

Allowing every tuple of $q$ columns from a set $Y$ creates $|Y|^q$ inputs.
We replace each product by an explicit list of polynomially many tuples
with positive rational weights. For an absolute constant $D$, the list
has exactly the same joint weights as the full product on every set of
at most $D$ coordinates. This is \emph{limited independence}. We build
and enumerate the entire list deterministically; the weights do not
represent random choices made by the reduction.

The construction extends tuples one coordinate at a time and removes
tuples by adjusting their weights. Associate with each tuple a zero--one
vector recording the assignments it realizes on every set of at most
$D$ coordinates, together with a coordinate recording total weight.
The required joint weights specify the weighted sum of these vectors.
If the list is longer than the vector dimension, a linear dependence
lets us change the weights without changing that sum. Since total
weight is fixed, the change has both positive and negative components;
we can continue until one tuple has weight zero and remove it, keeping
all other weights nonnegative. Compressing after each new coordinate
keeps the intermediate lists short, avoiding enumeration of an
exponentially large product. \Cref{det:code} proves the support and rational
bit-length bounds. The original entry rules still apply to every listed
tuple, so the upper-bound protocols are preserved.

The lower bound is subtler: conditioning on earlier messages can destroy
independence. For example, retaining only tuples whose first two
coordinates agree makes those coordinates dependent. Our conditioning
lemma shows that, after retaining any event of weight at least $v$,
there are only $\mathrm{poly}(1/v)$ exceptional coordinates. On every
set of at most eight other coordinates, the conditional distribution
is close in total variation to the original product distribution
(\cref{det:conditioning}). A moment bound proves this: too many
disjoint sets of coordinates with substantially changed distributions
would contradict the weight of the conditioning event.

We choose the interlacing widths and deletion thresholds so that enough
coordinates remain outside the exceptional set. Projecting onto these
coordinates gives actual inputs of the retained matrix, with enough
weight to apply the deletion bounds. Thus the proof accounts for the
dependencies created by messages.

If the base input set has $m$ labels, each list has at most
$(1+qm)^D$ tuples, and its rational weights have polynomial bit length
(\cref{det:code}). Each round therefore takes time polynomial in the
source graph size $n$ and the encoding length of the current matrix,
with an exponent independent of the round. After $k$ rounds, size and
running time are at most $n^{O(A^k)}$ for an absolute constant $A>1$.
Every fixed gap therefore gives a deterministic polynomial-time
reduction (\cref{det:reduction,intro:deterministic}).

A suitable choice of $k=\Theta(\log\log N)$ for sparse formulas makes
the construction quasipolynomial in $N$, while $\log\log T=O(k)$.
An estimator with a sufficiently small constant times $\log\log T$
error would therefore decide sparse $3$-SAT in quasipolynomial time
(\cref{det:time}).

\subsection{Two further consequences of keeping fewer inputs}

\paragraph{Sampling at the end gives a stronger randomized bound.}
For the randomized result, we keep full products of tuples throughout the
amplification. The table is then too large to print, but we can compute
single entries from short descriptions of the inputs. After all $k$
rounds, we use the input weights as sampling probabilities: we draw $m$ rows
and $m$ columns independently and print the $m^2$ entries at their intersections.
Here the reduction uses randomness; the communication protocols being measured
remain deterministic.

In the NO case, a protocol may be chosen after the sample is drawn, so the
sample must defeat every protocol of depth one less than the NO cost
simultaneously. We call a set of columns \emph{compatible} with a protocol
and a set of rows if the protocol is correct on them. Over all protocols of
depth $t$ and all row sets, a matrix of real rank $r$ has at most
$2^{O(r\log(r+2)\,2^t)}$ inclusion-maximal compatible column sets
(\cref{thm:rank-count}). The count depends on rank and protocol depth,
not on the original number of input labels. For row sets of weight $\ge3/5$, the
deletion bounds give each compatible column set weight $<3/5$. A union
bound over these sets then shows that, with high probability,
$m=\Theta(r\log r\cdot2^t)$ random columns are not all contained in any
one of them. Once such a column sample is fixed, every row set compatible
with a short protocol on the sampled columns has weight $<3/5$, so the same
union bound applies to the sampled rows (\cref{thm:compression}). The
sample need only preserve the gap, since no further amplification follows. Each interlace multiplies the rank by at
most its width, so the rank after $k$ rounds is at most $2^{d_0+kh}$
(\cref{prop:iteration}). For a sparse formula on $N$ variables the graph has
$N^{1+o(1)}$ vertices and $h=(1+o(1))\log N$,
so taking $k\approx N^\alpha/\log N$ gives
$\log T=(4+o(1))k\log N$ and
$k=(\alpha/4+o(1))\log T/\log\log T$.
The construction then takes time $2^{O(N^\alpha)}$.
An estimate within additive error $c\log T/\log\log T$ distinguishes the
two cases whenever $\alpha>8c$, giving \cref{thm:quantitative}.

\paragraph{Strict alternation needs a different encoding.}
The coloring protocol above uses consecutive bits from the same player.
Inserting dummy turns can cost more than the gap, so it does not give
hardness under strict alternation.

We instead start from the hypergraph-coloring reduction of Guruswami and
Lee~\cite{guruswami2018balanced}. For a power of two $s$, each edge has
$s^2$ vertices. In the YES case there is an $s$-coloring in which every
edge contains at least $s-1$ vertices of every color. In the NO case,
every coloring with at most $2^s$ colors has a monochromatic edge
(\cref{yao:source}).

Alice holds a vertex and a tuple of edge addresses. Bob holds an edge,
a coordinate selecting one address, and a Boolean array. The selected
address and the vertex's position in Bob's edge identify the array bit
to output; the output is zero if the vertex is outside the edge
(\cref{def:alternating-table}).
Write $s=2^t$, and let $b$ bits name an edge. In the YES case,
the players exchange the color, address and edge information in
alternating turns. The coloring leaves at most $2s-1$ candidate
vertices of Alice's color in Bob's edge, giving total cost
$2b+4t+2$ (\cref{prop:alternating-upper}).

For the NO lower bound, record Alice's responses to \emph{every}
possible sequence of Bob's earlier bits during her first $t$ turns.
There are at most $2^{2^t-1}\le2^s$ such response rules. Assign each
vertex a rule used by many of its address tuples. The NO promise gives
a monochromatic edge, so many rows sharing one early response rule
must still identify many different array positions. Counting only the
transcript of one interaction would not suffice, because Bob's
messages may depend on his entire array.

We construct a polynomial-size explicit list of arrays in which no
prescribed pattern on many positions occurs too often. Following a
hypothetical short protocol while retaining large sets of rows and
arrays leaves more possible arrays than its remaining response
patterns can cover
(\cref{yao:array-list,yao:small-bias-escape,yao:lift}).
This gives a NO lower bound of $2b+6t-6$, and hence a gap of $2t-8$.
Allowing either starter and receiver-only output loses at most two
bits (\cref{yao:comparison}). Choosing $t$ as a sufficiently large
constant therefore gives a deterministic polynomial-time reduction
for every fixed gap in Yao's convention.

\section{Related work}

\paragraph{Which optimization problem is hard?}
For explicitly given total tables, Kushilevitz and Weinreb proved that
approximating communication complexity is hard under cryptographic
assumptions~\cite{kushilevitz2009complexity}.
Ilango, Loff and Oliveira proved NP-hardness and strong
inapproximability for \emph{partial} matrices, where the protocol is
required to be correct only on specified input
pairs~\cite{ilango2020multioutput}. Our tables are \emph{total}: every
pair of a row and a column has a specified bit, and the protocol must
be correct on all of them.

Restrictions on the conversation also matter. For total matrices,
Hirahara, Ilango and Loff proved additive hardness for the number of
bits in a three-message protocol, and hardness of approximating the
number of leaves of a protocol tree for every fixed number $k\ge3$ of
messages~\cite{hirahara2021constant}. A message is a block of
consecutive bits from one player, and a leaf is a possible end of the
conversation. Neither restricting the number of messages nor counting
leaves gives the same optimization problem as minimizing bits under
strict alternation or unrestricted speaking order. The exact-hardness
results for the unrestricted
model~\cite{gaspers-he-mackenzie-npcc,hirahara2025communication} left open
both hardness with one bit of additive error and Yao's original
alternating model, as Hirahara, Ilango and Loff
noted~\cite{hirahara2025communication}. Our fixed-gap theorems resolve both questions.

\paragraph{Ingredients of the constructions.}
For the unrestricted reductions, we obtain hard graph-coloring instances
by combining reductions of Dinur and of Bogdanov, Obata and
Trevisan~\cite{dinur2007pcp,bogdanov2002lower}. In these graphs, each
vertex meets only a constant number of edges. The alternating reduction
uses the hypergraph-coloring promise of Guruswami and
Lee~\cite{guruswami2018balanced} described in \cref{sec:overview}.

Both reductions also replace large input sets by small explicit
ones that keep some of their properties. For the deterministic
interlacing construction, we follow Koller and
Megiddo~\cite[Section~3]{koller1994sample}: we extend a list of tuples one
coordinate at a time and delete tuples without changing the joint
weights of the values on every set of a bounded number of coordinates. For strict alternation, the
array list is based on the small-bias construction of
\citet{alon1992smallbias}.

Bounds on the number of column sets on which short protocols can be
correct, expressed in terms of matrix rank, appear in graph algorithms~\cite{bui2010rankwidth,oum2014faster}
and in work on hard submatrices~\cite{hrubes2025hard}. To count such sets, we use the
row-sum method of Oum, S{\ae}ther and
Vatshelle~\cite{oum2014faster}: summing Boolean rows identifies the
columns in which all those rows are zero. Mackenzie's hardness
condensation~\cite{mackenzie2026condensation} also counts the input sets
on which short protocols succeed, to shrink a communication problem while
keeping its cost high relative to the new input length. In contrast, our
sampling theorem (\cref{thm:compression}) keeps the number of bits above a
prescribed bound, and we prove it directly.



\paragraph{Reading guide.}
\Cref{sec:models} fixes the communication conventions and the use of input
weights. We then state the amplification theorem and its
iteration, prove the graph encoding and the two further interlaces, and
construct the short input lists used by the deterministic reduction.
We then give the final-sampling argument and, separately, the strictly
alternating construction. The appendices contain the remaining interlacing
estimates, list algorithms, parameter checks and coloring reductions.


\clearpage
\section{Preliminaries}
\label{sec:models}

\subsection{Communication and estimation}

An estimator receives the whole table, whereas each player knows only its own input.

\begin{definition}\label{def:table-estimation}
A table $M:X\times Y\to\{0,1\}$ is total, and its input sets $X$ and $Y$ are
finite and nonempty sets of labels. Equal rows or columns may have different indices.
Its size is $T=|X||Y|$.
A deterministic public-output protocol is a binary tree whose internal
nodes specify a speaker and a bit-valued function of that player's input
and the preceding transcript. Each leaf has a fixed output, and the
protocol must compute every entry correctly. Its cost is its maximum
depth, charging every bit, including fixed bits. Local computation is free.
Write $\D(M)$ for the least cost with unrestricted speaking order.

An additive-$a(T)$ estimator returns $\widehat D(M)$ satisfying
$|\widehat D(M)-\D(M)|\le a(T)$. Its output must admit polynomial-time
comparison with polynomial-length binary rationals. A bounded-error
randomized estimator satisfies this guarantee with probability at least
$2/3$ on every input. Polynomial time means worst-case time polynomial
in $T$. Estimators for the alternating costs below follow the same conventions.
\end{definition}

In the alternating model, the starting player and the output convention
can each affect the cost.

\begin{definition}[Alternating costs]\label{yao:model}
In strict alternation the players speak one bit at a time, in turn.
Let $D_A^{\mathrm{alt}}$ and $D_B^{\mathrm{alt}}$ be the minimum
public-output costs when Alice or Bob starts, respectively, and let
$D_{\mathrm{free}}^{\mathrm{alt}}=\min\{D_A^{\mathrm{alt}},D_B^{\mathrm{alt}}\}$.
The starter is chosen independently of the inputs. Every bit is charged,
and stopping is public: the transcript determines whether communication
has ended.

Yao's receiver-output cost $Y_{\mathrm{free}}$ permits the same choice of
starter, but the receiver of the last bit determines the answer from
that player's input and the transcript. At depth zero, the protocol
designates one player whose input determines the answer. Stopping is
still public and correctness is required on every input.
\end{definition}

Our quantitative statements use sparse $3$-SAT: the input has $N$ variables
and at most $c_sN$ clauses, where $c_s$ is fixed in advance. Running times
are worst-case bounds, including all random tapes for randomized algorithms;
bounded error means success probability at least $2/3$ on every formula.
Constants may depend on $c_s$ and on the estimator. We do not apply a
sparsification step with an uncharged exponential cost: all stated running
times refer directly to these sparse inputs.

All communication protocols in this paper are deterministic and exact.
Only estimators and randomized reductions use random bits. All logarithms
are binary unless written $\ln$. Write $[q]=\{1,\ldots,q\}$ and
$\operatorname{pow}_2(a)=2^{\ceil{\log a}}$ for the least power of two
at least $a\ge1$.

Every $n\times n$ table has $\D(M)\le\ceil{\log n}+1$: Alice sends her
row index and Bob sends the entry.

\subsection{Relabeling and input weights}

We keep the matrices small in two ways, by explicit short lists and by
sampling. In both, each new input decodes to an original input, as follows.

If $f:X'\to X$, $g:Y'\to Y$ and $F'(x',y')=F(f(x'),g(y'))$, each player can compute $f(x')$ or $g(y')$
from its own input and run a protocol for $F$. Hence $F'$ has at most the
communication cost of $F$ in every model above, and at most its rank. Both
are equal when $f$ and $g$ are onto. Transposition preserves rank and unrestricted
cost, and exchanges $D_A^{\mathrm{alt}}$ and $D_B^{\mathrm{alt}}$.

In particular, duplicating rows and columns makes a table square with
power-of-two side length without changing its cost. We may also
require the side length to be at least a prescribed bound $b$ polynomial in
the input size. The side length is then less than $2\max\{b,|X|,|Y|\}$, so
the enlarged table can be printed in polynomial time.

For the unrestricted construction, we assign the labels in each input set
positive rational weights summing to one. For row weights $\mu_X$ and column weights $\mu_Y$, a set's
weight is the sum of its labels' weights, for example
$\mu_X(A)=\sum_{x\in A}\mu_X(x)$. Uniform weights give ordinary fractions
of labels. Unequal weights fix the shares assigned to different input families
and preserve the required counts when we shorten tuple lists. The weights
are auxiliary data for the construction and its lower-bound proof;
communication cost depends only on the table, and every input must be computed
correctly. The final deterministic reduction discards the weights.
Distinct labels keep separate weights, even when their rows or columns agree.

If $X'$ has weights $\mu'$, a decoder $f:X'\to X$ induces weights on $X$ by adding the weights of
labels with the same image, and labels outside the image get weight zero. A set $A\subseteq X$ therefore has induced
weight $\mu'(f^{-1}(A))$. In particular, a retained set $A'\subseteq X'$
has image of induced weight at least $\mu'(A')$.

Under product weights, a tuple has the product of the weights of its coordinates.
Some counting arguments are shorter in probability notation: we regard a label
as drawn according to its weights and use the standard terms distribution,
marginal, and independence. This is an auxiliary calculation, not a random
choice by a communication protocol. An event is a set of labels.
Conditioning on a retained set $A$ of
positive weight means renormalizing: $(\mu\mid A)(B)=\mu(A\cap B)/\mu(A)$.
For each projection below, we specify whether we measure its size with the
original weights or with these conditional weights.

\paragraph{Counting surviving blocks.}
Later proofs need many blocks of a partition to retain enough weight.
We derive this from the total retained weight by the following count.
Suppose we partition an input set into $q$ blocks, each of weight $1/q$, such as the rows $\{i\}\times X$ of $\Int{q}{R}$ in
\cref{def:interlace} below. If a restriction
retains total weight $\beta$, the fractions retained within the blocks sum
to $q\beta$. We set aside at most $e$ blocks and suppose that at most $r$ others
retain fraction $\ge t$. Bounding these $r+e$ fractions by one and
every remaining fraction by $t$ gives
\[
q\beta\le r+e+qt.
\]
Thus, if $q\beta>r+e+qt$, more than $r$ blocks not set aside retain
fraction $\ge t$. If the blocks instead partition only a part of the
input set, of total weight $\lambda$, that part retains fraction at least
$(\beta-(1-\lambda))/\lambda$ of its weight, since the rest of the input set
contributes at most $1-\lambda$ to $\beta$.

\subsection{Interlacing and restrictions}
\label{sec:interlacing}

In an interlace, a row is a row of the base matrix together with an
index $i$, a column is a tuple of base columns, and the entry is read from
the $i$th coordinate of the tuple.

\begin{definition}[Interlacing]\label{def:interlace}
For a Boolean matrix $R:X\times Y\to\{0,1\}$,
the full $q$-fold interlace $\Int{q}{R}$ has rows $(i,u)\in[q]\times X$,
columns $(y_1,\ldots,y_q)\in Y^q$, and entry
\[
\Int{q}{R}((i,u),(y_1,\ldots,y_q))=R(u,y_i).
\]
When $R$ has normalized input weights $\mu_X,\mu_Y$, give row $(i,u)$
weight $\mu_X(u)/q$ and column $(y_1,\ldots,y_q)$ weight
$\prod_{j=1}^q\mu_Y(y_j)$.
\end{definition}

\paragraph{Why column deletion matters.}
Rows of the interlace with different indices $i$ read different
coordinates of Bob's tuple, but this distinction may disappear when
columns are deleted. For example, consider the $2\times2$ base matrix
\[
R=\begin{array}{c|cc}
     & 1 & 2\\ \hline
  a  & 0 & 1\\
  b  & 1 & 0
\end{array}.
\]
The left-hand table in \cref{fig:lookup} is its two-fold interlace.
We obtain the right-hand table by keeping only columns $(1,1)$ and $(2,2)$.
For each fixed $u\in\{a,b\}$, the rows $(1,u)$ and $(2,u)$ then coincide.

\begin{figure}[H]
\centering
\begin{tikzpicture}[>=Stealth,font=\small,
 every node/.style={align=center}]
\matrix (before) at (0,0) [matrix of math nodes,
 nodes={minimum width=9mm,minimum height=6mm,inner sep=1pt},
 column 1/.style={nodes={minimum width=7mm}},
 column 2/.style={nodes={minimum width=6mm}},
 column 4/.style={nodes={fill=black!10}},
 column 5/.style={nodes={fill=black!10}},
 row 1/.style={nodes={font=\footnotesize}}] {
 u&i&(1,1)&(1,2)&(2,1)&(2,2)\\
 a&1&0&0&1&1\\
 a&2&0&1&0&1\\
 b&1&1&1&0&0\\
 b&2&1&0&1&0\\
};
\matrix (after) at (8,0) [matrix of math nodes,
 nodes={minimum width=10mm,minimum height=6mm,inner sep=1pt},
 column 1/.style={nodes={minimum width=7mm}},
 column 2/.style={nodes={minimum width=6mm}},
 row 1/.style={nodes={font=\footnotesize}}] {
 u&i&(1,1)&(2,2)\\
 a&1&0&1\\
 a&2&0&1\\
 b&1&1&0\\
 b&2&1&0\\
};
\node[above=3mm of before,font=\small\bfseries] {Rows may depend on $i$};
\node[above=3mm of after,font=\small\bfseries] {Rows no longer depend on $i$};
\draw[->,semithick] ([xshift=4mm]before.east)
 -- node[above=3mm,font=\footnotesize] {keep $y_1=y_2$}
    ([xshift=-4mm]after.west);
\draw (before-1-1.south west) -- (before-1-6.south east);
\draw (after-1-1.south west) -- (after-1-4.south east);
\draw[black!35] (before-3-1.south west) -- (before-3-6.south east);
\draw[black!35] (after-3-1.south west) -- (after-3-4.south east);
\draw[draw=answergreen,rounded corners=2pt]
 ([xshift=-1mm,yshift=-.6mm]after-2-3.north west) rectangle
 ([xshift=1mm,yshift=.6mm]after-3-4.south east);
\draw[draw=answergreen,rounded corners=2pt]
 ([xshift=-1mm,yshift=-.6mm]after-4-3.north west) rectangle
 ([xshift=1mm,yshift=.6mm]after-5-4.south east);
\node[right=3mm of after,align=left,font=\footnotesize] {identical rows\\for fixed $u$};
\node[below=3mm of before,font=\footnotesize] {Shaded columns are deleted.};
\end{tikzpicture}
\caption{Deleting columns with $y_1\ne y_2$ removes the dependence on
the row index $i$. For each fixed $u$, the rows $(1,u)$ and $(2,u)$ agree
on every retained column.}
\label{fig:lookup}
\end{figure}

The same happens for every base matrix $R$. On the columns
$(y,\ldots,y)$, the entry of $\Int{q}{R}$ is $R(u,y)$ for every row index
$i$. Hence a lower bound for the full interlace need not survive arbitrary
column deletion. For this reason, \cref{def:profile} requires lower bounds only for
restrictions that retain a prescribed weight of rows and of columns.

\paragraph{Keeping fewer column tuples.}
Full interlacing creates $|Y|^q$ columns, too many for an explicit
polynomial-time reduction. We therefore
replace $Y^q$ by a short list with rational weights. For every assignment of
values to at most $D$ coordinates, we preserve the total weight of tuples
matching that assignment. The definition below calls these exact counting
identities $D$-wise independence. \Cref{det:code} constructs such lists
deterministically (\cref{sec:small-lists,sec:deterministic}).

\begin{definition}[Weighted lists of tuples]\label{def:exact-tuples}
Fix integers $q,D\ge1$. For a finite alphabet $Y$ with prescribed
normalized rational weights $\mu$, an indexed tuple list consists of a
finite nonempty label set $\mathcal L$, positive normalized rational
weights $\nu$ on its labels, and a decoder $z:\mathcal L\to Y^q$.
It is \emph{exactly $D$-wise independent} if the weights induced on any
set $J\subseteq[q]$ of at most $D$ decoded coordinates equal the product
weights $\mu^J$. These induced coordinate weights are its marginals.
No independence on larger sets is required. A set of labels may contain one label
but not another with the same decoded tuple.

For a base matrix $R:X\times Y\to\{0,1\}$ with column weights $\mu$,
the interlace on this list has a column for each label $\ell$, with entry
$R(u,z_i(\ell))$ at row $(i,u)$. Rows keep the weights of \cref{def:interlace}, and columns
get the weights $\nu$. For both kinds of interlace, \emph{block} $i$ is the row set
$\{i\}\times X$, and we measure row weight within a block by the base row
weights, without division by $q$.
\end{definition}

The full product is the list with labels $Y^q$, identity decoder and
product weights. For a label set $K$ and a set $J\subseteq[q]$, write
$\pi_JK=\{z(\ell)_J:\ell\in K\}$ for the projection of its decoded tuples.

To apply a lower bound to a few surviving blocks, we project the retained
column tuples onto those blocks. The next observation realizes the
projected tuples as actual columns of the restriction.

\begin{observation}[Restrictions projected to a set of blocks]\label{obs:representatives}
Let $G$ be a restriction of an interlace of $R$, on the full product or on
a tuple list, with retained label set $K$, and let $J$ be a set of blocks.
For each $i\in J$, the rows in block $i$ read only coordinate $i$ of the
decoded tuple. Choose one label of $K$
for each tuple in $\pi_JK$. The restriction of $G$ to these labels and to the
retained rows of the blocks in $J$ is then the full-product interlace on $J$,
restricted to the columns $\pi_JK$ and to the same rows. Every lower bound for this matrix
therefore holds for $G$. Its column weight is $\mu^J(\pi_JK)$, and each
block keeps its base row weights. If the marginal on $J$ is $\mu^J$, as for
full products and for exactly $|J|$-wise independent lists, then
$\mu^J(\pi_JK)\ge\nu(K)$, since $K\subseteq z_J^{-1}(\pi_JK)$. We never use
the weights of the chosen representatives, and we assert no independence of
the retained labels.
\end{observation}


\section{A reusable gap construction}
\label{sec:unrestricted}

In this section we state the one-round amplifier and then iterate it. In
the NO case the amplifier raises the cost by one bit more than in the YES
case, and its output again satisfies the lower bounds under input deletion
that the next round needs.

\subsection{Lower bounds that survive restrictions}

The deletion bounds use two column-weight thresholds, the second half the
first. The reason is that after Bob sends one bit, the heavier branch retains
at least half the column weight. At the lower threshold we allow the bound to
drop by one bit.

\begin{definition}[Deletion bounds]
\label{def:profile}
Let \(F:X\times Y\to\{0,1\}\) have positive normalized rational input weights. For
\(0<\rho\le1\) and an integer \(d\), \(F\) satisfies the
\emph{deletion bounds} \(\RP{F}{\rho}{d}\) if, for every \(j\in\{0,1\}\) and all \(U\subseteq X\), \(V\subseteq
Y\),
\[
\mu_X(U)\ge\rho\ \text{and}\ \mu_Y(V)\ge(3/5)2^{-j} \quad\Longrightarrow\quad
\D(F[U,V])\ge d-j.
\]
Here \(F[U,V]\) is the restriction of \(F\) to \(U\times V\).
No matching upper bound is required.
\end{definition}

Whether the source graph is three-colorable determines on which side of
the gap the constructed matrix lies. For an unsatisfiable formula, we need the graph
to stay non-three-colorable after a small fraction of its vertices is deleted,
since otherwise deleting inputs could erase the lower bound. The reductions
of Dinur and of Bogdanov, Obata and Trevisan provide such
graphs~\cite{dinur2007pcp,bogdanov2002lower}.

\begin{theorem}[Bounded-degree coloring source]
\label{src:gap-colouring}
There are constants \(\Delta\ge1\) and rational \(\gamma>0\), and a
deterministic polynomial-time reduction from \(3\)-SAT to nonempty
loopless multigraphs \(G=(V,E)\) with no isolated vertices and maximum degree at
most \(\Delta\), such that:
\begin{enumerate}
\item a satisfiable formula produces a three-colorable graph;
\item an unsatisfiable formula produces a graph for which every three-coloring has at least \(\gamma|E|\) monochromatic edge occurrences.
\end{enumerate}
For sparse formulas on \(N\ge1\) variables, the reduction can use a fixed
positive integer \(p\) such that, for \(n=|V|\),
\[
N\le |E|,\qquad n,|E|\le N(\log(N+2))^p,\qquad n/2\le |E|\le\Delta n/2.
\]
\end{theorem}

The proof in \cref{app:graph-source} composes these reductions and checks
the stated graph-size bounds and conventions.

We call graphs from satisfiable formulas YES and those from unsatisfiable
formulas NO. We build the amplifier in two ways, which satisfy the same
guarantee. In the \emph{full-product case}, used for the randomized
reduction, every interlace uses all column tuples, and
\cref{lem:full-schedule} fixes the parameters. In the \emph{explicit case},
used for the deterministic reduction, the interlaces use the small lists of
\cref{det:construction}, and \cref{lem:generated-schedule} fixes the
parameters. In each case the last step restores the deletion bounds at full
cost for the next round, by \cref{cl:bridge} in the full-product case and by
\cref{det:restore} in the explicit case.

\begin{theorem}[Amplification under input deletion]
\label{thm:amplifier}
Let \(G=(V,E)\) be a YES or NO graph produced by \cref{src:gap-colouring},
and put \(n=|V|\). In either case there are a public row-weight threshold
\(\rho\), a public integer \(h\), and a deterministic, implicitly represented
construction \(R\mapsto R^+\) on base matrices \(R\) with positive normalized
rational input weights. In the full-product case, \(n\) must exceed an
effectively computable constant, and \(\rho=2^{-L}\) with \(L=\ceil{\log n}\), and
\(h=L+O_{\Delta,\gamma}(\log L)\). In the explicit case, \(\rho=1/\mathsf N\)
and \(h=10\log\mathsf N\), with \(\mathsf N=s^4\) as in
\cref{lem:generated-schedule}. In both cases:
\begin{enumerate}
\item[(i)] On either source promise,
\[
\D(R^+)\le\D(R)+h+1,
\]
and, if \(G\) is YES, the stronger bound \(\D(R^+)\le\D(R)+h\) holds.
\item[(ii)] If \(G\) is NO and the base matrix \(R\) satisfies \(\D(R)=d\ge3\)
and \(\RP{R}{\rho}{d}\), then
\[
\D(R^+)=d+h+1,\qquad \RP{R^+}{\rho}{d+h+1}.
\]
\item[(iii)] On either source promise, \(\rank(R^+)\le2^h\rank(R)\).
\end{enumerate}
\end{theorem}

The proof is assembled at the end of \cref{sec:amplification-proof}
(p.~\pageref{proof:thm:amplifier}). It uses the interlacing estimates of
\cref{sec:interlacing-details} and, in the explicit case, the small lists of
\cref{sec:small-lists}. The rank bound follows from \cref{lem:amplifier-rank}.


\subsection{Iteration}

One round adds $h$ bits to the YES upper bound and $h+1$ bits to the NO
lower bound. If we start from the initial matrix below, of known cost
$d_0$, the two bounds therefore differ by $i$ bits after $i$ rounds.

\begin{lemma}[Initial matrix]\label{cl:seed}
For $L\ge1$, let $p_0=\operatorname{pow}_2(64(L+1))$ and
$R_0(z,j)=z_j$ on $\{0,1\}^{p_0}\times[p_0]$, with uniform input weights. Then
\[
d_0:=\log p_0+1,\qquad \D(R_0)=d_0,\qquad
\rank(R_0)=p_0=2^{d_0-1},\qquad \RP{R_0}{2^{-L}}{d_0}.
\]
For $L=\ceil{\log n}$ and $n\ge4$, the initial matrix has polynomial size in $n$
and can be printed in polynomial time.
\end{lemma}
The proof is in \cref{sec:initial-template-proofs}
(p.~\pageref{proof:cl:seed}).

\begin{proposition}[Iterated amplification]\label{prop:iteration}
Fix one case of \cref{thm:amplifier}, whose threshold is $\rho$ and whose
overhead is $h$, and a source graph to which it applies. Iterate the theorem from the
matrix $R_0$ of \cref{cl:seed} with $L=\log(1/\rho)$, using the same
parameters in every round, and let $R_i$ be the matrix after $i\ge0$ rounds.
Put $B_i=d_0+ih$, $H_i=B_i+i$, and $r_i=2^{B_i-1}$. A YES graph
gives $\D(R_i)\le B_i$. A NO graph gives
$\D(R_i)=H_i$ and $\RP{R_i}{\rho}{H_i}$. On either promise,
$\rank(R_i)\le r_i$ and $\D(R_i)\le H_i$.
\end{proposition}
\begin{proof}
\Cref{cl:seed} gives the case $i=0$, with $d_0\ge3$. By
\cref{thm:amplifier}, each round raises the YES bound by $h$, the NO cost
and the general upper bound by $h+1$, and rank by a factor at most $2^h$.
In the NO case, (ii) restores the deletion bounds at threshold $\rho$, so
the induction continues.
\end{proof}

\section{Proof of amplification under input deletion}
\label{sec:amplification-proof}

We now prove \cref{thm:amplifier}. The template bounds in
\cref{src:profiles} force the initial prefix via \cref{cl:separation};
the collision argument then gives \cref{src:near-full}, the near-full
deletion bound for $N$. For $U$, the template bounds and
\cref{det:conditioning,det:rareboost} give the stopping-node assignments
of \cref{res:destinations}. The transfer bounds
\cref{res:incoming,res:foreign}, together with source hardness, give
\cref{res:alignment}: full cost when column deletion is small, and
deletion bounds one bit below full cost. Finally, \cref{cl:bridge}
for full products and \cref{det:restore} for explicit lists use these
two conclusions to restore the deletion bounds at full cost.

\subsection{One choice of widths and tolerances}

We need enough coordinates to encode the graph, and enough interlaced
copies that the lower bounds still hold after a protocol discards some
inputs. We use one choice of parameters for both cases. Its two scales are $\rho$, the retained
fraction of base rows for which hardness must survive between rounds, and
$s$, the number of bins used to group inputs within a round.

\begin{definition}[Amplification parameters]\label{def:amplification-parameters}
Let $n$ be the source graph size. Choose powers of two $s,\rho^{-1}$ with
$n\le\rho^{-1}$, $\rho\le s^{-4}$, and $s\ge s_0$, the effective constant
of \cref{lem:common-schedule}. Set
\[
\begin{gathered}
q_1=s-2,\qquad q_2=s^2/\rho,\qquad w=s^8,\qquad Q=s^{24},\\
x=\rho^8,\qquad \eta=s^{-2},\qquad
\theta=(512Q)^{-1},\qquad \delta_U=\eta/(16w).
\end{gathered}
\]
The source construction uses $q_1,q_2$ coordinates, and the two further
interlaces have widths $w,Q$. The deletion bounds of \cref{src:profiles}
for the matrices of the source construction use row-weight threshold $x$,
and the lower bound for the source matrix $N$ survives
deleting an $\eta$ fraction of each input set's weight
(\cref{src:near-full}). For the transpose $W$ of the first further interlace
(\cref{def:source-matrix}), $\theta$
is the minimum retained row weight and $\delta_U$ is the allowed fraction of
deleted column weight.
\end{definition}

In the full-product case, we take $L=\ceil{\log n}$, $\rho=2^{-L}$, and
$s=\operatorname{pow}_2(2^{16}(L+1))$. \Cref{lem:full-schedule} shows that
these choices meet the constraints above once $n$ exceeds an effective
constant. In the explicit case, we take $\rho=s^{-4}$ for a sufficiently
large dyadic $s$ with $s^4\ge n$. \Cref{sec:parameters} verifies the
numerical requirements shared by both cases, namely the bounds
\eqref{eq:width-margins}, and the additional conditions of each case. All
these choices are fixed across rounds.

\subsection{Encoding the graph}

We construct the source matrix of \cref{sec:overview} from interlaced copies
of the base matrix $R$ and rows encoding the graph. In a NO instance, every
assignment of items to bins has a \emph{collision}: a coordinate at which one
bin contains two active items and an inactive item. We show in the rest of
the proof that a collision forces the extra bit of the NO case. Because of
the template rows, defined below, every protocol cheaper than the NO-case cost
must begin with $\log s$ bits from Alice, and these bits distinguish the $s$
bins (\cref{cl:separation}). The item rows encode the graph.

\begin{definition}[Source construction]
\label{def:source-matrix}
Let \(R:X\times Y\to\{0,1\}\) have positive normalized rational input weights, and let
\(s\ge4\) be dyadic, meaning a power of two. Coordinates ranging over integer sets such as \([q_2]\) are uniform below,
and product sets carry the product of the specified factor weights.

\smallskip\noindent\textit{Active items.}
We repeat the graph in several groups so that the proof can discard groups
that lose too many items.
An item is a vertex or auxiliary label in one of \(4s\) groups:
\(\mathcal I=[4s]\times(V\sqcup[s-3])\).
In each group, each edge occurrence has a coordinate activating its two
endpoints and all \(s-3\) auxiliary items. All other items are inactive.
Pad these \(4s|E|\) coordinates with inactive coordinates to length
\(q_0=\operatorname{pow}_2(4s|E|)\), then repeat each coordinate
\(q_2/q_0\) times, where \(q_2\) is a dyadic multiple of \(q_0\).
Write \(v_i(\alpha)\) for the bit marking item \(i\) at index \(\alpha\in[q_2]\),
and \(\pi_\alpha\) for the injection of the at most \(s-1\) active items into
\([s]\) that preserves a fixed order on \(\mathcal I\). These bits do not depend on any assignment of items to bins.

\smallskip\noindent\textit{Template interlaces.}
Put \(q_1=s-2\), leaving two message values for active and inactive items.
Define
\[
\Gamma=Y^{2s-1},\qquad J=[q_1]\times X,\qquad S=J^{q_2},\qquad
J_2=[q_2]\times\Gamma.
\]
The template rows use the first \(q_1\) positions of a list
\(\gamma\in\Gamma\), and we reserve its last \(s+1\) positions for the item
rows. The two successive interlace matrices are
\[
A(\gamma,(r,u))=R(u,\gamma_r),\qquad
R_{\rm temp}(z,(\alpha,\gamma))=A(\gamma,z_\alpha).
\]
Thus $A$ has rows indexed by $\Gamma$ and columns by $J$, and $R_{\rm temp}$
has rows indexed by $S$ and columns by $J_2$. Moreover, $A^{\mathsf T}$ is an
interlace of $R$ with repeated column labels,
and $R_{\rm temp}^{\mathsf T}$ is an interlace of $A$. A template label \(z\in S\)
is a tuple of pairs \(z_\alpha=(r,u)\), one for each index.
The pair specifies position $r$ and a base row $u$, giving the
matrix entry \(R(u,\gamma_r)\).

\smallskip\noindent\textit{Matrix entries.}
The matrix \(B_{\rm raw}\) has template rows \((p,z)\), indexed by a bin
\(p\in[s]\) and a template label, and item rows \((i,u)\), indexed by an item
and a base row:
\[
\mathcal K=([s]\times S)\sqcup(\mathcal I\times X).
\]
We give the template rows most of the weight, so that deleting an \(\eta\)
fraction of the row weight leaves most of their weight intact.
Set \(\lambda=255/256\) and \(\mu=1/256\).
Template rows have total weight \(\lambda\), and vertex rows and auxiliary rows
each have total weight \(\mu/2=1/512\). Within each of these three types, bin or item labels are uniform
and the other labels follow their specified weights.
For a column \(Z=(\xi,j_1,\ldots,j_s)\), where \(\xi\in Y^{s+1}\) and
\(j_p=(\alpha_p,\gamma_p)\in J_2\), set
\[
B_{\rm raw}((p,z),Z)=R_{\rm temp}(z,j_p),
\]
\[
B_{\rm raw}((i,u),Z)=
 \begin{cases}
 R(u,\xi_{\pi_\alpha(i)}),
   &\alpha_1=\cdots=\alpha_s=\alpha,\ v_i(\alpha)=1,\\
 R(u,\xi_{s+1}),&\text{otherwise}.
 \end{cases}
\]
Thus an item row uses its assigned position when the indices
\(\alpha_1,\ldots,\alpha_s\) agree and the item is active there, and otherwise
it uses the last entry of the shared list.

\smallskip\noindent\textit{Two column families.}
The first column family \(\Omega=Y^{s+1}\times J_2^s\) has product weights.
The second column family is
\[
\mathcal D = \{(\operatorname{tail}(\gamma),j,\ldots,j):j=(\alpha,\gamma)\in
J_2\}, \qquad \operatorname{tail}(\gamma)=(\gamma_{q_1+1},\ldots,\gamma_{2s-1}),
\]
with the weights induced from \(J_2\). We give the two families
the weights
\[
\begin{array}{ccl}
\text{Family}&\text{Weight}&\text{Role}\\
\Omega&\lambda&\text{Force Alice's initial messages}\\
\mathcal D&\mu&\text{Read the three positions of a collision}.
\end{array}
\]
Thus the first column family has weight \(\lambda\), like the template rows. Keep the two family tags distinct even when their raw columns
coincide. At every bin \(p\), the label \(j_p\) has the weights prescribed on \(J_2\) in
either family. We call the columns of the two families \emph{$\Omega$-inputs} and
\emph{$\mathcal D$-inputs}. After transposition they become rows, and we
then also call them \emph{$\Omega$-rows} and \emph{$\mathcal D$-rows}.

\smallskip\noindent\textit{The two interlaces.}
Transpose the matrix and restrict its rows to the first column family:
\[
N=B_{\rm raw}^{\mathsf T},\qquad
N_{\rm ref}=N[\Omega,\mathcal K],
\]
where the weights on \(\Omega\) are normalized. The matrix \(N_{\rm ref}\) is used only in
the lower-bound analysis. Finally set
\[
U=\Int{w}{N},\qquad W=U^{\mathsf T},\qquad
R^+=(\Int{Q}{W})^{\mathsf T}.
\]
\end{definition}

As in \cref{sec:overview}, the $\Omega$-inputs force Alice's first
\(\log s\) bits to name a bin, and this fixes an assignment of the items to
bins. Against a $\mathcal D$-input with label \((\alpha,\gamma)\), the rows
of two items active at \(\alpha\) and of one inactive item read three
distinct positions of the same list \(\gamma\). A collision puts all three
items in one bin. The interlace \(U\) extends the lower bound for
\(N\) to restrictions that keep only a small fraction of one input set. This
costs one bit in the deletion bounds (\cref{res:alignment}), and the final
interlace restores that bit.

Throughout, we write \(d=\D(R)\). With the parameters of
\cref{def:amplification-parameters}, we define the costs and the threshold
used in the lower-bound analysis (in the explicit case, we drop the
subscript \(\rm g\) that \cref{sec:deterministic} attaches to its
parameters):
\begin{equation}\label{eq:source-costs}
H_1:=d+\log s,\qquad H_2:=H_1+\log q_2,\qquad C:=H_2+\log s+1,\qquad
x_{\rm cap}:=2s\rho.
\end{equation}
We will show that \(A\), \(R_{\rm temp}\) and \(N\) have NO-case lower bounds
\(H_1\), \(H_2\) and \(C\). The collision bound of \cref{src:profiles}
requires each position involved to retain base row weight at least
\(x_{\rm cap}\) within its block.

The following example shows the extra position created by a collision.
Fix \(s=4\) and a coordinate \(\alpha\) at which \(u,v\) are active
and \(h\), from another group, is inactive. Suppose these three items
lie in the same bin, with \(\pi_\alpha(u)=1\) and
\(\pi_\alpha(v)=2\). On a \(\mathcal D\)-input, their rows use
\(\gamma_3,\gamma_4,\gamma_7\), respectively. Together with the
bin's \(s-2=2\) template positions, they expose \(s+1=5\) positions
of the same list.

\begin{figure}[H]
\centering
\begin{tikzpicture}[
  >=Stealth,
  font=\small,
  collisioncell/.style={
    draw=black!65, line width=.45pt,
    minimum width=16mm, minimum height=8.5mm,
    inner sep=0pt, outer sep=0pt
  },
  collisionactive/.style={
    circle, draw=answergreen, fill=answergreen!12,
    line width=.7pt, minimum size=6.5mm, inner sep=0pt
  },
  collisioninactive/.style={
    rectangle, draw=boborange, fill=boborange!7, dashed,
    line width=.7pt, minimum size=6.5mm, inner sep=0pt
  }
]

% One bin: the template rows and all three witness items.
\draw[rounded corners=3pt, line width=.65pt]
  (-.4,1.15) rectangle (11.6,2.55);

\node[font=\small\bfseries] at (5.6,2.92)
  {One bin at a fixed coordinate \(\alpha\)};

\node[
  draw=aliceblue,
  fill=aliceblue!7,
  rounded corners=2pt,
  line width=.65pt,
  minimum width=27mm,
  minimum height=8mm
] (templates) at (1.6,1.78) {Template rows};

\node[text=answergreen] at (4.8,2.28) {Active items};
\node[collisionactive] (u) at (4.0,1.70) {\(u\)};
\node[collisionactive] (v) at (5.6,1.70) {\(v\)};

\node[text=boborange] at (10.4,2.28) {Inactive item};
\node[collisioninactive] (h) at (10.4,1.70) {\(h\)};

% Seven adjacent cells form one list, not seven bins.
\node[collisioncell, fill=aliceblue!10]
  (g1) at (0.8,0) {\(\gamma_1\)};
\node[collisioncell, fill=aliceblue!10]
  (g2) at (2.4,0) {\(\gamma_2\)};

\node[collisioncell, fill=answergreen!12]
  (g3) at (4.0,0) {\(\gamma_3\)};
\node[collisioncell, fill=answergreen!12]
  (g4) at (5.6,0) {\(\gamma_4\)};

\node[collisioncell, fill=black!4, text=black!45]
  (g5) at (7.2,0) {\(\gamma_5\)};
\node[collisioncell, fill=black!4, text=black!45]
  (g6) at (8.8,0) {\(\gamma_6\)};

\node[collisioncell, fill=boborange!12]
  (g7) at (10.4,0) {\(\gamma_7\)};

% The arrows indicate which list positions these rows read.
\draw[->, aliceblue, line width=.7pt]
  (templates.south) -- (g1.north);
\draw[->, aliceblue, line width=.7pt]
  (templates.south) -- (g2.north);

\draw[->, answergreen, line width=.7pt]
  (u.south) -- (g3.north);
\draw[->, answergreen, line width=.7pt]
  (v.south) -- (g4.north);

\draw[->, boborange, dashed, line width=.7pt]
  (h.south) -- (g7.north);

\node at (5.6,-.78)
  {One shared list \(\gamma\) on a \(\mathcal D\)-input};

\node[font=\normalsize] at (5.6,-1.68) {\(
  \underbrace{2}_{\text{template}}
  +\underbrace{2}_{\text{active}}
  +\underbrace{1}_{\text{inactive}}
  =\underbrace{5}_{\text{positions}}>4=s
\)};

\end{tikzpicture}
\caption{A collision exposes \(s+1\) positions, illustrated for \(s=4\).
On a \(\mathcal D\)-input with label \((\alpha,\gamma)\), the template
rows in the depicted bin use positions \(1,2\), the active items
\(u,v\) use positions \(3,4\), and the inactive item \(h\) uses
position \(7\). Activity and the positions read are determined by
\(\alpha\), independently of the bin assignment. Only the positions
used by the displayed rows are highlighted; other items are omitted.}
\label{fig:source-collision}
\end{figure}

The next lemma shows that, for a NO graph, collisions as in
\cref{fig:source-collision} survive deletion. After we delete a few items,
every assignment still has collisions at a constant fraction \(\zeta\) of
the coordinates.

\begin{lemma}
\label{src:item-witness}
A YES graph admits an assignment of \(\mathcal I\) to \([s]\) with at most one
active item in each bin at every coordinate. For a NO graph, delete fewer than
\((s/2)\delta_v n\) vertex items and fewer than \(s/2\) auxiliary items, where
\[
\delta_v=\frac{\gamma|E|}{2\Delta n},\qquad \zeta=\frac{\gamma}{24}.
\]
Every assignment of the remaining items to the \(s\) bins then has at least
\(\zeta q_2\) coordinates at which one bin contains two active items and an
inactive item.
\end{lemma}

We prove the collision guarantee in \cref{sec:collision-conversation}.
First, the YES assignment gives a short protocol.

\begin{proposition}
\label{src:upper}
If \(\D(R)\le y\), then \(\D(N)\le y+2\log s+\log q_2+1\). For a YES graph, the
final \(+1\) can be omitted. Both bounds remain valid after restricting either
input set.
\end{proposition}

\begin{proof}
For a YES graph, fix a public bin assignment $\chi$ from
\cref{src:item-witness}. We describe a protocol for $B_{\rm raw}$ and
exchange the players at the end. Alice holds a template row $(p,z)$ or an
item row $(i,u)$, and Bob holds $Z=(\xi,j_1,\ldots,j_s)$ with
$j_p=(\alpha_p,\gamma_p)$. They exchange:
\[
\begin{array}{lll}
\text{Alice} & \text{bin }p\text{ (use }p=\chi(i)\text{ for an item row)}
  & \log s\text{ bits},\\
\text{Bob} & \alpha_p & \log q_2\text{ bits},\\
\text{Alice} & \text{one of }q_1\text{ template positions, active, inactive}
  & \log s\text{ bits}.
\end{array}
\]
The third message has $q_1+2=s$ possibilities. For a template row, write
$z_{\alpha_p}=(r,u)$ and send code $r$, and the answer is
$R(u,(\gamma_p)_r)$. For an item row, send active precisely when
$v_i(\alpha_p)=1$. If all $\alpha_j$ agree and the code is active, Bob
identifies the unique active item in bin $p$ and uses
$\xi_{\pi_{\alpha_p}(i)}$. Otherwise he uses $\xi_{s+1}$.
Bob can test agreement from his own input. For both row types, the players
then finish by running a cost-$y$ protocol for $R$ once.

For an arbitrary graph, we use any public assignment and replace the active
code by the $s$ possible values of $\pi_{\alpha_p}(i)$. There are now
$q_1+s+1=2s-1$ positions, so the third message costs one additional bit.
The same rules work for every column tag. Exchanging the players gives the
protocol for $N=B_{\rm raw}^{\mathsf T}$, and restriction preserves its cost.
\end{proof}

\subsection{A collision survives the conversation}
\label{sec:collision-conversation}

The lower bound uses the NO promise only through \cref{src:item-witness}. We need its
collision count for every assignment of the retained items, because Alice's
first messages fix the assignment before any index is revealed.

\begin{proof}[Proof of \cref{src:item-witness}]
\phantomsection\label{proof:src:item-witness}
For a YES graph, let \(c:V\to[3]\) be a proper three-coloring given by
\cref{src:gap-colouring}. Assign each vertex \(v\) to bin \(c(v)\), and assign
the \(s-3\) auxiliary labels to the remaining bins. Every edge coordinate
activates that edge's two endpoints and those auxiliaries in one group, so all
its active items occupy distinct bins. Padded coordinates activate no items.

Now suppose that \(G\) is NO. By the deletion hypotheses in
\cref{src:item-witness}, fewer than \(s/2\) groups lose an auxiliary item and
fewer than \(s/2\) groups lose more than \(\delta_v n\) vertices. Hence more
than \(3s\) groups retain every auxiliary item and all but at most \(\delta_v
n\) vertices.

Fix such a group. If two auxiliaries share a bin, they are active together at every
edge coordinate. Otherwise the auxiliaries occupy \(s-3\) distinct bins,
leaving three free bins. Color the retained vertices in those free bins by
their bins, and extend the coloring arbitrarily to all other vertices. The
count below holds for every extension.
At least \(\gamma|E|\) edge occurrences are monochromatic by the NO promise,
and at most \(\Delta\delta_v n=\gamma|E|/2\) meet a deleted vertex.
Every remaining monochromatic edge gives two active items in one bin: an endpoint
in an auxiliary bin shares that bin with its auxiliary, and if neither endpoint is in
an auxiliary bin, the two endpoints share a free bin. Thus this group has
at least \(\gamma|E|/2\) edge coordinates with two active items in one bin.

It remains to find an inactive item in the same bin. For each counted
coordinate, we choose one bin containing the two active items.
We discard a group from the count if one of its chosen bins contains retained
items from that group alone. Every discarded group has an exclusive chosen bin, and two groups
cannot have the same exclusive bin, so at most \(s\) groups are discarded.
In each of the more than \(2s\) remaining groups, every chosen
bin also contains a retained item from another group. That item is inactive
at the chosen coordinate. We have therefore found more than
\(\gamma s|E|\) coordinates with a collision.
Since \(q_0<8s|E|\), their fraction after padding exceeds
\(\gamma/8\), and hence the stated \(\zeta=\gamma/24\).
Uniform repetition to width \(q_2\) preserves that fraction.
\end{proof}

The next estimate is used to prove the template-layer deletion bounds
below for full-product interlaces. For explicit lists we instead use
\cref{det:exact}. Both applications are in the appendix; the collision
argument uses the resulting template-layer bounds.

\begin{lemma}\label{lem:many-product-blocks}
Suppose $\RP{M}{x}{H}$ holds for a rationally weighted matrix $M$, where
$H\ge3$. Let $u\ge16$ be dyadic with $ux\le1$. Every full-product interlace
restriction retaining at least $9u/16$ blocks, each of base row weight at least
$ux$, and a column set of weight at least $2^{-u/1024}$, has cost at least
$H+\log u$.
\end{lemma}

The proof is in \cref{sec:interlace-induction}, on
p.~\pageref{proof:lem:many-product-blocks}.

For the collision argument, we need deletion bounds for $A$, $R_{\rm temp}$
and $N_{\rm ref}$, and a \emph{collision bound}, which says that the template positions
together with the three positions of a collision cost one bit more than $H_1$.

\begin{lemma}[Template-layer deletion bounds]\label{src:profiles}\label{det:profiles}
Consider either case of \cref{thm:amplifier}: the full-product construction of
\cref{def:source-matrix} with the parameters of \cref{lem:full-schedule}, where
$\rho=2^{-L}$, or the explicit construction of \cref{det:construction} with the
parameters of \cref{lem:generated-schedule}, where $\rho=\rho_{\rm g}=1/\mathsf N$
and the subscript $\rm g$ is suppressed. Suppose $\D(R)=d\ge3$ and $\RP{R}{\rho}{d}$,
let $x=\rho^8$ be the block threshold, and define $H_1,H_2,C,x_{\rm cap}$ by
\eqref{eq:source-costs}. Then
\[
\RP{A}{x}{H_1},\qquad \RP{R_{\rm
temp}}{x}{H_2},\qquad \RP{N_{\rm ref}}{x}{C-1}.
\]
In the interlace of $R$ on $\Gamma$ (on $\Gamma_{\rm g}$ in the explicit case),
the first $q_1$ positions together with any three distinct tail positions
require at least $H_1+1$ bits whenever every position retains base row weight
at least $x_{\rm cap}$ within its block and the retained set of list labels
has weight greater than $1/2$.
\end{lemma}

The proof is in \cref{sec:initial-template-proofs}, after the initial-matrix
proof (p.~\pageref{proof:src:profiles}).

Three proofs below follow retained input sets along a protocol prefix and
count them where the prefix stops. A mark lets a set count twice until it
splits into two sets. In the collision application, the mark records the
extra lower bound supplied by a collision.

\begin{lemma}[Kraft count]\label{lem:kraft-count}
Let $q$ be a power of two, and let $V$ be a set of nodes of a binary protocol
tree, none an ancestor of another, all at depth at most $\log q$. For $v\in V$
at depth $r_v$ put $\kappa_v=q2^{-r_v}$. Suppose each node $v\in V$ holds $h_v$
objects, $z_v\le h_v$ of them marked, and $\sum_{v\in V}(h_v+z_v)>q$. Then some $v$
has $h_v+z_v>\kappa_v$. For this $v$, if $\kappa_v\ge2$ then
$h_v\ge\kappa_v/2+1$, and if $\kappa_v=1$ then $h_v\ge2$ or $h_v=z_v=1$.
\end{lemma}
\begin{proof}
Kraft's inequality gives $\sum_v\kappa_v\le q$, so some $v$ has
$h_v+z_v>\kappa_v$, and $z_v\le h_v$ gives $2h_v>\kappa_v$. If
$\kappa_v\ge2$, then $\kappa_v$ is even and $h_v\ge\kappa_v/2+1$. If
$\kappa_v=1$, then $h_v+z_v\ge2$, so $h_v\ge2$ or $h_v=z_v=1$.
\end{proof}

In the applications, the objects are retained input sets. We follow them
through the messages of one player, starting from sets of weight at least
$t$, by the following \emph{tracking rule}, which keeps every tracked set at
weight at least $t2^{-r}$ after $r$ such messages. An unmarked set follows a child retaining at least half its weight.
A marked set splits into two unmarked sets if both children meet the next
threshold, and otherwise it follows a child meeting that threshold and keeps
its mark. Thus every mark splits at most once. Before it splits, a marked
set loses less than $\sum_{r\ge0}t2^{-r-1}=t$ of its weight. A split
replaces the contribution $2$ of a marked set to $\sum_v(h_v+z_v)$ by $1+1$. Hence
$\sum_v(h_v+z_v)$ equals the number of initial sets plus the number of
initial marks.

\begin{lemma}[Following collision witnesses]\label{lem:collision-tracking}
Let $q\ge2$ be dyadic, let $0\le b<q$ be an integer, and let $x>0$
satisfy $2qx\le1/8$.
Suppose a protocol with at most
$H+\log q$ bits left has $q-b$ indices, each with a retained set of
list inputs of weight at least $3/4$. Each set consists of inputs with that
fixed index. Mark $w_0>b$ of these sets.
Assume the following lower bounds, where all set weights are computed from
the original input weights.
\begin{enumerate}
\item[(i)] Until Alice speaks, a set of retained lists at one index with weight at least $x$, against the
unchanged set of Alice's inputs, costs at least $H>0$.
\item[(ii)] At the first Alice bit, one child retains a set of Alice's inputs on which,
for every dyadic $u\ge2$ with $ux\le1$, retained sets at any $u/2+1$ distinct indices, each
of weight at least $ux$, cost at least $H+\log u$.
\item[(iii)] A marked set retaining weight greater than $1/2$ costs at least
$H+1$ before Alice speaks.
\end{enumerate}
Then no such protocol exists.
\end{lemma}
\begin{proof}
We track the retained list sets of each index through Bob bits by the
tracking rule with $t=2qx$, and we stop at the first Alice bit. A marked set loses less
than $2qx\le1/8$ before it splits, so an unsplit mark retains weight greater
than $1/2$. By (i), a tracked set cannot end at a leaf, nor reach Bob depth
$\log q+1$: it would still have weight at least $x$ but only $H-1$ bits left.
The stopping nodes therefore satisfy the hypotheses of \cref{lem:kraft-count},
where $v$ holds $h_v$ tracked sets and $z_v$ unsplit marks, and
$\sum_v(h_v+z_v)=q-b+w_0>q$. Let $v$ be the node given by that lemma, with
$\kappa=\kappa_v$. Its sets have distinct indices, because descendants of a
split never rejoin, and each has weight at least $2\kappa x$.

If $\kappa\ge2$, then (ii) with $u=\kappa$ requires $H+\log\kappa$ bits after
Alice's bit, but only $H+\log\kappa-1$ remain. If $\kappa=1$ and $h_v\ge2$, use
(ii) with $u=2$ and the available weights $2x$. If $h_v=z_v=1$, then by (iii)
the unsplit marked set needs $H+1$ bits at $v$, where only $H$ remain.
\end{proof}

\begin{theorem}[Source hardness after deletion]\label{src:near-full}
Use either the full-product construction of \cref{def:source-matrix} or the
explicit-list construction of \cref{det:construction}, with a NO graph and
$\D(R)=d\ge3$, and define $H_1,H_2,C,x_{\rm cap}$ by \eqref{eq:source-costs}.
Suppose the deletion bounds and the collision bound stated in
\cref{src:profiles} hold. Assume the numerical inequalities
\eqref{eq:width-margins} of \cref{lem:common-schedule}, that the weights on $S$ give exact eight-wise
independence with the prescribed weights on $J$ as marginal, and that the first column family's bin tuple has
exactly eight-wise independent weights with the prescribed $J_2$ marginal.
Then $\D(N)=C$, and every restriction retaining at least $1-\eta$ of each
input set's total weight has cost at least $C$. In either case of
\cref{thm:amplifier} these hypotheses follow from \cref{src:profiles},
\cref{lem:common-schedule}, and the exact marginals of the construction.
\end{theorem}
\begin{proof}
We first outline the proof. Alice's first $\log s$ bits assign the items
to bins. The collisions of this assignment give the marked sets of
\cref{lem:collision-tracking}, and \cref{src:profiles} gives its three
lower bounds. That lemma then rules out depth $C-1$.

Suppose a restriction of $N$ retaining at least $1-\eta$ of each input set's
weight has a protocol of depth at most $C-1$.
We transpose it to a protocol for $B_{\rm raw}$. Each template bin loses at most
$\eta s/\lambda<1/1000$ of its total $S$ weight. The first column family $\Omega$
retains normalized weight at least $1-\eta/\lambda>99/100$.
Template entries ignore the tail \(\xi\). We therefore project the retained
$\Omega$-inputs onto their bin tuple $(j_1,\ldots,j_s)$ (in the explicit
case, onto their $\mathcal J_{\rm g}$ label) and choose one retained
completion for each image label. These representatives form a restriction of an
interlace of $R_{\rm temp}$ on the template rows, with column weight at
least $99/100$ (\cref{obs:representatives}).

We apply \cref{cl:separation} to the protocol inherited on these
representatives, without pruning its tree. We take $M=R_{\rm temp}$,
$d=H_2$, $u=s$, $x$ in place of its $\rho$, and row threshold $999/1000$ in
place of its $\xi$. Its hypotheses follow
from the two deletion bounds $\RP{R_{\rm temp}}{x}{H_2}$, the depth
$C-1=H_2+\log s$, and the numerical inequalities
$999/1000>2sx$ and $999/1000-2(s-1)x>998/1000$.
By that theorem, every binary history of length below $\log s$ is reached by
template inputs, and every such node is Alice's. They are actual nodes of the
original protocol, so Alice sends its first $\log s$ bits on all retained
inputs, including item rows and $\mathcal D$-inputs. At the resulting $s$
nodes, the respective template bins retain sets $A_1,\ldots,A_s\subseteq S$,
each of weight greater than $998/1000$.
We needed the representatives only to establish this prefix. From now on we use
all originally retained columns, with $C-1-\log s=H_2$ bits left.

We next fix the item assignment before choosing an index. We set aside
items whose rows retain less than half of their base row weight in the
restriction. By the row weights in \cref{def:source-matrix}, each item type has
total weight \(1/512\), so we set aside fewer than \(4096\eta sn\) vertex
items and fewer than \(4096\eta s(s-3)\) auxiliary items. These numbers are below
the deletion thresholds of \cref{src:item-witness}, because \eqref{eq:width-margins}
gives \(4096\eta sn\le\delta_v n/16\) and \(4096\eta s(s-3)<s/2\), as
checked in the proof of \cref{lem:common-schedule}. By averaging over these \(s\) nodes, we can assign every remaining item to
a node that receives at least \(1/(2s)\) of its base weight. The parameter choices give
\(4s^2\rho\le1\), so
\[
\frac1{2s}\ge2s\rho=x_{\rm cap}.
\]
Thus the assigned part retains at least \(x_{\rm cap}\) of the item's base row
weight. This assignment is fixed before the index is chosen. By
\cref{src:item-witness}, at least \(\zeta q_2\) coordinates contain two active
items and an inactive item in one bin. We assign each such coordinate to one node \(j\in[s]\)
at which its collision occurs and write \(U_j\subseteq[q_2]\) for the resulting
disjoint sets.

For each \(j\), we discard two kinds of indices \(\alpha\): those at which
the template labels in \(A_j\) project onto too little base row weight, and
those at which too little $\mathcal D$-weight is retained. First, we discard
\(\alpha\) for \(A_j\) if one of the
\(q_1\) blocks of the \(J\)-projection of \(A_j\) at \(\alpha\) retains base row
weight less than \(x_{\rm cap}\). Since \(x_{\rm cap}\le1/8\) by
\eqref{eq:width-margins} and each block has weight \(1/q_1\) in \(J\), the
\(J\)-projection then misses more than \(7/(8q_1)\) of the weight of \(J\).
If there were at least $8q_1$ such indices, the number $Z$ of $8q_1$ of them at which
$z_\alpha$ lies outside this projection, for $z\in S$ drawn from its weights,
would satisfy $\mathbb EZ>7$ and, by pairwise independence,
$\operatorname{Var}Z\le\mathbb EZ$. Every label in $A_j$ has $Z=0$, so
Chebyshev's inequality would give $\Pr(A_j)\le\Pr(Z=0)<1/7$, contrary to
$\Pr(A_j)>998/1000$. Thus fewer than \(8q_1\) indices are discarded for \(A_j\).
Second, we discard \(\alpha\) if the $\mathcal D$-inputs with index \(\alpha\)
retain less than \(3/4\) of their weight. The $\mathcal D$-inputs have total
weight \(\mu\), and their index is uniform on \([q_2]\). Each such \(\alpha\)
therefore loses more than \(\mu/(4q_2)\) of the column weight. Since the
restriction deletes at most \(\eta\) of the column weight, we discard fewer
than \(4\eta q_2/\mu\) such indices. Let \(B_j\) be the union of these two
sets of discarded indices and put \(b_j:=|B_j|\). The inequalities in
\eqref{eq:width-margins} give
\[
\sum_{j=1}^s b_j<8sq_1+4s\eta q_2/\mu<\zeta q_2/4.
\]
Since \(\sum_j|U_j|\ge\zeta q_2\), some \(j\) satisfies \(|U_j|>4b_j\). Fix this
\(j\), write \(b:=b_j\), and define \(w_0:=|U_j\setminus B_j|\). Then
\(w_0\ge|U_j|-b>3b\).

Fix a collision index \(\alpha\in U_j\setminus B_j\). Let
\(i_1,i_2\) be two active items and \(i_3\) an inactive item in one bin at
\(\alpha\), as given by \cref{src:item-witness}. We choose retained template
rows for the first \(q_1\) projected blocks and retained rows of \(i_1,i_2,i_3\).
By the entry rule in \cref{def:source-matrix}, against a
$\mathcal D$-input with label \((\alpha,\gamma)\) these rows read the
following entries, where \(u\) is the base row of the respective row:
\[
R(u,\gamma_\ell)\quad(\ell\in[q_1]),\qquad
R(u,\gamma_{q_1+\pi_\alpha(i_1)}),\qquad
R(u,\gamma_{q_1+\pi_\alpha(i_2)}),\qquad R(u,\gamma_{q_1+s+1}).
\]
The three tail positions are distinct by the collision definition and the tail
convention in \cref{def:source-matrix}. Each selected row group has base weight at
least \(x_{\rm cap}\), because \(\alpha\notin B_j\) and the item assignment
above was fixed with that threshold. The selected rows depend only on the
projected labels and not on \(\gamma\), so they form a restriction of the
interlace of \(R\) on \(\Gamma\) in \cref{src:profiles}. The collision bound of
\cref{src:profiles} therefore requires \(H_1+1\) bits whenever the retained
$\mathcal D$-inputs at index \(\alpha\) have weight greater than \(1/2\) in \(\Gamma\).

We now apply \cref{lem:collision-tracking} with $q=q_2$ and $H=H_1$.
Its $q_2-b$ indices are those outside $B_j$, and its $w_0>3b$ marked sets are
those at the indices in $U_j\setminus B_j$. The numerical bounds
\eqref{eq:width-margins} give $q_2x\le10^{-4}$, hence $2q_2x\le1/8$.
We now verify its hypotheses (i), (ii) and (iii).

Until Alice speaks, the set $A_j$ of template labels is unchanged. Its projection onto
each $J$-coordinate has weight at least $\Pr(A_j)>998/1000>3/5$, so (i)
follows from the deletion bounds of $A$ (\cref{obs:representatives}).
At the first Alice bit, one child retains weight greater than
$499/1000>9/25$. After transposition, the entries of that child at indices with
retained list sets are $A(\gamma,z_\alpha)$, so they form the interlace of $A$ with
those sets of lists as row blocks and tuple list $S$. Since $S$ is exactly
eight-wise independent and this child retains template-label weight greater
than $9/25$, \cref{det:exact} with $M=A$, $d=H_1$ and $2^a=u$ gives cost at
least $H_1+\log u$ for any $u/2+1$ retained index sets of weight at least
$ux$. This gives (ii), with representatives retained at that child. The collision bound
established above gives (iii), which applies before that Alice bit. Hence
\cref{lem:collision-tracking} excludes depth $C-1$.
The upper-bound protocol of \cref{src:upper} has depth $C$, which proves
equality.
\end{proof}

\subsection{Lower bounds after restricting tuple inputs}
\label{sec:retained-tuples}

Restricting the tuple inputs can destroy the independence among their
coordinates (\cref{sec:overview}), so we use independence only in the original
tuple weights. Let $0<v\le1$ be a lower bound on the weight of an event, and
let $0<\tau\le1$ bound the allowed error. The next lemma shows that, after
conditioning on the event, every joint marginal on at most eight coordinates
is within total variation distance $\varepsilon(v,\tau)$ of the product
weights, provided these coordinates avoid a set of at most $e(v,\tau)$
exceptional coordinates. Set
\[
\begin{aligned}
\varepsilon(v,\tau)&:=\min\{\tau,2^{-30}v^{1/256}\},\\
e(v,\tau)&:=8\ceil{2^{20}\varepsilon(v,\tau)^{-2}v^{-1/512}},
\qquad e(v):=e(v,1).
\end{aligned}
\]
We use independence order $8192=8\cdot1024$ so that any $1024$ events on
pairwise disjoint sets of at most eight coordinates are independent under
the original weights, as a $1024$th-moment bound requires.

\begin{lemma}[Conditioning outside exceptional coordinates]\label{det:conditioning}
Let $\nu$ be the normalized weights of an indexed exactly $8192$-wise independent tuple list
(\cref{def:exact-tuples}) on $Z^q$ with marginal $\mu$, and let $0<v,\tau\le1$. For every label event
$Y$ with $\nu(Y)\ge v$, there is a set $E\subseteq[q]$, with
$|E|\le e(v,\tau)$, such that
\[
d_{\rm TV}\bigl((\nu\mid Y)_J,\mu^J\bigr)\le\varepsilon(v,\tau)
\qquad
(J\subseteq[q]\setminus E,\ |J|\le8).
\]
Here total variation is the maximum absolute difference between the two
distributions' probabilities of an event. The reduction need not compute the
event or the exceptional set.
\end{lemma}

The proof is in \cref{sec:conditioning-moments}.

Combined with \cref{obs:representatives}, this bounds the column weight of
projected restrictions after conditioning. If $K\subseteq Y$ and $J$
consists of at most eight coordinates outside $E$, then the inverse image of
$\pi_JK$ contains $K$, so
\[
\mu^J(\pi_JK)
\ge(\nu\mid Y)_J(\pi_JK)-\varepsilon(v,\tau)
\ge(\nu\mid Y)(K)-\varepsilon(v,\tau).
\]
Thus the projected restriction of \cref{obs:representatives} has column weight
at least $(\nu\mid Y)(K)-\varepsilon(v,\tau)$, and we need no product assumption on
$\nu\mid Y$.

The next two results apply to full products and to explicit lists. They
share a proof in \cref{sec:surviving-block-proofs}
(p.~\pageref{proof:det:exact}).

\begin{lemma}[Surviving blocks in a tuple list]\label{det:exact}
Suppose $d\ge2$ and $\RP{M}{x}{d}$. Let the tuple list for an interlace be
exactly eight-wise independent with the base-column marginal. If an integer
$a\ge1$ satisfies $2^ax\le1$, and a restriction retains at least
$2^{a-1}+1$ blocks, each of base row weight at least $2^ax$, together with a
tuple-label set of total weight at least $9/25$, then the restricted
interlace has complexity at least $d+a$.
\end{lemma}

After a protocol prefix, the retained tuple labels may weigh less than the
$9/25$ of \cref{det:exact}. We then condition on the event $Y$
selected by the prefix and discard its exceptional coordinates, as follows.

\begin{theorem}[Surviving blocks after conditioning]\label{det:rareboost}
Suppose $d\ge2$, $\RP{M}{x}{d}$, and an exactly $8192$-wise independent
tuple list has weights $\nu$ and the base-column weights as marginal. Let $\nu(Y)\ge v>0$,
choose the exceptional set from \cref{det:conditioning} with any error bound
$0<\tau\le1$, and let an integer $a\ge1$ satisfy $2^ax\le1$. If a restriction
retains at least $2^{a-1}+1$ blocks indexed outside that exceptional set, each of base row weight at least
$2^ax$, and a column set $K\subseteq Y$ with $(\nu\mid Y)(K)\ge1/2$, then
its complexity is at least $d+a$.
\end{theorem}


\subsection{Where displaced inputs can go}\label{sec:internal-interlace}

We now use $w$ copies of $N$ to obtain a lower bound when only a small
fraction of one input set remains, following the copies through a
hypothetical short protocol as in \cref{sec:overview}. The argument applies to
both constructions. In $U=\Int{w}{N}$ each copy has row weight $1/w$, but we
measure transferred weight in units of one complete copy, and losses among a
copy's $\Omega$-rows in weights normalized within $\Omega$. Recall that
$\Omega$ has weight $\lambda=255/256$ and $\mathcal D$ has weight $\mu=1/256$.

\paragraph{Common hypotheses.}
In this section we use only the following hypotheses, where $H_1,H_2,C$ are as in
\eqref{eq:source-costs} and $H_1,H_2\ge2$.
To bound arriving weight in \cref{res:incoming}, we will fix the
type/bin/item sequence and the base-label sequence by averaging.
The independence in (H3) leaves the template-label weights unchanged
until we condition on the retained template inputs.
\begin{enumerate}
\item[(H1)] The per-copy deletion bounds $\RP{A}{x}{H_1}$,
$\RP{R_{\rm temp}}{x}{H_2}$ and $\RP{N_{\rm ref}}{x}{C-1}$.
\item[(H2)] The matrix $N$ has exact cost $C$, including on every restriction
retaining $1-\eta$ of each input set's total weight.
\item[(H3)] The weights on $S$ give exact $8192$-wise independence across its
$J$-valued coordinates, and the bin tuple in $\Omega$ has the same
independence order and the prescribed $J_2$ weights as marginal.
The tuple weights $\nu$ in $U$ have
the original weights on $\mathcal K$ at each coordinate and a representation by three independent
sequences: type/bin/item labels, base labels with the $X$ weights, and
template labels with the $S$ weights. Each sequence has a value at
every copy and is exactly $8192$-wise independent across copies. The decoder uses
$(p,z_i)$ at a template position and $(a,u_i)$ at an item position.
Thus the decoded labels in $\mathcal K$ at any at most $8192$ copies have the product of the original $\mathcal K$ weights.
The tuple coordinates \(z_{i,\alpha}\), indexed by copy $i$ and coordinate
$\alpha$, are also exactly $8192$-wise independent: group at most
$8192$ coordinates by copy, then use independence between copies and among
the tuple coordinates within each template input. By independence, fixing some of the three complete
sequences leaves the remaining ones with their original weights.
\item[(H4)] The widths and deletion thresholds satisfy the numerical bounds in
\eqref{eq:width-margins}: each quantity in the first column is below the
corresponding bound in the last column. These bounds provide enough coordinates
to encode the graph and ensure that, after deleting inputs and setting aside
exceptional coordinates, enough inputs remain for the lower-bound arguments.
\end{enumerate}
In both cases, (H1) follows from \cref{src:profiles} and (H2) from
\cref{src:near-full}. \Cref{lem:common-schedule} verifies (H4) for the parameter
choices of \cref{def:amplification-parameters}. For (H3) with full products, we take all
complete sequences, including values ignored by the decoder, with product
weights across the three sequences. Decoding induces the original product
weights on $\mathcal K^w$, so pulling back preserves the weights of retained
sets and the entries and depth of protocols. We never need to enumerate
this auxiliary representation. In the explicit case the complete sequences are the
actual indexed input labels (\cref{det:construction,det:flat}).

By \cref{det:conditioning} with $v=\theta$ and $\tau=\eta/2$, conditioning on
a tuple event of weight at least $\theta$ makes at most $e(\theta,\eta/2)$
copies exceptional. Write $E_{\rm exc}=2e(\theta,\eta/2)$. The reduction need
not identify the exceptional copies.

We stop each path just before Bob's first bit or just after Alice's
$(\log w)$-th bit, whichever comes first. Up to these
\emph{stopping nodes} only Alice has spoken, so each row follows the same path
for every remaining column. By following the $\Omega$-rows, we assign distinct
stopping nodes to all but at most $E_{\rm exc}$ copies (\cref{res:destinations}).

\begin{lemma}\label{res:destinations}
Under the common hypotheses above, suppose a restriction of $U$ retains row
weight at least $1-\delta_U$ and a tuple event $Y'$ of original weight at least
$\theta$, with a protocol of depth at most $C+\log w-1$.
Apply \cref{det:conditioning} to this one event with error bound $\eta/2$, removing
at most $E_{\rm exc}/2$ copies. At an Alice node of depth $r$, put
$\kappa=w2^{-r}$ and follow each remaining set of $\Omega$-rows from one copy,
of normalized weight at least $2\kappa x$, into every child retaining weight
at least $\kappa x$. We count these row sets as follows. An unsplit row set
contributes one. Once a row set first splits into two children meeting the
threshold, each descendant contributes one-half. We add these
contributions at each node. They count row sets and are not input weight.

Stop at the first node where Bob speaks or at depth $\log w$, whichever comes first.
The latter stopping condition applies regardless of the next speaker.
Complete paths without row sets by fixed Alice bits. At a node where Bob
first speaks with $\kappa\ge2$, the total contribution is at most $\kappa/2$.
A node at depth $\log w$ has at most one row set. At least $w-E_{\rm exc}$
nodes at that depth are reached by unsplit row sets from distinct copies. Assign each
such node to the copy of its row set. Over all other stopping nodes,
the sum of $\kappa$ is at most $E_{\rm exc}$.
\end{lemma}
\begin{proof}
Since $\delta_U=\eta/(16w)$, initial deletion removes at most
$w\delta_U=\eta/16$ complete copies in total, hence at most $\eta/16$ of any
single copy, so every copy starts with normalized weight on $\Omega$ at least
$1-\eta/(16\lambda)>2wx$, by \eqref{eq:width-margins}. A row set always has a child retaining at
least half its input weight. As with the marks in the tracking rule after
\cref{lem:kraft-count}, the first split of a row set preserves the total count,
since $1=\frac12+\frac12$, and further splits can only increase it. Descendants of a copy
never rejoin. Each of the at least $w-E_{\rm exc}/2$ copies outside the
exceptional set starts with an unsplit row set, so the total count at the
stopping nodes is at least $w-E_{\rm exc}/2$.

Until Bob speaks, the tuple set is the fixed $Y'$. For every copy whose
coordinate lies outside the exceptional set,
its projection has original $\mathcal K$ weight at least $1-\eta/2>3/5$
(the display after \cref{det:conditioning} with $K=Y=Y'$) and gives a restriction of $N_{\rm ref}$
(\cref{obs:representatives}). Since a row set
has row weight $\ge2x$, the bound
$\RP{N_{\rm ref}}{x}{C-1}$ therefore requires at least $C-1>0$ further bits,
and no row set reaches a leaf.
After we pad paths without row sets, every root-to-leaf path has one
stopping node, so the sum of $\kappa$ over these nodes is $w$.

At a node where Bob first speaks with $\kappa\ge2$, suppose there were at
least $\kappa/2+1$ row sets. They come from distinct copies whose coordinates
lie outside the exceptional set for the same event $Y'$. The Bob child retaining more column weight has conditional weight at least one-half.
We apply \cref{det:rareboost} to $N_{\rm ref}$ and this child with
$a=\log\kappa$. Here $a$ is an integer, $\kappa x\le wx\le1$, and
each row set has weight $2\kappa x\ge\kappa x$. The child requires
$C-1+\log\kappa$ bits, but only $C+\log\kappa-2$ remain.
Consequently there are at most $\kappa/2$ row sets, and their total
contribution is at most $\kappa/2$.

At a node at depth $\log w$, two row sets of weight $2x$ would require $C$ bits by the same theorem
with $a=1$ and all of $Y'$, but only $C-1$ bits remain.
Thus a node at depth $\log w$ has at most one row set. If the set is unsplit, we
assign the node to the copy supplying it. Otherwise, the set's
contribution is at most one-half. If $c$ is the number of assigned nodes,
the count at every other node is at most half its value of $\kappa$,
whence
\[
 w-E_{\rm exc}/2\le c+\tfrac12(w-c),\qquad w-c\le E_{\rm exc}.
\]
Each assigned node has $\kappa=1$. An unsplit row set follows only one path,
so the assigned nodes correspond to distinct copies. The sum of $\kappa$
over all other stopping nodes is bounded as claimed.
\end{proof}

Rows displaced from their copy's assigned node must reach either an
unassigned node or a node assigned to another copy. We use the first bound
for unassigned nodes and the second for nodes assigned to another copy.
Their proofs are in
\cref{sec:incoming-bound-proof,sec:foreign-bound-proof}, respectively.

\begin{lemma}\label{res:incoming}
For the protocol and stopping nodes of \cref{res:destinations}, the total
weight from $\mathcal D$ across all copies reaching a node where Bob first
speaks before depth $\log w$ is less than $2s\kappa$. At a node at depth
$\log w$ it is less than $2s$.
Weight is measured in units of one complete copy of $N$.
\end{lemma}

\begin{lemma}\label{res:foreign}
In the same protocol, for each copy $i$ with an assigned node, let
$\varepsilon_i$ be the normalized weight of copy $i$'s $\Omega$-rows missing
at that node. Every
other copy that has an assigned node sends at most $\mu\delta_i$ weight from
$\mathcal D$ to the node assigned to copy $i$, where
\[
 \delta_i=2q_2x+4\varepsilon_i.
\]
\end{lemma}

These bounds lead to the contradiction shown in
\cref{fig:reference-destinations}, in which $P_i$ is the $\mathcal D$-weight of
copy $i$ surviving outside its assigned node and $a_i$ is its initially
deleted $\mathcal D$-weight.

\begin{figure}[H]
\centering
\begin{tikzpicture}[
  >=Stealth,font=\footnotesize,
  box/.style={draw,rounded corners=2pt,align=center,inner sep=5pt,minimum height=12mm},
  src/.style={box,text width=28mm},
  dst/.style={box,text width=38mm},
  flow/.style={->,semithick}]
\node[src] (origin) at (0,0) {Rows from $\mathcal D$\\in copies with assigned nodes};
\node[src,draw=black!55] (deleted) at (0,-1.75)
  {Initially deleted\\$\sum_i a_i\le\eta/16$};
\node[dst,draw=black!55] (own) at (4.7,1.6)
  {Stay at their assigned node\\not counted in $P_i$};
\node[dst,fill=aliceblue!6,draw=aliceblue] (outgoing) at (4.7,-.4)
  {Reach other nodes\\[3pt] $\sum_iP_i>\eta w/2$};
\node[dst] (other) at (10.1,.6)
  {Other assigned nodes\\total $\le\eta w/32$};
\node[dst] (exceptional) at (10.1,-1.4)
  {Nodes not assigned\\to a copy\\total $\le\eta w/16$};
\coordinate (firstsplit) at (2.1,0);
\draw[semithick] (origin.east) -- (firstsplit);
\draw[flow,black!55] (origin.south) -- (deleted.north);
\draw[flow,black!55] (firstsplit) |- (own.west);
\draw[flow] (firstsplit) |- (outgoing.west);
\coordinate (secondsplit) at (7.4,-.4);
\draw[semithick] (outgoing.east) -- (secondsplit);
\draw[flow] (secondsplit) |- (other.west);
\draw[flow] (secondsplit) |- (exceptional.west);
\node[box,fill=black!4,text width=127mm] at (4.7,-3.75)
  {\textbf{There is nowhere for all the outgoing weight to go.}\\[3pt]
   Required outgoing weight $>\eta w/2
   \;>\;3\eta w/32=\eta w/32+\eta w/16$
   (total receiving bound).};
\end{tikzpicture}
\caption{Where the displaced $\mathcal D$-weight would have to go. Every surviving $\mathcal D$-row outside the node assigned to its copy must reach one of the two
groups of stopping nodes on the right, whose combined bound on arriving weight
is too small. The second group contains nodes not assigned to copies, either
where Bob first speaks before depth $\log w$ or at that depth. Rows deleted
initially and rows staying at the node assigned to their copy are counted separately.
The proof of \cref{res:alignment} derives these bounds. All weights are in
complete-copy units, and arrow widths are not to scale.}
\label{fig:reference-destinations}
\end{figure}

For $W$, the deletion bounds we obtain are one bit below the full cost. The next
subsection recovers that bit.

\begin{theorem}[Lower bounds for $W$]\label{res:alignment}
Under the common hypotheses of \cref{sec:internal-interlace},
\[
 \D(W)=C+\log w,\qquad \RP{W}{\theta}{C+\log w-1}.
\]
Every restriction of $W$ retaining row weight at least $\theta$ and column
weight at least $1-\delta_U$ has cost at least $C+\log w$.
\end{theorem}
\begin{proof}
Suppose a restriction of $W$ with row weight $\ge\theta$ and column
weight $\ge1-\delta_U$ has a protocol of depth $\le C+\log w-1$, and
transpose the restriction and its protocol to $U$.
Its retained tuple event $Y'$ has weight at least $\theta$.
By \cref{res:destinations} and $E_{\rm exc}<w/4$ from \cref{lem:common-schedule},
at least $w-E_{\rm exc}\ge3w/4$ copies
have assigned nodes, each with only $C-1$ bits left. These copies lie outside
the exceptional set, so conditioning on $Y'$ with
error bound $\eta/2$ gives each of them a column projection of original
weight at least $1-\eta/2$, which gives a restriction of $N$
(\cref{obs:representatives}). By \cref{src:near-full}, the assigned node must miss more than
$\eta$ of the row weight in its assigned copy.

An unsplit set of $\Omega$-rows discards less than
$\sum_{r=0}^{\log w-1}wx2^{-r}=2(w-1)x$ on its path, since at depth $r$
the discarded child retains less than $\kappa x=wx2^{-r}$.
Initial deletion removes at most $\eta/16$ complete copies. Summing over copies
with assigned nodes and using $w^2x/\eta\le1/(32\lambda)$ from
\eqref{eq:width-margins}, we obtain
\[
 \sum_i\varepsilon_i\le\frac{\eta}{16\lambda}+2w(w-1)x
 \le\frac{\eta}{8\lambda}.
\]
For these copies, the missing $\Omega$-rows, initially
deleted $\mathcal D$-rows, and $\mathcal D$-rows surviving elsewhere partition
the part absent from their own assigned node. Thus
$P_i+a_i+\lambda\varepsilon_i>\eta$, and
\[
\sum_iP_i>\frac{3\eta w}{4}-\frac{\eta}{16}-\frac{\eta}{8}
>\frac{\eta w}{2}.
\]
We now show that the other stopping nodes cannot receive this much weight. The sum of $\kappa$
over nodes not assigned to copies is at most $E_{\rm exc}$, so
\cref{res:incoming} bounds their arriving
weight by $2sE_{\rm exc}\le\eta w/16$. By \cref{res:foreign}, other assigned nodes receive
at most
\[
\mu w\sum_i(2q_2x+4\varepsilon_i)
\le\eta w\left(\frac1{64}+\frac1{510}\right)<\frac{\eta w}{32},
\]
using $4\mu wq_2x\le\eta/32$ and the bound on losses within $\Omega$ above.
Since only Alice has spoken, every surviving row reaches one stopping node.
The total arriving weight is therefore less than $3\eta w/32$, contradicting
the outgoing bound and proving the claim.

For the deletion bounds $\RP{W}{\theta}{C+\log w-1}$, we fix $j\in\{0,1\}$ and transpose a restriction of $W$
with row weight at least $\theta$ and column weight at least
$\beta_j=(3/5)2^{-j}$. In $U$ its tuple event $Y$ has weight at least
$\theta$. We discard the $\mathcal D$-rows, leaving a restriction of
the interlace of $N_{\rm ref}=N[\Omega,\mathcal K]$. The average retained
weight on $\Omega$, normalized within $\Omega$, is at least
$\beta'_j=(\beta_j-(1-\lambda))/\lambda$.
We put $\kappa=w2^{-j}$ and apply \cref{det:conditioning} to $Y$ with
error bound one. If $m$ copies have coordinates outside the exceptional set and
retain normalized weight on $\Omega$ at least $\kappa x$,
block averaging gives $w\beta'_j\le m+e(\theta)+w\kappa x$ and hence
\[
 \frac m\kappa\ge2^j\beta'_j-wx-\frac{4e(\theta)}w
 \ge\frac{(3/5)-4(1-\lambda)}\lambda-\frac1{256}-\frac1{32}
 =\frac{44}{75}-\frac1{256}-\frac1{32}>\frac12.
\]
Thus at least $\kappa/2+1$ blocks indexed outside the exceptional set remain. We apply
\cref{det:rareboost} to $N_{\rm ref}$ with $a=\log\kappa$, using all
of $Y$, of conditional weight one. The parameter choices give $\kappa\ge2$ and
$\kappa x\le1$. The lower bound is
$C-1+\log\kappa=C+\log w-1-j$, proving both bounds after transposition.
Applied to the whole matrix, the restriction bound proved first gives $\D(W)\ge C+\log w$.
Conversely, the player holding the copy index sends it in $\log w$ bits, and
the players then run the cost-$C$ protocol for $N$, giving equality.
\end{proof}

\subsection{Restoring the bound for the next round}

To repeat the construction, we need the deletion bounds at the exact cost. \Cref{cl:bridge}
recovers the missing bit from the bound of \cref{res:alignment} for column
weight at least $1-\delta_U$. Its proof follows the many-block estimate in
\cref{sec:interlace-induction} (p.~\pageref{proof:cl:bridge}).

\begin{theorem}[Restoring the deletion bounds]\label{cl:bridge}
Let $W$ have $m\ge1$ rows and $n\ge1$ columns, and suppose its exact cost
satisfies $\D(W)\ge3$. Let $\theta>0$ and $0<\delta\le1/64$. Assume that every
restriction retaining row weight at least $\theta$ and column weight at least
$1-\delta$ has cost at least $\D(W)$, and that $\RP{W}{\theta}{\D(W)-1}$ holds.
For an integer $b\ge1$, choose any power of two $Q=2^r$ satisfying
$Q\delta\ge80b$, and assume $Q\theta\le1/20$. Then $R^+=(\Int{Q}{W})^{\mathsf T}$
satisfies
\[
\D(R^+)=\D(W)+r,\qquad \RP{R^+}{2^{-b}}{\D(W)+r}.
\]
Independently of the lower-bound hypotheses, any upper bound on $\D(W)$ gives an
upper bound larger by $r$ on $\D(R^+)$. The matrix $R^+$ has exactly $n^Q$ rows
and $Qm$ columns.
\end{theorem}

The final sampling step of \cref{sec:sampling} produces a small table only
if the rank of $R^+$ grows slowly. The next lemma bounds this rank in both the
full-product and the explicit case.

\begin{lemma}\label{lem:amplifier-rank}
In the construction encoding the graph, let
$h=2\log s+\log q_2+\log w+\log Q$ and suppose
$|\mathcal I|\le2sq_2$. Then $\rank(R^+)\le2^h\rank(R)$.
The same bound holds if any tuple set is replaced by an indexed list
of tuples from that set, including products of such lists.
\end{lemma}
\begin{proof}
A width-$t$ interlace stacks $t$ blocks obtained by selecting or repeating
columns of its base matrix, and hence multiplies rank by at most $t$.
The template rows form nested interlaces of widths $q_1$, $q_2$ and $s$,
so the template part has rank at most $sq_1q_2\rank(R)$.
For each item $i$, its block has entries $R(u,\psi_i(Z))$, where the selected
column $\psi_i(Z)$ is independent of $u$, so its rank is at most $\rank(R)$.
After the outer interlaces, rank subadditivity gives
\[
\rank(R^+)\le wQ(sq_1q_2+|\mathcal I|)\rank(R)
\le wQsq_2(q_1+2)\rank(R)=2^h\rank(R),
\]
using $q_1+2=s$. Restricting tuple lists, duplicating indexed labels, and
transposing preserve this argument.
\end{proof}

\begin{proof}[Proof of \cref{thm:amplifier}]
\phantomsection\label{proof:thm:amplifier}
As in \cref{src:profiles}, we suppress the subscript $\rm g$ in the explicit
case. In both cases $h=2\log s+\log q_2+\log w+\log Q$. Beyond its effective
cutoff, \cref{lem:full-schedule} gives $h=L+O_{\Delta,\gamma}(\log L)$ in the
full-product case, and \cref{lem:generated-schedule} gives
$h=40\log s=10\log\mathsf N$ in the explicit case.

We first prove (ii), so let the graph be NO, $\D(R)=d\ge3$ and
$\RP{R}{\rho}{d}$. By \cref{src:profiles}, the per-copy deletion
bounds and the collision bound hold. The construction has exact marginals (in
the explicit case, those of $S_{\rm g}$ and $\mathcal J_{\rm g}$ in
\cref{det:construction}), and \cref{lem:common-schedule} gives the numerical bounds
\eqref{eq:width-margins}. Hence these bounds meet the
hypotheses of \cref{src:near-full}, so $\D(N)=C$ even after we delete at most
$\eta$ of each input set's weight. This gives (H1) and (H2) of
\cref{sec:internal-interlace}, and (H4) is \cref{lem:common-schedule}. We
verified (H3) for full products right after stating the hypotheses. In the
explicit case, \cref{det:construction} provides the three independent
sequences. \Cref{det:flat} shows that decoding the nested template labels gives
the exact $8192$-wise independence of (H3), and that fixing the other two
sequences leaves it unchanged.
Thus \cref{res:alignment} gives $\D(W)=C+\log w$ and
$\RP{W}{\theta}{C+\log w-1}$, and cost $\ge C+\log w$ on every restriction
with row weight $\ge\theta$ and column weight $\ge1-\delta_U$.
Only the restoration step differs between the cases.
\begin{itemize}
\item \emph{Full-product case.} By \cref{lem:full-schedule},
\(0<\delta_U\le1/64\), \(Q\delta_U\ge80L\), and \(Q\theta=1/512<1/20\), so
\cref{cl:bridge} applies with \(b=L\) and \(\delta=\delta_U\), at threshold
$2^{-L}=\rho$.
\item \emph{Explicit case.} \Cref{det:restore}, which restores the
deletion bounds at threshold $\phi$, applies with
$(\phi,\delta,Q)=(\rho,\delta_U,Q)$. \Cref{lem:generated-schedule} verifies its
numerical hypotheses, and \cref{det:construction} gives the outer list the
prescribed exact marginals.
\end{itemize}
In both cases we restore the deletion bounds at threshold $\rho$ and cost
\[
C+\log w+\log Q
=d+\underbrace{2\log s+\log q_2+\log w+\log Q}_{h}+1,
\]
so $\D(R^+)=d+h+1$ and $\RP{R^+}{\rho}{d+h+1}$.

For (i), on either graph, the players send the copy index at each of the two
interlaces, in $\log w+\log Q$ bits, and then run the protocol of
\cref{src:upper}. This gives \(\D(R^+)\le\D(R)+h+1\). On a YES graph we omit
the final $+1$ of \cref{src:upper}, giving \(\D(R^+)\le\D(R)+h\). In the explicit
case every input decodes locally to the corresponding full-product input with
the same entry, and the protocol acts only on decoded inputs, so it applies
unchanged.

For (iii), the bound $q_2\ge4s|E|$ from \eqref{eq:width-margins} and
$|E|\ge n/2$ imply $|\mathcal I|=4s(n+s-3)\le2sq_2$, so
\cref{lem:amplifier-rank} gives the rank bound in both cases.
\end{proof}

%% ---- section: sampling ----

\section{Interlacing on small input lists}
\label{sec:small-lists}
\subsection{Constructing exact marginals}

The deterministic reduction replaces each full product by an explicit tuple
list (\cref{def:exact-tuples}) whose induced weights on every set of at most $8192$
coordinates equal the product weights.

The construction adds one coordinate at a time. To each tuple we attach the
zero--one vector indicating its values on all sets of at most $8192$ coordinates.
After each extension, the construction moves weight along linear dependences
among these vectors. Each move removes at least one tuple, keeps all weights
nonnegative and leaves the joint weights on these coordinate sets unchanged.
The moves continue until no more tuples remain than there are constraints. This is the sequential
construction idea of \citet[Section~3]{koller1994sample}. The proof of the
next lemma, including the implementation and bit bound, is in
\cref{sec:support-algorithm}. We store the rational weights
explicitly and use them to measure retained input sets.

\begin{lemma}[Small rational tuple lists]\label{det:code}
Fix an integer independence order $\mathsf D\ge1$. Given normalized rational weights $\mu$ on an alphabet of
size $m$ and an integer $q\ge1$, a deterministic algorithm constructs an exactly
$\mathsf D$-wise independent tuple list $\mathsf{Code}_{\mathsf D}(q,\mu)$ (\cref{def:exact-tuples}) with
a label set of size at most $(1+qm)^{\mathsf D}$.
For fixed $\mathsf D$, the construction time, total output length and
rational-weight bit lengths are polynomial in $q$, $m$ and the binary encoding
length of $\mu$.
\end{lemma}

Below we always use $\mathsf{Code}_{8192}$.
Keeping the weights as binary-encoded rational numbers avoids expanding them
into repeated labels: a denominator with polynomially many bits can represent
far too many copies to print in polynomial time.

\subsection{Restoring the deletion bounds}

The matrix $W$ of \cref{res:alignment} satisfies the deletion bounds one bit below
its cost, and has full cost on restrictions retaining almost all column weight. One more interlace
recovers the missing bit. The proof distinguishes whether Bob has spoken. Before
his first message the bound for almost all columns applies. Afterward,
\cref{det:suffix} gives the missing bit.

\begin{theorem}\label{det:restore}
Let $0<\phi\le1$, $0<\delta\le1/64$ and $\theta>0$. Suppose $W$, with input weights
$\mu_X,\mu_Y$, satisfies $\D(W)=C\ge3$, $\RP{W}{\theta}{C-1}$, and
\[
\D(W[A,B])\ge C \quad\text{whenever }\mu_X(A)\ge\theta
\text{ and }\mu_Y(B)\ge1-\delta.
\]
Let $Q$ be a power of two satisfying
\[
Q\ge128\bigl(e(\phi,\delta/2)+1\bigr), \qquad Q\theta\le1/20,
\]
where $e(\cdot,\cdot)$, defined before \cref{det:conditioning}, bounds the
number of exceptional coordinates in that lemma.
Interlace $W$ $Q$ times on any exactly $8192$-wise independent tuple list with
the base-column weights, and transpose the result to obtain $R^+$. Then
\[
\D(R^+)=C+\log Q, \qquad \RP{R^+}{\phi}{C+\log Q}.
\]
Independently of the lower-bound hypotheses, a protocol of depth $U$ for $W$
yields one of depth $U+\log Q$ for $R^+$.
\end{theorem}
\begin{proof}
To prove the lower bounds, we fix $j\in\{0,1\}$ and a restriction of $R^+$ retaining row weight
$\ge\phi$ and column weight $\ge(3/5)2^{-j}$. After undoing the transpose, the retained
columns form a label event $Y$ of weight at least $\phi$, and the within-block
weights of the retained rows average at least $(3/5)2^{-j}$ over the $Q$ blocks. We apply
\cref{det:conditioning} to $Y$ with error bound $\delta/2$ and put
\[
e=e(\phi,\delta/2),\qquad k=Q2^{-j},\qquad a=\log k.
\]
Since $Q$ is a power of two with $Q\ge128(e+1)>2$, $a$ is a positive integer. More than $k/2$ blocks
have coordinates outside the exceptional set and retained base-row weight
at least $k\theta$: otherwise the
block-averaging observation gives
$(3/5)Q2^{-j}\le k/2+e+Qk\theta$, and hence
\[
\frac35\le\frac12+\frac{2^je}{Q}+Q\theta
\le\frac12+\frac1{64}+\frac1{20}<\frac35.
\]
We retain these blocks and all of $Y$.

Suppose the restriction to these blocks and $Y$ has a protocol of depth at most $C+a-1$,
and let $\nu$ denote the tuple-label weights. We follow the protocol for $a$
messages, starting from conditional column weight one, by the weight-halving
walk from the proof of \cref{det:rareboost}, which there runs for $a-1$ messages. At each Alice
message, we assign every current block to a child retaining at least half its row
weight and keep a child assigned at least half the blocks. At each Bob message, we
keep a child retaining at least half the current column weight under $\nu\mid Y$.
If the protocol ends earlier, we pad with dummy Alice messages. If $\alpha$ messages are Alice's and $\beta$
are Bob's, then $\alpha+\beta=a$, and the surviving restriction has a column set $K'\subseteq Y$ and
\[
\begin{gathered}
\text{more than }2^{\beta-1}\text{ blocks outside the exceptional set},\qquad
\text{row weight per block at least }2^\beta\theta,\\
(\nu\mid Y)(K')\ge2^{-\beta},\qquad
\text{at most }C-1\text{ bits remaining}.
\end{gathered}
\]
If $\beta=0$, the columns are still all of $Y$. At least one block indexed
outside the exceptional set remains, and its row weight is at least $\theta$. Its projected column set has base
weight at least $1-\varepsilon(\phi,\delta/2)\ge1-\delta$, by
\cref{det:conditioning}. Choosing retained representatives
(\cref{obs:representatives}), we obtain a restriction of $W$ that requires $C$ bits by the hypothesis for column weight $\ge1-\delta$.

If $\beta\ge1$, we apply the short-continuation bound \cref{det:suffix} with
$d=C-1$, $b=\beta-1$, and $x=\theta$. Its hypotheses hold. The number of blocks is an integer exceeding
$2^{\beta-1}$, so it is at least $2^b+1$. Each block has row
weight at least $2^{b+1}x$, and $(\nu\mid Y)(K')\ge2^{-b-1}$. Also $d\ge2$ and
$2^{b+1}x=2^\beta\theta\le k\theta\le1$. The lemma therefore requires $d+1=C$ bits.
In both cases this exceeds the $C-1$ remaining bits.

Thus each restriction requires at least $C+a=C+\log Q-j$ bits. Transposing
back proves both deletion bounds and the lower bound on the whole matrix. For the
upper bound, send the block index in $\log Q$ bits, decode that coordinate,
and run the depth-$U$ protocol for $W$. Exchanging the players gives the
same cost for $R^+$. Taking $U=C$ proves equality.
\end{proof}

\section{Deterministic gap tables}\label{sec:deterministic}
\label{sec:deterministic_amplification}

The shared amplification proof applies to explicit lists once their
marginal weights and independence properties are verified. We now construct
these lists and bound their encoding length, making each round a deterministic
polynomial-time operation. \Cref{det:reduction} gives the resulting gap and
encoding bound after \(k\) rounds.

The subscript $\rm g$ marks parameters and matrices of the explicit-list
construction. We choose a
power of two $s_{\rm g}$ above the fixed cutoff in
\cref{lem:generated-schedule}, with $\mathsf N:=s_{\rm g}^4\ge|V(G)|$ for the
source graph $G$.
Set $L_{\rm g}=\log\mathsf N$ and $\rho_{\rm g}=1/\mathsf N$. The choices in
\cref{def:amplification-parameters} then fix every other parameter as a
function of $s_{\rm g}$. In particular, since $\rho_{\rm g}=s_{\rm g}^{-4}$, $q_{2,\rm g}=s_{\rm g}^2/\rho_{\rm g}=s_{\rm g}^6$,
$w_{\rm g}=s_{\rm g}^8$, $Q_{\rm g}=s_{\rm g}^{24}$, and the block threshold is
$x_{\rm g}=\rho_{\rm g}^8$. These values are the same in every round, although the base
matrix $R$ changes from round to round.

\begin{corollary}[Explicit case]\label{det:amplifier}
For the construction \(R\mapsto R_{\rm g}^{+}\) of \cref{det:construction} with the
parameters of \cref{lem:generated-schedule}, parts (i) and (ii) of
\cref{thm:amplifier} hold with \(L=L_{\rm g}\), so that \(2^{-L}=\rho_{\rm g}\), and with
\[
h=h_{\rm g}:=2\log s_{\rm g}+\log q_{2,\rm g}+\log w_{\rm g}+\log Q_{\rm g}
=40\log s_{\rm g}=10\log\mathsf N.
\]
The construction is deterministic and polynomial in \(\mathsf N\) and the
explicit rational encoding length of \(R\). Every generated input decodes to an
input of the corresponding full-product construction with the same entries.
\end{corollary}

After defining the lists below, we prove this corollary in
\cref{sec:deterministic-endpoints}.

\subsection{Replacing products by explicit lists}

A coordinate of a tuple in one small list can itself be a label of another
small list. The next lemma shows that decoding such nested labels preserves
the product weights on every small set of coordinates. It also records why
fixing independently weighted labels of other lists does not change these weights.
As in \cref{sec:models}, inputs with the same decoded value keep separate
labels and weights.

\begin{lemma}\label{det:flat}
Let an indexed list on \(Z^{q_{\rm in}}\) be exactly \(8192\)-wise independent
with marginal \(\mu_Z\). Choose labels of this inner list at \(q_{\rm out}\) positions using
an exactly \(8192\)-wise independent outer list whose marginal equals the
inner label weights.
Then the decoded coordinates, indexed by outer--inner pairs, are exactly
\(8192\)-wise independent with marginal \(\mu_Z\). If the outer-list label
has product weights with the tuple of labels from finitely many other lists,
fixing those other labels at any joint value of positive weight leaves its
decoded weights unchanged.
\end{lemma}

\begin{proof}
Given at most $8192$ prescribed decoded coordinates, we group them by
outer position. Outer independence factors their joint weight over
these groups, and inner independence factors each group into its prescribed
coordinate marginals. Each factorization involves at most $8192$ coordinates,
as both independence hypotheses require.
Fixing the other labels leaves the outer-label weights unchanged, and hence
also leaves their decoded weights unchanged.
\end{proof}

\begin{definition}\label{det:construction}
Let \(R:X\times Y\to\{0,1\}\) be a rationally weighted base matrix and use the
parameter choices of \cref{lem:generated-schedule} and construct the items
and their values at each coordinate as in \cref{def:source-matrix}.
Replace the product alphabets \(\Gamma\), \(S\), and \(J_2^s\) by the indexed
lists \(\Gamma_{\rm g}\), \(S_{\rm g}\), and \(\mathcal J_{\rm g}\) of
\cref{det:code}, defined by
\[
\Gamma_{\rm g}=\mathsf{Code}_{8192}(2s_{\rm g}-1,\mu_Y),\qquad J_{\rm
g}=[q_{1,\rm g}]\times X,\qquad S_{\rm g}=\mathsf{Code}_{8192}(q_{2,\rm
g},\mu_{J_{\rm g}}),
\]
\[
J_{2,\rm g}=[q_{2,\rm g}]\times\Gamma_{\rm g},\qquad \mathcal J_{\rm
g}=\mathsf{Code}_{8192}(s_{\rm g},\mu_{J_{2,\rm g}}).
\]
Integer factors are uniform, and all other factors have their base or indexed
weights. On decoded labels use
\[
A_{\rm g}(\gamma,(r,u))=R(u,\gamma_r),\qquad R_{\rm temp,g}(z,(\alpha,\gamma))=A_{\rm g}(\gamma,z_\alpha).
\]
Use the raw entry rules, family tags, and type weights of
\cref{def:source-matrix}. For the family corresponding to $\Omega$, the
column weights are the product of the weights of a $\mathcal J_{\rm g}$ label
and the weights induced on the tail by $\Gamma_{\rm g}$. For the family
corresponding to $\mathcal D$, the weights are induced from $J_{2,\rm g}$ under
the map defining $\mathcal D$: the tail and all repeated labels use the same
$\gamma$. This defines $B_{\rm raw,g}$, its transpose $N_{\rm g}$, and its
normalized restriction $N_{\rm ref,g}$ to inputs with the $\Omega$ family tag.

The lower-bound proof fixes the type/bin/item sequence and the base-label
sequence before restricting the template labels; the latter must still have
their original weights, as required by (H3).
To interlace \(w_{\rm g}\) copies of \(N_{\rm g}\), we therefore use three
tuple lists with product weights between their labels. Each list consists
of sequences of length \(w_{\rm g}\) and is exactly \(8192\)-wise independent
under its label weights.
Type/bin/item labels have the weights induced by those labels in the raw rows,
base labels have weights \(\mu_X\), and template labels have the \(S_{\rm g}\)
weights. We enumerate the Cartesian product of these three label sets.
At each position, the type label in the first sequence determines whether the
decoded row is a template row \((p,z)\), from the bin and template labels, or a row
\((i,u)\) indexed by an item, from the item and base labels. Since the three sequences
are independent and each is exactly \(8192\)-wise independent, their
coordinatewise combination gives an exactly \(8192\)-wise
independent tuple list with the raw-row marginal. Interlacing \(N_{\rm g}\) on
this list defines \(U_{\rm g}\). Set \(W_{\rm g}=U_{\rm
g}^{\mathsf T}\). We interlace \(W_{\rm g}\) at width \(Q_{\rm g}\) using
\(\mathsf{Code}_{8192}\) with its column weights and then transpose to obtain \(R_{\rm
g}^{+}\). We store all input labels and rational weights explicitly, keeping inputs with
equal decoded values separate.
\end{definition}

As in \cref{obs:representatives}, we obtain projected restrictions from
chosen representatives. The per-copy deletion bounds and the collision
bound are the explicit case of \cref{src:profiles}. Their proofs use
\cref{det:conditioning,det:rareboost} to handle retained events of weight
as small as $x_{\rm g}$. With these bounds, the collision and stopping-node
arguments of \cref{src:near-full,sec:internal-interlace} apply to the
indexed construction.

\subsection{Iteration, encoding size and deterministic running time}
\label{sec:deterministic-endpoints}

To prove \cref{det:amplifier}, we need only add the encoding bound to the
explicit case of \cref{thm:amplifier}.

\begin{proof}[Proof of \cref{det:amplifier}]
\phantomsection\label{proof:det:amplifier}
We suppress the subscript $\rm g$. Each width is a fixed polynomial
in $\mathsf N$. Therefore the fixed number of calls to \cref{det:code}, all of order
$8192$, and enumeration of the Cartesian products of their label sets take time bounded by a fixed polynomial
in $\mathsf N$ and the encoding length of $R$. The indexed labels and rational
weights needed by the next round have encoding length within the same bound.

Parts (i) and (ii) follow from the proof of \cref{thm:amplifier}, with
\cref{det:restore} used in place of \cref{cl:bridge}. We apply \cref{det:restore}
with $(\phi,\delta,Q)=(\rho,\delta_U,Q)$, and \cref{lem:generated-schedule}
verifies its numerical hypotheses. The upper-bound
protocol uses only the decoded inputs and not their weights, so it also works
for the indexed construction.
\end{proof}


\begin{theorem}\label{det:reduction}
There are effective absolute constants \(\mathsf A\ge2\) and \(C_0>0\) and a
deterministic algorithm with the following guarantee. Given a YES or NO
source graph \(G\) from \cref{src:gap-colouring}
and an integer \(k\ge1\), it produces an explicit total Boolean matrix
\(F_{G,k}\) and a public integer \(B_k^{\rm g}\) such that
\[
\text{YES: }\D(F_{G,k})\le B_k^{\rm g},\qquad \text{NO: }\D(F_{G,k})=B_k^{\rm
g}+k.
\]
For every \(\mathsf N=s_{\rm g}^4\ge|V(G)|\) meeting the parameter conditions
of \cref{lem:generated-schedule}, its running time and full
binary output length are at most \(\mathsf N^{C_0\mathsf A^k}\). The constants
depend only on the fixed source reduction and construction. The output may be
square with power-of-two side length at least \(\mathsf N\). The public threshold is
\[
B_k^{\rm g}=d_{0,\rm g}+kh_{\rm g},
\]
where \(d_{0,\rm g}\) is the cost of the initial matrix in \cref{cl:seed}
with \(L=L_{\rm g}\):
\[
d_{0,\rm g}=\log\operatorname{pow}_2(64(L_{\rm g}+1))+1,\qquad L_{\rm
g}=\log\mathsf N.
\]
\end{theorem}

\begin{proof}
We choose the parameters of \cref{lem:generated-schedule} and let $R_0$ be the
initial matrix of \cref{cl:seed} with $L=L_{\rm g}=\log\mathsf N$. We apply
\cref{det:construction} $k$ times, using each output as the next base matrix
and storing each indexed table with its rational
input weights. Finally, we forget the weights
and duplicate inputs to make a square table with power-of-two side at least $\mathsf N$.

By \cref{cl:seed}, the initial matrix has exact cost $d_{0,\rm g}\ge3$ and
satisfies the deletion bounds at threshold $\rho_{\rm g}=2^{-L_{\rm g}}$.
\Cref{prop:iteration}, in the explicit case (\cref{det:amplifier}), therefore gives
\[
\D(R_k)\le d_{0,\rm g}+kh_{\rm g}\quad\text{on YES},\qquad
\D(R_k)=d_{0,\rm g}+k(h_{\rm g}+1)\quad\text{on NO}
\]
for the matrix $R_k$ after $k$ rounds. Forgetting weights and duplicating rows
and columns to make the table square change neither cost. This proves the gap
with $B_k^{\rm g}=d_{0,\rm g}+kh_{\rm g}$.

For size and time, let $Z_i$ be $\mathsf N+2$ plus the complete encoding
length after round $i$, including labels and normalized rational weights. The polynomial
bound for the initial matrix and \cref{det:amplifier} give effective constants $C_1$ and a
power of two $\mathsf A\ge2$, independent of $i$ and $k$, such that
\[
Z_0\le\mathsf N^{C_1},\qquad
Z_{i+1},\ \text{time for round }i+1\le Z_i^{\mathsf A},\qquad
Z_i\le\mathsf N^{C_1\mathsf A^i}.
\]
After enlarging the effective constant $C_0$ if necessary, we bound the sum of
the $k$ round times and the final square padding by $\mathsf N^{C_0\mathsf A^k}$.
Computing the public threshold and the parameters takes time within the same bound.
\end{proof}

We now derive \cref{intro:deterministic} by taking a fixed number $k$ of rounds,
and \cref{det:time} by taking $k=\Theta(\log\log N)$ rounds for a sparse formula on
$N$ variables.

\begin{proof}[Proof of \cref{intro:deterministic}]
\phantomsection\label{proof:intro:deterministic}
For fixed $k$, we apply the graph reduction of \cref{src:gap-colouring} and then
\cref{det:reduction}. The least value $\mathsf N=s_{\rm g}^4$ meeting the parameter conditions obeys
$\mathsf N\le16\max\{|V(G)|,\mathsf N_0\}$ for an effective constant
$\mathsf N_0$, by \cref{lem:generated-schedule}. Since the graph reduction runs in
polynomial time and $k$ is fixed,
$\mathsf N^{C_0\mathsf A^k}$ is polynomial in the formula length.
The resulting public threshold separates cost $\le B_k^{\rm g}$ from
cost exactly $B_k^{\rm g}+k$, on a square table with power-of-two side.
\end{proof}

\begin{proof}[Proof of \cref{det:time}]
\phantomsection\label{proof:det:time}
We choose the number $k$ of rounds so that the gap $k$ exceeds twice the
allowed error $c_*\log\log T$. At the same time, $k$ must be small enough that the
construction and the estimator run in quasipolynomial time in the size of the sparse formula.
Let $\mathsf A$ be the effective power of two from \cref{det:reduction}, and
set $c_*=1/(8\log\mathsf A)$. Suppose such an estimator exists.
For a sparse formula on $N$ variables, \cref{src:gap-colouring} and
\cref{lem:generated-schedule} give a value
$\mathsf N=s_{\rm g}^4\ge\max\{N,|V(G)|\}$ satisfying the parameter conditions.
Since $|V(G)|\le N(\log(N+2))^p$ for the fixed $p$ of
\cref{src:gap-colouring}, we have $\mathsf N=N^{O(1)}$.
For sufficiently large $N$, we choose
\[
k=\floor{\frac{\log\log N}{\log\mathsf A}}.
\]
Then $k=\Theta(\log\log N)$ and $\mathsf A^k\le\log N$. The square output of
\cref{det:reduction} has threshold $B_k^{\rm g}$ and $T$ entries, where
\[
N^2\le T\le\mathsf N^{C_0\mathsf A^k},\qquad
\log\log T\le k\log\mathsf A+O(\log\log N).
\]
For sufficiently large $N$, the second bound is at most $3k\log\mathsf A$:
indeed it is at most $2\log\log N+O(1)$, whereas
$k\log\mathsf A=\log\log N+O(1)$.
Consequently
\[
2c_*\log\log T\le 3k/4<k.
\]
Comparing the estimate with the public midpoint
$B_k^{\rm g}+k/2$ therefore decides satisfiability.
Both construction and the estimator's fixed polynomial running time in $T$
are bounded by
\[
2^{O((\log N)^2)}.
\]
The estimator's guarantee holds for all sufficiently large tables, so it applies here because $T\ge N^2$.
For the finitely many remaining values of $N$, including those with $k<1$,
exhaustive search over assignments decides satisfiability.
This proves the claimed deterministic running-time implication.
\end{proof}

%% ---- section: alternation ----

\section{Final sampling and growing error}
\label{sec:quantitative}

In the full-product case of \cref{prop:iteration}, the iterated matrix $R_k$
has the desired gap but is too large to print, so we replace it by a sampled
submatrix. A single sample must be hard for every protocol of depth below the NO-case
bound. A small sample suffices because the rank bound $\rank(R_k)\le r_k$
limits the number of column sets on which such
protocols can be correct (\cref{thm:rank-count}).

\subsection{Counting column sets on which short protocols are correct}

For the union bound in \cref{sec:sampling}, we count the sets of columns on
which a short protocol can be correct on a retained row set. A single count
covers all row sets at once. We fix the tree, Alice's messages and the leaf
outputs, and leave Bob's messages free to depend on the column.

\begin{definition}[Columns on which Bob can complete the protocol]\label{def:compatibility}
Let a depth-$t$ tree $P$ specify a complete binary tree of
depth $t$, the speaker at each internal node, all Alice maps (from her row to
her bit at each of her nodes), and all leaf outputs. Bob's maps are unspecified. For a Boolean matrix
$F:X\times Y\to\{0,1\}$ and $U\subseteq X$, let $G_P(U)$ be the set of columns
for which some assignment of bits to Bob's nodes makes the tree correct on every
row in $U$. The assignment may depend on the column, but not on Alice's row.
\end{definition}

The following characterization is due to \citet[Lemma~3.2]{mackenzie2026condensation}.

\begin{observation}[Short protocols as containments]\label{obs:compatible-protocol}
For retained sets $U\subseteq X$, $V\subseteq Y$ and an integer $t\ge0$, with correctness on an empty restriction understood
vacuously,
\begin{equation}\label{eq:compatibility}
 \D(F[U,V])\le t
 \quad\Longleftrightarrow\quad
 V\subseteq G_P(U)\text{ for some depth-}t\text{ protocol tree }P\text{ of this kind}.
\end{equation}
\end{observation}

\Cref{fig:compatible-trees} uses two one-bit trees to show why leaf
requirements combine by intersection at an Alice node and by union at a Bob
node. We label each leaf by the columns on which its output is correct for
the rows reaching it.

\begin{figure}[H]
\centering
{\small
The same matrix in both trees: $F(x,y)=x_y$.

\smallskip
\begin{tabular}{c@{\qquad}*{3}{>{\centering\arraybackslash}p{8mm}}}
\toprule
&\multicolumn{3}{c}{Bob's column $y$}\\
\cmidrule(l){2-4}
Alice's row $x$&$1$&$2$&$3$\\
\midrule
$001$&$0$&$0$&$1$\\
$011$&$0$&$1$&$1$\\
\bottomrule
\end{tabular}
\par}
\medskip
\captionsetup[subfigure]{justification=centering}
\begin{subfigure}[t]{.48\linewidth}
\centering
\begin{tikzpicture}[font=\small,>=Stealth,
 box/.style={draw,rounded corners=2pt,align=center,inner sep=4pt},
 leaf/.style={box,text width=2.55cm}]
\node[box,draw=aliceblue] (a) at (0,0) {Alice sends $x_2$};
\node[leaf,anchor=north] (a0) at (-1.65,-1.55) {row $001$\\output $0$\\columns $\{1,2\}$};
\node[leaf,anchor=north] (a1) at (1.65,-1.55) {row $011$\\output $1$\\columns $\{2,3\}$};
\coordinate (am) at (0,-.95);
\coordinate (al) at (-1.65,-.95);
\coordinate (ar) at (1.65,-.95);
\draw[aliceblue] (a.south) -- (am);
\draw[aliceblue] (am) -- node[midway,above=3pt,inner sep=1pt] {$0$} (al);
\draw[aliceblue] (am) -- node[midway,above=3pt,inner sep=1pt] {$1$} (ar);
\draw[->,aliceblue] (al) -- (a0.north);
\draw[->,aliceblue] (ar) -- (a1.north);
\end{tikzpicture}
\caption{\textbf{Both required: intersection.}\\
$\{1,2\}\cap\{2,3\}=\{2\}$.}
\end{subfigure}\hfill
\begin{subfigure}[t]{.48\linewidth}
\centering
\begin{tikzpicture}[font=\small,>=Stealth,
 box/.style={draw,rounded corners=2pt,align=center,inner sep=4pt},
 leaf/.style={box,text width=2.55cm}]
\node[box,draw=boborange] (b) at (0,0) {Bob chooses using $y$};
\node[leaf,anchor=north] (b0) at (-1.65,-1.55) {rows $001,011$\\output $0$\\columns $\{1\}$};
\node[leaf,anchor=north] (b1) at (1.65,-1.55) {rows $001,011$\\output $1$\\columns $\{3\}$};
\coordinate (bm) at (0,-.95);
\coordinate (bl) at (-1.65,-.95);
\coordinate (br) at (1.65,-.95);
\draw[boborange] (b.south) -- (bm);
\draw[boborange] (bm) -- node[midway,above=3pt,inner sep=1pt] {$0$} (bl);
\draw[boborange] (bm) -- node[midway,above=3pt,inner sep=1pt] {$1$} (br);
\draw[->,boborange] (bl) -- (b0.north);
\draw[->,boborange] (br) -- (b1.north);
\end{tikzpicture}
\caption{\textbf{Either works: union.}\\
$\{1\}\cup\{3\}=\{1,3\}$.}
\end{subfigure}
\caption{Alice's branch depends on the row, so a column on which the tree is correct must
satisfy both leaf conditions. Bob's branch depends on the column, so either
leaf may work, but his chosen leaf must answer both rows correctly.}
\label{fig:compatible-trees}
\end{figure}

The count below follows from \citet[Lemma~4.3]{mackenzie2026condensation},
using the row-sum argument of Oum, S{\ae}ther and
Vatshelle~\cite[Theorem~4.2]{oum2014faster}: the columns on which selected
Boolean rows all vanish are exactly the zero coordinates of the integer sum of
those rows. Counting these sums instead of the row subsets gives a bound that depends
only on the rank and not on the number of rows. As a result, the final sample can be small.

\begin{theorem}[Counting the sets of columns]\label{thm:rank-count}
Let $r\ge1$ and $t\ge0$ be integers. If a Boolean matrix $F$ has real rank at most
$r$, then, over all depth-at-most-$t$ trees $P$ specified as in
\cref{def:compatibility} and all retained row subsets $U$, there are at most
\begin{equation}\label{eq:rank-count}
 2^{(r+1)\ceil{\log(r+2)}\,2^{t+1}}
\end{equation}
distinct sets $G_P(U)$. Repeated row and column labels are allowed.
\end{theorem}

Mackenzie's bound for depth exactly $t$ is $2^{B(r)2^t-1}$ with
$B(r)=(r+1)\ceil{\log(r+1)}+2$. Since $B(r)\le2(r+1)\ceil{\log(r+2)}$ for
$r\ge1$, and padding turns a tree of depth $t'<t$ into a depth-$t$ tree with the
same sets $G_P(U)$, this gives \eqref{eq:rank-count}. Repeated labels are
handled by passing to distinct values~\cite[Section~2]{mackenzie2026condensation}.

\subsection{Sampling columns and then rows}
\label{sec:sampling}

We sample the same number $m$ of rows and columns, and we take $m$ to be a
power of two, so the output has the square power-of-two form of
\cref{thm:quantitative}. In the proof, we first show that, with high
probability, the sampled columns preserve the assumed lower bound $H$ on every
row set of original weight $\ge3/5$. We then analyze the independent row sample.

\begin{theorem}[Sampling rows and columns]\label{thm:compression}
Let $F:X\times Y\to\{0,1\}$ have real rank at most $r$, where $r\ge1$ is an
integer, and let $\mu_X,\mu_Y$ be input distributions on nonempty finite indexed sets.
Suppose an integer $H\ge1$ satisfies, for all $U\subseteq X$ and $V\subseteq Y$,
\[
\mu_X(U)\ge3/5,\quad\mu_Y(V)\ge3/5
\quad\Longrightarrow\quad \D(F[U,V])\ge H.
\]
Put
\[
m=\operatorname{pow}_2\bigl(16(r+1)\ceil{\log(r+2)}\,2^H\bigr),
\]
where $\operatorname{pow}_2(u)$ is the least power of two at least $u$.
Draw $m$ rows and $m$ columns independently, retaining separate indices for
repeated draws. The resulting $m\times m$ matrix has cost at least $H$ except
with probability less than $1/100$. Every protocol for $F$ also computes the
sampled matrix, so the sampled matrix has cost at most $\D(F)$ on every sample outcome.
\end{theorem}
\begin{proof}
Repeated labels can always be decoded locally. Conversely, a protocol on
repeated labels can be restricted to one occurrence of each label.

For every $U$ with $\mu_X(U)\ge3/5$ and every depth-$(H-1)$ tree
$P$ specified as in \cref{def:compatibility}, we have $\mu_Y(G_P(U))<3/5$: otherwise
\cref{obs:compatible-protocol} with $V=G_P(U)$ gives $\D(F[U,V])\le H-1$,
contradicting the restriction hypothesis. As $U$ and
$P$ vary, these sets form a family of size at most
$2^{m/16}$ by \cref{thm:rank-count}. By a union bound, the probability that all sampled
columns lie in one of them is at most $2^{m/16}(3/5)^m$.
Outside this event, the sampled-column matrix
$F_Y(x,j)=F(x,y_j)$ has cost $\ge H$ on every original row set of weight
$\ge3/5$.

We condition on such a column outcome. The transpose $F_Y^{\mathsf T}$ still has
rank at most $r$. Consider the sets $G_P(U)$ for this transpose, where $U=[m]$ is
its full row set and $P$ has depth $H-1$. Each has $\mu_X(G_P(U))<3/5$: otherwise
\cref{obs:compatible-protocol} and transposition give
$\D(F_Y[G_P(U),[m]])\le H-1$, contradicting the preceding guarantee. There are again at most $2^{m/16}$ such sets.
The row draws remain independent with distribution $\mu_X$, so the probability that
they all lie in one of these sets is at most $2^{m/16}(3/5)^m$. A protocol of depth at most
$H-1$ on the final matrix would give such a containment. Since $m\ge32$,
the total failure probability is at most
\[
2\cdot2^{m/16}(3/5)^m\le2^{1-7m/16}\le2^{-13}<1/100.
\]
The universal upper bound follows by decoding sampled indices and running a
protocol for $F$.
\end{proof}

We sample only once, after all amplification rounds on full products.
\Cref{thm:compression} bounds the failure probability but gives no test
certifying that a particular sample succeeds.

\subsection{Printing the sample and comparing errors}
\label{sec:reduction}
\label{sec:randomised_reduction}

We generate the sampled input labels and evaluate only the corresponding
entries of $R_k$.

We take a source graph on $n$ vertices above the fixed effective cutoff in
\cref{lem:full-schedule}, use the parameter choices of that lemma at every
round, and write
\[
L=\ceil{\log n},\qquad B_k=d_0+kh,\qquad H_k=B_k+k,\qquad r_k=2^{B_k-1}
\]
as in \cref{prop:iteration}, which gives $\rank(R_k)\le r_k$. The reduction
never computes a rank. A \emph{random tape} is the sequence of random bits used
by the algorithm. The next lemma bounds the running time on every tape, not
only on likely ones.

\begin{lemma}\label{lem:sample-access}
For the level-$k$ iteration from the initial matrix in \cref{prop:iteration}, use
$r_k,H_k$ in \cref{thm:compression}. A joint sampler and entry evaluator use $n^{O(k+1)}$
time and random bits on every tape, with an absolute implied constant
independent of $k$. Every output label is an actual input of $R_k$, and the
sample counts and tape length are fixed. The joint distribution of the generated labels is within total
variation less than $1/100$ of the ideal independent sample of \cref{thm:compression}.
\end{lemma}

Here total variation is the largest difference between the probabilities
that two distributions assign to an event. The proof in
\cref{sec:sampling-algorithm} gives the recursive construction with a fixed
number of random bits. Sampling is approximate,
but every printed entry is exact. Combining the lemma with
\cref{thm:compression}, we obtain the reduction.

\begin{theorem}[Randomized reduction]\label{thm:uniform-reduction}
One randomized algorithm, given a source graph on $n$ vertices above the fixed cutoff in \cref{lem:full-schedule} and
an integer $k\ge1$, outputs a total Boolean matrix $M$ and an integer $B$.
The table is square with power-of-two side length. On every tape the algorithm
uses $n^{O(k+1)}$ time and random bits, with an absolute implied
constant independent of $k$. For a YES source, $\D(M)\le B$ on every tape.
For a NO source,
\[
\Pr[\D(M)=B+k]\ge49/50.
\]
Here $B=B_k$ and the tape-independent entry count is
\[
m=\operatorname{pow}_2\bigl(16(r_k+1)\ceil{\log(r_k+2)}\,2^{H_k}\bigr),
\qquad T=m^2.
\]
\end{theorem}

\begin{proof}
The algorithm computes the parameters of \cref{lem:full-schedule} and the
initial matrix of \cref{cl:seed}, then $B_k,H_k,r_k,m$. It generates $m$ columns
and $m$ rows with \cref{lem:sample-access}, evaluates all $m^2$ row--column
pairs (keeping repeated labels as distinct indices), and outputs $(M,B_k)$.
The time and random-bit bounds follow from \cref{lem:sample-access}, and the
output is $m\times m$ with $m$ a power of two.
On a YES source, every generated label is an input of $R_k$. Decoding the
labels and running a protocol for $R_k$ therefore shows that every sample has cost at most $B_k$.

On a NO source, \cref{prop:iteration} gives
\[
\D(R_k)=H_k,\qquad \RP{R_k}{2^{-L}}{H_k},\qquad \rank(R_k)\le r_k.
\]
Because $2^{-L}\le1/2$, the $j=0$ lower bound applies whenever both retained
weights are at least $3/5$. Thus \cref{thm:compression} gives ideal sampling
failure less than $1/100$. By \cref{lem:sample-access}, using the implemented
sampler instead of the ideal independent sample adds less than $1/100$ to this
failure probability. Therefore, except with probability less than $1/50$, the
output has cost $\ge H_k$. Decoding the inputs gives cost $\le H_k$ on every
tape, so the cost is then exactly $B_k+k$.
\end{proof}

\subsection{From an approximation algorithm to a $3$-SAT algorithm}
\label{sec:time-tradeoff}

We now combine the two constructions with an estimator, proving the general
implication stated in \cref{thm:time-tradeoff}. Each construction produces
a known threshold $B$ and a gap $k$. Its size bound determines which
approximation errors are small enough to distinguish the two cases.

\begin{theorem}[From approximation algorithms to $3$-SAT algorithms]
\label{thm:time-tradeoff}
There is an effective absolute integer $\mathsf A\ge2$ and, for each fixed
clause-density bound $c_s$, an effective integer $C\ge1$ with the following
property. Suppose a deterministic or
bounded-error randomized polynomial-time estimator approximates $\D(M)$
with additive error at most a nondecreasing function $a(T)$ on every
sufficiently large explicit total Boolean table with $T$ entries.
Let $k=k(N)\ge1$ be an integer function computable in time polynomial in
$N+k(N)$. Choose either construction and its size bound:
\[
\begin{array}{c|c}
\text{Construction} & S(N,k)\\ \hline
\text{Randomized} & (N+2)^{C(k+1)}\\[2pt]
\text{Deterministic} & (N+2)^{C\mathsf A^k}
\end{array}
\]
If, for all sufficiently large $N$,
\[
2a\bigl(S(N,k(N))\bigr)<k(N),
\]
then $3$-SAT on $N$ variables and at most $c_sN$ clauses has an algorithm
running in time $S(N,k(N))^{O(1)}$. It is deterministic when both the
construction and the estimator are deterministic, and bounded-error
randomized otherwise. The exponent may depend on the estimator and the
algorithm computing $k$. It suffices for the estimator to work on square
tables whose side length is a power of two.
\end{theorem}

\begin{proof}[Proof of \cref{thm:time-tradeoff}]
\phantomsection\label{proof:thm:time-tradeoff}
Apply \cref{src:gap-colouring}, then \cref{thm:uniform-reduction} in the
randomized case or \cref{det:reduction} in the deterministic case. Since
the source graph and the deterministic size parameter are $N^{O(1)}$, a
fixed $C$ bounds both the table size $T$ and construction time by the
chosen $S(N,k)$. The constant $\mathsf A$ is the absolute constant from
\cref{det:reduction}. The tables tend to infinity with $N$.

Because $a$ is nondecreasing, $a(T)\le a(S(N,k))<k/2$. An accurate
estimate is therefore below the midpoint $B+k/2$ on a YES output and
above it on a successful NO output. If the estimator is randomized,
repeat this comparison independently a constant number of times and take
the majority answer, reducing its error below $1/100$. The randomized
construction adds failure probability at most $1/50$; the deterministic
construction adds none. Thus the combined algorithm has error less than
$1/3$ whenever randomness is used, and is exact otherwise.

Construction, computing $k$, estimation, and comparison with the binary
rational midpoint all take time $S(N,k)^{O(1)}$: in particular,
$S(N,k)\ge N+k$ absorbs the assumed polynomial cost of computing $k$.
Exhaustive search handles the finitely many smaller sizes.
\end{proof}

To read this in terms of a proposed $3$-SAT running time $R(N)=2^{g(N)}$,
suppose $g(N)/\log(N+2)$ grows without bound and the chosen number of rounds
can be computed within the stated time. A budget $2^{O(g(N))}$ allows
\[
\begin{aligned}
k&=\Theta\!\left(\frac{g(N)}{\log(N+2)}\right)
&&\text{in the randomized construction},\\
k&=\Theta\!\left(\log\frac{g(N)}{\log(N+2)}\right)
&&\text{in the deterministic construction}.
\end{aligned}
\]
The constant in the second choice is chosen so that
$\mathsf A^k\log(N+2)=O(g(N))$.
Accuracy better than half this gap, at the corresponding table sizes,
therefore yields a $3$-SAT algorithm within that budget. If such a $3$-SAT
algorithm is excluded by an explicit hypothesis, the same statement gives
inapproximability. To quantify a further improvement over a known algorithm,
choose a smaller target budget $R$. The theorem does not treat the best
known upper bound as optimal.

For \cref{thm:quantitative}, we use the exact sample size rather than
the coarser bound $S(N,k)$. This gives the coefficient $1/8$ in the
randomized error bound; the algorithmic implication is the same.

\begin{proof}[Proof of \cref{thm:quantitative}]
\phantomsection\label{proof:thm:quantitative}
Suppose a bounded-error randomized polynomial-time estimator achieves additive
error $a(T)=c\log T/\log\log T$ for some fixed $0\le c<1/8$ on every
sufficiently large square table of power-of-two side length. Fix
$8c<\alpha<1$. We obtain a bounded-error $2^{O(N^\alpha)}$ algorithm for
sparse $3$-SAT.

For a formula with $N$ variables, \cref{src:gap-colouring} produces a source
graph with $n$ vertices, where
\[
2N/\Delta\le n\le N(\log(N+2))^p
\]
for the fixed integer $p$ of that theorem. The lower bound follows from
$N\le|E|\le\Delta n/2$, also given there.
Thus the graph is above the fixed cutoff for every sufficiently large $N$, and
\[
L=\ceil{\log n}=\log N+O(\log\log N).
\]
Choose a fixed rational $\beta$ with $8c<\beta<\alpha$, and put
$k=\floor{N^\beta/\ceil{\log N}}$. This integer is computable in polynomial
time, using an integer root for the fixed rational power. By
\cref{lem:full-schedule,cl:seed}, $h=L+O(\log L)$ and $d_0=O(\log L)$, so
\[
B_k=d_0+kh=(1+o(1))k\log N.
\]
Substituting $r_k=2^{B_k-1}$ and $H_k=B_k+k$ into
\cref{thm:uniform-reduction} gives
$\log m=2B_k+k+O(\log(B_k+1))$. Hence the printed size satisfies
\begin{align}
\log T&=2\log m=4B_k+2k+O(\log(B_k+1))
             =(4+o(1))k\log N,\label{eq:quantitative-size}\\
\log\log T&=(\beta+o(1))\log N,
\qquad k=\left(\frac\beta4+o(1)\right)\frac{\log T}{\log\log T}.
\end{align}
Consequently $k>2a(T)$ for all sufficiently large $N$. The table already has
the square power-of-two form required by the assumed estimator.

For each estimator run, we test whether its output is below the rational
midpoint $B_k+k/2$. An accurate estimate answers this test correctly on every
YES output and every successful NO output. We repeat the run independently a
constant number of times and take the majority answer, which reduces the
decision error on each such table below $1/100$. This uses only comparison with a binary rational, which the
estimator model supports (\cref{def:table-estimation}). The reduction
fails with probability at most $1/50$, so the decision error is less than
$3/100<1/3$.

Both construction time and random-bit usage are
$n^{O(k+1)}=2^{O(N^\beta)}$ on every tape. The estimator also takes
$2^{O(N^\beta)}$ time because $\log T=O(N^\beta)$.
For all sufficiently large $N$, we have $k\ge1$, the graph is above the fixed cutoff, and $T$ meets the estimator's size threshold.
The finitely many smaller values of $N$ are decided by exhaustive search within
a fixed bound. Since $\beta<\alpha$, the total running time is
$2^{O(N^\alpha)}$, as claimed.
\end{proof}

\section{Strictly alternating protocols}\label{sec:alternation}
\label{sec:alternating}

We now prove the result that settles Yao's original question. If we simulate
the multi-bit messages of the unrestricted construction by inserting the missing
turns, the cost can double and the additive gap is lost. We therefore use a
separate matrix designed for strict alternation. The end of \cref{sec:overview} sketches
the argument. A balanced hypergraph coloring gives a short alternating protocol
(\cref{prop:alternating-upper}). In the NO case, a short protocol
induces a coloring of the vertices by Alice's early response rules. A monochromatic edge then gives many
array positions on which her remaining messages produce too few answer patterns
(\cref{yao:small-bias-escape}). With an explicit array list (\cref{yao:array-list}),
the table has polynomial size for each fixed gap, and no repeated amplification is
needed.

\subsection{Starter, output and matrix inputs}

Changing who starts or who must know the answer changes the cost by a
constant, so the gap transfers between the conventions of \cref{yao:model}.

\begin{lemma}\label{yao:comparison}
On nonempty input sets, each cost
$C\in\{D_A^{\rm alt},D_B^{\rm alt},D_{\rm free}^{\rm alt},Y_{\rm free}\}$
satisfies
\[
|C(M)-D_A^{\rm alt}(M)|\le2.
\]
\end{lemma}
\begin{proof}
Prepending one fixed bit sent by the new starter changes the starter, so
$|D_A^{\rm alt}-D_B^{\rm alt}|\le1$ and
\[
D_A^{\mathrm{alt}}-1\le D_{\mathrm{free}}^{\mathrm{alt}}\le D_A^{\mathrm{alt}}.
\]
Public output also gives receiver output. Conversely, at a nonzero-depth
terminal transcript the receiver is the next scheduled speaker and can announce
the answer in one bit. Public stopping tells both players when to make this
announcement. At depth zero, choose the designated output player as the starter
and let that player announce the answer. Thus
\[
Y_{\mathrm{free}}\le D_{\mathrm{free}}^{\mathrm{alt}}\le Y_{\mathrm{free}}+1,
\]
and the claimed two-bit bound follows. These transformations preserve public
stopping and use an input-independent starter.
\end{proof}

A hypergraph coloring is proper if no edge is monochromatic. In the YES case
the hypergraph has a balanced coloring with few colors. In the NO case it has no
proper coloring even with the larger number of colors that a short protocol would induce. We use the balanced
conversion of Guruswami and Lee~\cite{guruswami2018balanced}, in a form whose
edges have distinct vertices.

\begin{proposition}[Hypergraph coloring promise]\label{yao:source}
Fix an integer $t\ge4$, and set
\[
s_{\rm a}=2^t,\qquad K_{\rm a}=s_{\rm a}^2,\qquad S_{\rm a}=2^{s_{\rm a}}.
\]
There is a deterministic polynomial-time reduction from $3$-SAT to
$K_{\rm a}$-uniform hypergraphs with nonempty vertex sets and distinct vertices
in every edge such that:
\begin{enumerate}
\item a satisfiable formula has an $s_{\rm a}$-coloring with at least $s_{\rm a}-1$ vertices of every color in every edge;
\item an unsatisfiable formula has no proper $S_{\rm a}$-coloring.
\end{enumerate}
The polynomial may depend on $t$.
\end{proposition}

The proof in \cref{app:hypergraph-source} derives this promise from the cited
result and ensures that every edge has distinct vertices.

Alice holds a vertex and a tuple of addresses. Bob holds an array, a
coordinate selecting one of Alice's addresses, and an edge index. If the vertex
is outside the selected edge, or the edge index is unused, the answer is zero.

\begin{definition}\label{def:alternating-table}
Let $\mathcal H=(V,E)$ be a source hypergraph with an ordering of the vertices
in each edge. Set
\[
q_{\rm a}=2^{t-1}=s_{\rm a}/2, \qquad H_{\rm a}=2^b,
\]
where $H_{\rm a}$ is the least power of two at least $\max\{1,|E|\}$. Index the
edges by distinct $\alpha\in[H_{\rm a}]$, and let $\iota_\alpha(v)$ be the
position of $v$ in that edge, or zero if $v$ is absent or the index is unused.
Alice's address tuples are all of
\[
\mathcal C=[H_{\rm a}]^{q_{\rm a}},
\]
each used once.
For any labeled Boolean array list $\mathcal W=(w^\ell)$ on
\[
\Omega_{\rm a}=[H_{\rm a}]\times[K_{\rm a}],
\]
define
\[
M_{\rm a}((v,c),(g,\alpha,\ell)):=
\begin{cases}
w^\ell(c_g,\iota_\alpha(v)),&\iota_\alpha(v)>0,\\
0,&\iota_\alpha(v)=0.
\end{cases}
\]
Rows are indexed by $(v,c)\in V\times\mathcal C$. Column labels lie in
$[q_{\rm a}]\times[H_{\rm a}]\times[|\mathcal W|]$. Equal arrays retain their
separate labels.
\end{definition}

In \cref{fig:yao-entry}, \(g\) selects an address from Alice's tuple, which
gives the row of Bob's array, and the position of her vertex in Bob's edge gives
the column.

\begin{figure}[H]
\centering
\begin{tikzpicture}[>=Stealth,every node/.style={font=\small},
 data/.style={draw,rounded corners=2pt,align=left,inner sep=5pt}]
\matrix (bits) at (8,0) [matrix of math nodes,
 nodes={minimum width=8mm,minimum height=6mm,inner sep=1pt},
 column sep=0pt,row sep=0pt,
 row 1/.style={nodes={font=\footnotesize}},
 column 1/.style={nodes={font=\footnotesize}}] {
  &1&2&3&4\\
 1&1&0&1&0\\
 2&0&1&0&0\\
 3&0&1&1&1\\
 4&1&0&0&1\\
};
\begin{scope}[on background layer]
  \node[fill=aliceblue!8,fit=(bits-2-2)(bits-2-5),inner sep=0pt] {};
  \node[fill=boborange!8,fit=(bits-2-4)(bits-5-4),inner sep=0pt] {};
\end{scope}
\node[draw=aliceblue,fit=(bits-2-1),inner sep=0pt] {};
\node[draw=boborange,fit=(bits-1-4),inner sep=0pt] {};
\node[draw=answergreen,very thick,fit=(bits-2-4),inner sep=0pt] {};
\node[data,draw=aliceblue,fill=aliceblue!5,text width=48mm,anchor=east]
 (address) at ([xshift=-17mm]bits-2-1.west)
 {Alice's tuple: $c=(3,\boxed{1})$\\
  Bob selects $g=2$\\[2pt]
  Address: $c_g=1$};
\node[data,draw=boborange,fill=boborange!5,text width=54mm,anchor=south]
 (position) at ([yshift=15mm]bits-1-4.north)
 {Bob's edge: $e_\alpha=(u,v,\boxed{z},r)$\\
  Alice's vertex: $z$\\[2pt]
  Position: $\iota_\alpha(z)=3$};
\draw[->,aliceblue,semithick] (address.east) -- (bits-2-1.west);
\draw[->,boborange,semithick] (position.south) -- (bits-1-4.north);
\node[below=5mm of bits,align=center]
 {Bob's array $w^\ell$\\[2pt]
  \textcolor{answergreen}{Requested bit: $w^\ell(1,3)=1$}};
\node[data,text width=29mm,anchor=north west]
 (outside) at ([yshift=-12mm]address.south west)
 {Vertex outside\\the selected edge};
\node[data,anchor=west] (zero) at ([xshift=9mm]outside.east) {output $0$};
\draw[->,semithick] (outside.east) -- (zero.west);
\end{tikzpicture}
\caption{The address selects a row and the vertex's edge position selects a
column, giving \(w^\ell(1,3)=1\). A vertex outside the edge gives zero without
reading the array. The indices are illustrative, and the arrows show how the
entry is selected, not a protocol's messages.}
\label{fig:yao-entry}
\end{figure}

\subsection{The protocol given by a coloring}

Once Bob has sent the edge, a balanced coloring lets Alice name her
position among the few positions of her color. The schedule in the proof of
\cref{prop:alternating-upper} exchanges color and coordinate, then edge and
address, then this position, and costs
\[
B_{\rm a}:=(2t-1)+2b+2(t+1)+1=2b+4t+2.
\]

\begin{proposition}\label{prop:alternating-upper}
For the table of \cref{def:alternating-table}, an $s_{\rm a}$-coloring with at
least $s_{\rm a}-1$ vertices of every color in every edge gives
\[
D_A^{\rm alt}(M_{\rm a})\le B_{\rm a}.
\]
\end{proposition}
\begin{proof}
Fix a row $(v,c)$, a column $(g,\alpha,\ell)$ and a balanced coloring $\chi$.
An edge contains at least
$s_{\rm a}-1$ vertices of each of the other $s_{\rm a}-1$ colors. Consequently
the vertices whose color is $\chi(v)$ occupy at most
\[
K_{\rm a}-(s_{\rm a}-1)^2=2s_{\rm a}-1
\]
positions.
We add one message code for ``outside the edge.'' There are then at most
$2s_{\rm a}=2^{t+1}$ possible message values, encoded in $t+1$ bits. We fix their ordering publicly
for every color and edge index. An unused edge index has only the value for ``outside the edge.''
Alice determines her message value from $v$, $\alpha$ and $\chi(v)$.

First Alice sends the $t$ bits of $\chi(v)$, while Bob sends
the $t-1$ bits of $g$ between consecutive Alice bits. This uses $2t-1$ bits and
leaves Alice knowing $g$. Next Bob sends the $b$ bits of $\alpha$, while Alice
sends the $b$ bits of $c_g$, alternating. This uses $2b$ bits, after which
Alice knows $\alpha$ and Bob knows $c_g$. Finally Bob sends $t+1$ dummy bits and Alice sends
the $t+1$ bits encoding her message value, after which Bob sends one answer bit. He answers zero for
``outside the edge'' and otherwise reads the specified array entry. The three
stages use $B_{\rm a}$ bits in total.
\end{proof}

\subsection{Few response patterns cannot cover many array positions}

Bob's early messages may depend on the whole array, so the transcript on one
column does not determine how Alice behaves on other columns. To analyze an arbitrary
protocol, we therefore also consider Bob histories that do not occur on a given column.
For a fixed row, we record Alice's $r$th bit after every possible history of
$r-1$ Bob bits, including hypothetical histories. Her own preceding bits
are then determined. For this count, we pad branches that stop early with
fixed bits in the alternating order. For example, her first
two turns require three response bits: the initial bit and one response to each
possible first Bob bit.
In general, her first $t$ turns require
$\sum_{r=1}^t2^{r-1}=2^t-1$ response bits, giving at most
$2^{2^t-1}$ possible response rules for these turns.

Among all Boolean arrays, the fraction agreeing with prescribed bits on
$k$ distinct positions is $2^{-k}$. We need only an upper bound of
$2^{-k}+\varepsilon$ on this fraction, and a small explicit list already
satisfies this weaker bound. Here we count all array labels equally, keeping
repeated arrays as separate labels. For a list
$\mathcal W=(w^\ell)_{\ell\in\mathcal L}$ on a finite position set $\Omega$,
with $\mathcal L$ finite and nonempty, we require
\begin{equation}\label{eq:alternating-pattern}
\frac{|\{\ell\in\mathcal L:w^\ell|_I=z\}|}{|\mathcal L|}
\le2^{-|I|}+\varepsilon
\qquad
(I\subseteq\Omega,\ z\in\{0,1\}^I).
\end{equation}
Adding these counts, at most a $P(2^{-|I|}+\varepsilon)$ fraction of all
labels can match any of $P$ patterns on $I$. Every fraction below is measured
against the original list size $|\mathcal L|$, even after a transcript has
removed some columns.

\Cref{yao:small-bias-escape} rules out a short protocol when many requested
positions share Alice's early response rules. The early messages cannot distinguish
those positions, so many array labels remain after them. The remaining
messages can produce only a few answer patterns, which cannot cover all
the surviving labels. In the lemma, $m$
is the total number of Alice turns, and $t$ is the number of initial Alice
turns during which her response rules agree across the rows for every possible
Bob history.

\begin{lemma}[Lower bound for alternating protocols]\label{yao:small-bias-escape}
Let $m,t$ be integers with $m\ge2t\ge2$. Suppose the nonempty labeled
array list $\mathcal W$ satisfies \cref{eq:alternating-pattern} with
$\varepsilon\le2^{-4m}$. If $J\subseteq\Omega$ has at least
$2^{m-t+1}$ positions, then the matrix
\[
L_J(j,\ell):=w^\ell(j)
\]
has no Alice-first alternating protocol of depth at most $2m-1$ in which all
rows give Alice the same response rules for her first $t$ turns on every possible Bob history.
\end{lemma}
\begin{proof}
Suppose such a protocol exists. We pad every shorter branch to depth $2m-1$
with fixed alternating bits. At the final Alice node, her row and the transcript
determine the old output leaf, so we replace her last bit by that output. This
uses no extra bit and, since $m>t$, preserves her first $t$ responses.

We stop with about $\log m$ turns per player remaining, so that Bob's
remaining response rules require fewer bits to describe than there are
surviving array positions.
Put $h=\ceil{\log m}\le m-t$ and $Q=2^h$. The inequality follows from
$m\le2(m-t)\le2^{m-t}$. We follow the first $2(m-h)-1$ messages and keep the
Bob child with more column labels and the Alice child with more rows.
Alice's first $t$ turns discard no rows because these response rules agree.
Each of her remaining $m-h-t$ turns retains at least half the rows,
and each of the $m-h-1$ Bob turns retains at least half the columns.
The resulting rectangle $X\times Y$, with $Y\subseteq\mathcal L$, therefore satisfies
\[
|X|\ge2^{h+1}=2Q,\qquad
\frac{|Y|}{|\mathcal L|}\ge2^{-(m-h-1)}.
\]
Exactly $h$ Bob--Alice pairs remain. We record Bob's next bit after every
hypothetical history of earlier Alice bits, which gives
$1+2+\cdots+2^{h-1}=Q-1$ bits. Fixing these bits and Alice's response rules determines
her final output on every retained row. On any $2Q$ of those rows, there are
therefore at most $2^{Q-1}$ output patterns.

Since $m\le Q<2m$, we have $\varepsilon\le2^{-2Q}$. Every column in $Y$
must match one of these patterns. Applying \cref{eq:alternating-pattern}
to each pattern and adding the bounds gives
\[
\frac{|Y|}{|\mathcal L|}\le2^{Q-1}(2^{-2Q}+\varepsilon)\le2^{-Q}.
\]
This contradicts the lower bound on the surviving fraction because
$Q\ge m>m-h-1$.
\end{proof}

Since \cref{eq:alternating-pattern} holds for every set of positions, one
array list works against every protocol. The bound applies to fractions of
the original list; it need not hold within the subset left by a transcript.
We use the accuracy
$\epsilon_{\rm a}=2^{-4(b+3t)}=2^{-12t}H_{\rm a}^{-4}$, which is inverse
polynomial in $H_{\rm a}$ for each fixed $t$. At this accuracy, the
finite-field construction of \cref{sec:alternating_implementation} gives such a
list with polynomially many arrays in $H_{\rm a}$, and we use it in place of the
set of all Boolean arrays in $M_{\rm a}$.

\begin{corollary}\label{yao:array-list}
For every fixed $t\ge4$, a nonempty labeled array list on $\Omega_{\rm a}$
satisfying \cref{eq:alternating-pattern} with
\[
\epsilon_{\rm a}:=2^{-4(b+3t)}
\]
can be generated in deterministic time polynomial in $H_{\rm a}$.
\end{corollary}

The construction and proof are in \cref{sec:arrays}.

\subsection{From a hypergraph to an explicit gap}

Given a short protocol, we color each vertex by the most common response
rules among its address tuples. A monochromatic edge then gives the many distinct
array positions that \cref{yao:small-bias-escape} requires.

\begin{proposition}\label{yao:lift}
Use the complete address list $\mathcal C$ and the explicit array list of
\cref{yao:array-list} in $M_{\rm a}$. A balanced source coloring gives
\[
D_A^{\rm alt}(M_{\rm a})\le B_{\rm a},
\]
whereas the absence of a proper $S_{\rm a}$-coloring gives
\[
D_A^{\rm alt}(M_{\rm a})\ge B_{\rm a}+2t-8.
\]
For each fixed $t$, the full table can be printed in deterministic time
polynomial in $|V|+H_{\rm a}$.
\end{proposition}
\begin{proof}
The upper bound is \cref{prop:alternating-upper}. For the lower bound, we put
$m=b+3t-3$, so $2m-1=B_{\rm a}+2t-9$, and suppose that an Alice-first protocol
has depth at most $2m-1$. We pad it to that depth and, for each vertex $v$,
partition its address tuples by Alice's response rules on her first $t$ turns.
By the count of response bits above, there are at most
$2^{2^t-1}\le S_{\rm a}$ choices of these rules. A largest class $C_v$ therefore has
\[
|C_v|\ge H_{\rm a}^{q_{\rm a}}/S_{\rm a}.
\]
We color $v$ by this class's response rules. Since the hypergraph has no proper
$S_{\rm a}$-coloring, some edge $e_\alpha$ is monochromatic, that is, all its
vertices receive the same response rules.

For $v\in e_\alpha$ and $g\in[q_{\rm a}]$, let
$J_{v,g}=\{c_g:c\in C_v\}$ be the set of addresses occurring in coordinate
$g$. Since $C_v$ is contained in the product of these projections,
\[
\prod_{g=1}^{q_{\rm a}}|J_{v,g}|\ge |C_v|
\ge H_{\rm a}^{q_{\rm a}}/S_{\rm a}.
\]
A column $(g,\alpha,\ell)$ reads only coordinate $g$ of Alice's tuple, so we
use a single coordinate $g$ for all vertices of the edge. AM--GM, summed over the vertices of $e_\alpha$, gives a coordinate $g$ with
\[
\begin{aligned}
N:=\sum_{v\in e_\alpha}|J_{v,g}|
&\ge K_{\rm a}H_{\rm a}S_{\rm a}^{-1/q_{\rm a}}\\
&=\frac{K_{\rm a}H_{\rm a}}4=2^{m-t+1},
\end{aligned}
\]
where $S_{\rm a}=4^{q_{\rm a}}$.
For each $v\in e_\alpha$ and $u\in J_{v,g}$, we choose a tuple $c(v,u)\in C_v$
with $c(v,u)_g=u$. We retain the corresponding rows and all columns
$(g,\alpha,\ell)$. Their entries are
\[
M_{\rm a}((v,c(v,u)),(g,\alpha,\ell))
=w^\ell(u,\iota_\alpha(v)).
\]
These are $N$ distinct positions, because different vertices occupy different
positions in the edge and different addresses give different first coordinates.
All retained rows give Alice the same response rules for her first $t$ turns. Since $m=b+3t-3\ge2t$
and $\epsilon_{\rm a}\le2^{-4m}$, the restriction to these rows and columns contradicts
\cref{yao:small-bias-escape}. Integer depths give
$D_A^{\rm alt}(M_{\rm a})\ge2m=B_{\rm a}+2t-8$.

There are $|V|H_{\rm a}^{q_{\rm a}}$ rows and
$q_{\rm a}H_{\rm a}|\mathcal W|$ columns. For fixed $t$, the number
$q_{\rm a}$ is constant and \cref{yao:array-list} constructs the arrays in
polynomial time. Enumerating all tuples and evaluating the displayed entry rule
therefore prints the entire table in polynomial time.
\end{proof}

To finish, we choose $t$ so that the gap $2t-8$ of \cref{yao:lift} is at least
the requested gap. We then transfer the resulting hardness to the other starter
and output conventions with \cref{yao:comparison}.

\begin{proof}[Proof of \cref{yao:main,yao:bpp}]
\phantomsection\label{proof:yao:main}\label{proof:yao:bpp}
Given a fixed requested gap $k\ge1$, we choose $t=\ceil{k/2}+4$. Then
$2t-8=2\ceil{k/2}\ge k$.
Composing the reduction of \cref{yao:source} with \cref{yao:lift}, we obtain,
in deterministic polynomial time, a table whose cost $D_A^{\rm alt}$ is at
most $B_{\rm a}$ in the YES case and at least $B_{\rm a}+k$ in the NO case, proving \cref{yao:main}.

An additive-$u$ estimate of any cost in \cref{yao:comparison} is, by the
triangle inequality, an additive-$(u+2)$ estimate of $D_A^{\rm alt}$ with the
same success probability. We choose any fixed integer $k>2(u+2)$ and compare
the estimate with $B_{\rm a}+k/2$. This gives
$\mathsf P=\NP$ for a deterministic estimator and $\NP\subseteq\BPP$ for a
bounded-error randomized estimator, proving \cref{yao:bpp}.
\end{proof}

The address length $q_{\rm a}=2^{t-1}$ and the degree of the source
reduction may both grow with $t$. Hence the construction is polynomial only for each
fixed gap, and this argument gives no growing-error result for strict
alternation.

\section{Remaining approximation questions}
\label{sec:conclusion}

We obtained fixed-error NP-hardness in both models and algorithmic implications
for growing errors in the unrestricted model. The quantitative consequences
depend on how fast the explicit table grows with the requested gap. Smaller
tables would give faster SAT algorithms from the same approximation guarantee,
and stronger inapproximability under assumptions excluding those SAT algorithms.

For unrestricted protocols, our deterministic construction outputs a table
of size at most $N^{2^{O(k)}}$ for an $N$-variable sparse formula and gap $k$
(\cref{det:reduction}). Can smaller lists preserve the restriction bounds needed for the
next round, or can another construction improve the $\log\log T$ error scale?

For strict alternation, a construction whose size and running time depend
more mildly on the gap could let the gap grow with the table size, giving
analogous approximation-to-SAT implications in Yao's model.


\appendix
\section{Interlacing estimates and the initial matrix}\label{sec:interlacing-details}
\label{sec:interlacing_details}

We first prove the projection and small-block estimates used in
\cref{sec:surviving-block-proofs}. The surviving-block bound proved there
then supplies the separation theorem and the column-loss induction below.
We use these estimates to prove \cref{lem:many-product-blocks,cl:bridge},
and finish with the initial matrix and template-layer bounds
(\cref{cl:seed,src:profiles}).

\subsection{Projecting column tuples and retained weight}

After a protocol prefix, a restriction of an interlace may retain many blocks
but only a small fraction of its column weight. We select a few blocks whose
projected tuples still have enough weight for a lower bound. We measure the
loss by the logarithm $\Delta_p$ below, which increases by one each time the
retained column weight is halved.

Fix a Boolean matrix $M:X\times Y\to\{0,1\}$ with positive normalized input weights $\mu$ and
$\nu$. Every weight below uses these original weights or their products.
For a nonempty $K\subseteq Y^p$, define
$\Delta_p(K)=\log(1/\nu^p(K))$. The projection onto no coordinates has loss $\Delta_0=0$, since $Y^0$ is
a single point of weight one.

We show that projection cannot increase the loss, and that projecting onto a
suitable set of $l$ of the $p$ coordinates leaves at most an $l/p$ fraction of it.

\begin{lemma}\label{cl:projection}
Let $K\subseteq Y^p$ be nonempty, and write $\pi_JK$ for its projection onto the
coordinates in $J\subseteq[p]$. Then
\begin{align*}
 \Delta_{\card J}(\pi_JK)&\le\Delta_p(K),\\
 \Delta_{\card{J_0}}(\pi_{J_0}K)+\Delta_{\card{J_1}}(\pi_{J_1}K)
 &\le\Delta_p(K) &&(J_0\sqcup J_1=[p]).
\end{align*}
For every integer $1\le l\le p$, some $l$-element set $J$ satisfies
\[
\Delta_l(\pi_JK)\le\frac lp\Delta_p(K).
\]
Viewing $K$ as a column set of $\Int{p}{M}$, the interlace on the blocks in
$J$, restricted to the columns $\pi_JK$ and to any retained rows of these
blocks, is a restriction of $\Int{p}{M}$ with columns in $K$ and the same
block row weights (\cref{obs:representatives}).
\end{lemma}

\begin{proof}
The first two inequalities follow from
$K\subseteq\pi_J^{-1}(\pi_JK)$ and
$K\subseteq\pi_{J_0}K\times\pi_{J_1}K$, using the product weights.

For the averaged bound, we interpret the fixed weights as a probability
distribution and let $Z$ have distribution $Q=\nu^p(\cdot\mid K)$.
For distributions $P$ and positive $P'$ on a finite set, their relative entropy is
$D_{\rm KL}(P\|P')=\sum_zP(z)\log(P(z)/P'(z))$, and entropy is
$\mathsf H(P)=-\sum_zP(z)\log P(z)$, where zero terms are omitted.
Conditioning on $K$ gives
\[
D_{\rm KL}(Q\|\nu^p)=\Delta_p(K).
\]
Since $Q_J$ is supported on $\pi_JK$, nonnegativity of relative entropy gives
\[
\Delta_{|J|}(\pi_JK)\le D_{\rm KL}(Q_J\|\nu^J).
\]
For a uniformly chosen $l$-subset $J$, the entropy chain rule gives
\[
\mathbb E_J\mathsf H(Z_J)
\ge\mathbb E_J\sum_{i\in J}\mathsf H(Z_i\mid Z_1,\ldots,Z_{i-1})
=\frac lp\mathsf H(Z).
\]
Also, since $\nu^J$ is a product, each coordinate contributes to
$-\mathbb E\log\nu^J(Z_J)$ exactly when it lies in $J$, which happens with probability $l/p$. Subtracting the entropy
inequality therefore yields
\[
\mathbb E_J\Delta_l(\pi_JK)
\le\mathbb E_JD_{\rm KL}(Q_J\|\nu^J)
\le\frac lpD_{\rm KL}(Q\|\nu^p)=\frac lp\Delta_p(K).
\]
Hence some $J$ satisfies the desired bound. The conditional distribution $Q$ need not be a
product, so we compute all projection weights with the original product weights.
The last assertion is \cref{obs:representatives}.
\end{proof}

\subsection{Two or three surviving blocks}

Later arguments, such as the proof of \cref{det:exact}, reduce a large
collection of surviving copies to two or three. We therefore first prove lower
bounds for these small collections. Each proof splits into cases according
to the first speaker and uses both column-weight thresholds of \cref{def:profile}.

\begin{lemma}[Small block bounds]\label{lem:small-blocks}
Suppose $\RP{M}{x}{H}$ with $H\ge2$. Restrictions of a full-product interlace of $M$ with the
following numbers of blocks, row weights and product column weights have at
least the indicated costs whenever the displayed row threshold is $\le1$. Here we measure row weight separately in each block.
\[
\begin{array}{c|c|c|c}
\text{blocks}&\text{row weight per block}&\text{product column weight}&\text{cost}\\
\hline
2&2x&9/50&H\\
2&2x&9/25&H+1\\
3&4x&9/50&H+1
\end{array}
\]
\end{lemma}
\begin{proof}
For the first row, \cref{cl:projection} gives a singleton column projection
of weight $\ge\sqrt{9/50}>3/10$. Applying $\RP{M}{x}{H}$ with $j=1$ to
that block gives cost $\ge H-1\ge1$,
so the restriction is nonconstant. If Alice speaks first, we follow that block
into a child retaining base row weight $\ge x$. The retained columns are
unchanged, so the suffix costs at least $H-1$. If Bob speaks first, his denser child has column weight $\ge9/100$
and hence a singleton projection of weight $\ge3/10$. The suffix at that
child again costs at least $H-1$. Including the first bit gives cost at least $H$.

For the second row, the first row already gives nonconstancy. If Alice
speaks, some singleton projection has column weight $\ge3/5$, and a child
retaining half the row weight of that block costs at least $H$. If Bob speaks, his denser
child has column weight $\ge9/50$, so the first row gives suffix
cost at least $H$. Including the first bit gives cost at least $H+1$.

For the third row, projecting onto any two blocks and applying the first row
shows that the restriction is not constant. If Alice speaks, we assign each block to a child
retaining at least half its row weight. One child receives two blocks of weight $\ge2x$,
and the first row gives suffix cost $H$. If Bob speaks, his denser child
retains column weight $\ge9/100$. The averaged projection inequality gives
two of its three coordinates with product weight at least
\[
(9/100)^{2/3}>9/50.
\]
Their row weights are still $\ge4x\ge2x$, so the first row again gives
suffix cost $H$. Including the first bit proves the claim. By
\cref{obs:representatives}, each projection is a restriction of the original matrix.
\end{proof}

The next theorem shows that a short protocol for a width-$u$ interlace must
begin with $\log u$ Alice bits, and that each resulting rectangle keeps most of
the row weight of one block and little of every other block. Its proof uses \cref{det:exact}, whose proof in
\cref{sec:conditioning-details} does not use this theorem.

\begin{theorem}\label{cl:separation}
Suppose $\RP{M}{\rho}{d}$ with $d\ge2$. Let $u\ge2$ be dyadic and
$1\ge\xi>2u\rho$. Interlace $M$ at width $u$ using an exactly eight-wise independent
tuple list with the base-column marginal, as in \cref{def:exact-tuples}.
Every restriction retaining base row weight at least $\xi$ in every block
and tuple-label weight at least $18/25$ has the following property: every
protocol of depth at most $d+\log u$ for it begins with $\log u$ Alice bits
along every path.
The $u$ nodes at depth $\log u$ are in bijection with the blocks: each such
node retains more
than $\xi-2(u-1)\rho$ of its assigned block and less than $2\rho$ of every
other block. Full products are a special case, and fixed bits count as transmitted bits.
\end{theorem}

\begin{proof}
We maintain the following invariant at each protocol-tree node of depth
$r<\log u$: there are $\kappa=u2^{-r}$ assigned blocks, each retaining base row
weight $\ge\xi2^{-r}>2\kappa\rho$, and the retained column set is the original
one. The root has this property. By
\cref{obs:representatives}, the projection onto the coordinate of one assigned
block is a restriction of $M$ with column weight $\ge18/25>3/5$, so $\RP{M}{\rho}{d}$ gives cost at least $d\ge2$ and the
node is not a leaf. Since $1\ge\xi>2u\rho$, every row threshold $2^a\rho$ with
$1\le a\le\log u$ used below is less than one.

If Bob speaks, one child retains column weight $\ge9/25$ and all
$\kappa$ assigned blocks. By \cref{det:exact} with $a=\log\kappa$,
it needs $d+\log\kappa$ bits, but has only
$d+\log u-r-1=d+\log\kappa-1$ left. Thus Alice speaks. We assign each block
to a child retaining at least half its weight. If a child received more than
$\kappa/2$ blocks, it would contain at least $\kappa/2+1$ blocks of weight
$\xi2^{-r-1}>\kappa\rho$, and \cref{det:exact} with $a=\log\kappa$ would again
require $d+\log\kappa$ bits, exceeding the $d+\log\kappa-1$ left.
Both children therefore receive exactly $\kappa/2$ blocks. This proves the
invariant, and hence that Alice sends the first $\log u$ bits. In particular, a fixed Alice
bit is impossible, since it would assign every block to one child.

At depth $\log u$, each part has one assigned block of weight greater than
$2\rho$. If such a node retained two blocks of weight at least $2\rho$, they
would require $d+1$ bits by
\cref{det:exact} with $a=1$, whereas only $d$ remain. Every other block thus has
weight less than $2\rho$. Each original block is assigned to exactly one part, so
its weight elsewhere totals less than $2(u-1)\rho$, proving the final bound.
\end{proof}

\subsection{Induction on column loss}
\label{sec:interlace-induction}

For a restriction of $\Int{p}{M}$ with surviving columns $K$, Bob's messages
increase the loss $\Delta_p(K)$. The following induction gives lower bounds
that tolerate this loss, and it uses only lower bounds on the base matrix.

\begin{lemma}\label{cl:defect-induction}
Let $C\ge3$, $0<\theta\le1/4$, and $0<\delta\le1/64$.
Suppose $\RP{M}{\theta}{C-1}$ (\cref{def:profile}) and every restriction of $M$ retaining row
weight at least $\theta$ and column weight at least $1-\delta$ has cost at least $C$.
Let $e\in\{0,1\}$, let $a\ge2e$ be an integer with $2^a\theta\le1$, and let
$p>2^{a-1}$ be an integer. If a restriction $G$ of $\Int{p}{M}$ retains base
row weight at least $2^a\theta$ in every block and its nonempty column set $K$ satisfies
\[
\Delta_p(K)\le(p-2^{a-1})\delta+\frac{4e}{3},
\]
then $\D(G)\ge C+a-e$. Here $C$ is only a lower bound, and we do not assume
$\D(M)=C$.
\end{lemma}

\begin{proof}
Say that $G$ is in state $(a,e)$, or $(a,e,p)$, when it satisfies the
hypotheses for these values. We prove the lower bound by induction on
\[
\mathfrak m(a,e):=2a+1-e.
\]
An Alice bit halves the row thresholds and lowers $a$.
A Bob bit increases the loss by at most one. We cover this increase either
by raising $e$ to one, which adds $4/3$ to the allowed loss, or by projecting
onto about half the blocks with \cref{cl:projection}.

With $d=C-1\ge2$ and $x=\theta$, the bound $\RP{M}{\theta}{C-1}$ lets us apply
\cref{det:exact} to full products, and $4\theta\le1$. At column weight $9/25$, its cases $a=1,2$
give cost at least $C$ for two blocks of row weight $2\theta$, and at least
$C+1$ for three blocks of row weight $4\theta$. We call these the two-copy and
three-copy bounds. Every projection below with loss $\le3\delta+4/3$
has column weight $>9/25$, because
\[
3\delta+\frac43\le\frac{265}{192}<\frac75<-\log(9/25).
\]
The last inequality follows from \(2^7\cdot9^5<25^5\).

We first verify the base states. For \((a,e)=(0,0)\), the averaged
projection bound in \cref{cl:projection} gives a singleton \(J\subseteq[p]\)
with
\[
\Delta_1(\pi_JK)\le\frac{p-1/2}{p}\delta<\delta.
\]
Since \(2^{-v}\ge1-v\) for \(0\le v\le1\), applied with \(v=\Delta_1(\pi_JK)\), the projected column weight is
greater than \(1-\delta\). The block in \(J\) retains row weight
\(\ge\theta\), so the hypothesis gives cost at least \(C\).

For \((a,e)=(2,1)\), the loss bound and \cref{cl:projection} give a
three-coordinate projection with loss at most
\[
\frac3p\left((p-2)\delta+\frac43\right)
\le3\delta+\frac43<-\log(9/25).
\]
Its column weight exceeds \(9/25\), and each block retains weight
\(\ge4\theta\). The three-copy bound above gives cost at least
\(C+1=C+2-1\).
There is one further base case, \((a,e,p)=(3,1,5)\). Its loss is
\(\le\delta+4/3<-\log(9/25)\), and its rows have weight \(\ge8\theta\).
Its column weight therefore exceeds \(9/25\), and \cref{det:exact} with $a=3$
applies to these five blocks, giving $C+2=C+3-1$ bits.

We also handle the Bob-first case of \((a,e)=(1,0)\) now;
its Alice-first case is covered by the general step below. A singleton projection has loss
\(\le(p-1)\delta/p<\delta\), so its column weight exceeds \(1-\delta\).
Its rows retain weight \(\ge2\theta\ge\theta\), so the hypothesis gives
cost at least \(C\ge3\) and the restriction is nonconstant. After a Bob bit,
the denser child has a two-coordinate projection of loss at most
\[
\frac2p\left((p-1)\delta+1\right) \le2\delta+1<3\delta+\frac43,
\]
so the two-copy bound gives cost at least \(C\) after that bit. Every other
valid state has \(a\ge2\), \(p\ge3\), and base row weight \(\ge4\theta\). Some
three-coordinate projection has loss \(\le3\delta+4/3\), so by the
three-copy bound the restriction is not constant.

Consider any remaining case and the first bit of a protocol for it. If Alice
speaks, we assign every block to a child retaining at least half its row weight. For
child \(i\in\{0,1\}\), let \(p_i\) be the number of assigned blocks and let
\(L_i\) be the loss after projecting the retained column set onto those
blocks. Then \(p_0+p_1=p\), and the partition inequality in \cref{cl:projection}
gives \(L_0+L_1\le\Delta_p(K)\). Put
\[
s_i=\max\{p_i-2^{a-2},0\}.
\]
Since \(s_0+s_1\ge p-2^{a-1}>0\), at least one \(s_i\) is positive. The
loss bounds required of the children with \(s_i>0\) satisfy
\[
 \begin{aligned}
\sum_{i:s_i>0}\left(s_i\delta+\frac{4e}{3}\right)
 &\ge (p-2^{a-1})\delta+\frac{4e}{3}\\
&\ge\sum_{i:s_i>0}L_i.
 \end{aligned}
\]
At least one eligible child contributes the constant \(4e/3\), and the
final inequality uses the assumed loss bound on \(K\) and
\(L_0+L_1\le\Delta_p(K)\). Consequently some child with \(s_i>0\) has
\(L_i\le s_i\delta+4e/3\). Each assigned block retains weight \(\ge2^{a-1}\theta\), and
there are \(p_i>2^{a-2}\) such blocks. This is state \((a-1,e)\), and the exceptional
state \((a,e)=(2,1)\) was already proved. The induction hypothesis gives cost at least
\(C+a-1-e\) in that child, and Alice's bit gives \(C+a-e\). The same argument
covers the step from state \((1,0)\) to state \((0,0)\).

If Bob speaks, one child has at least half the current column weight, so its loss
is \(\le\Delta_p(K)+1\). When \(e=0\) and \(a\ge2\), this child satisfies
state \((a,1)\), because
\[
\Delta_p(K)+1\le(p-2^{a-1})\delta+\frac43.
\]
Its induction cost is at least \(C+a-1\), and adding Bob's bit proves the claim.

It remains to treat a Bob bit when \(e=1\) and \(a\ge3\). The extra term
\(4/3\) has then already been used, so we instead project the denser child onto
about half the blocks and show that the projection is in state \((a-1,1)\). The five-block base case already
handles \(p=5\), so \(p\ge6\). We keep \(p'=\ceil{p/2}\) blocks. Then
\[
\frac12\le\frac{p'}p\le\frac47,\qquad
\frac{p'}p(p-2^{a-1})\le p'-2^{a-2},\qquad
\frac47\left(\frac43+1\right)=\frac43.
\]
By the averaged projection bound in \cref{cl:projection}, some \(p'\)
coordinates of the denser child therefore have loss at most
\[
\frac{p'}p\left((p-2^{a-1})\delta+\frac43+1\right)
\le(p'-2^{a-2})\delta+\frac43.
\]
The projected child retains base row weight \(\ge2^{a-1}\theta\) in each
of its \(p'>2^{a-2}\) blocks, so it is state \((a-1,1)\). Its induction cost is at
least \(C+a-2\), and adding Bob's bit gives \(C+a-1\), as required.
\end{proof}

We prove \cref{lem:many-product-blocks} by applying \cref{cl:defect-induction}
to $p=9u/16$ blocks. The induction needs more than $u/2$ blocks, and the
extra $u/16$ blocks give slack $(9u/16-u/2)\delta=u/1024$ at $\delta=1/64$ in the
loss bound. The retained column weight may therefore be as small as $2^{-u/1024}$.

\begin{proof}[Proof of \cref{lem:many-product-blocks}]
\phantomsection\label{proof:lem:many-product-blocks}
Let $K\subseteq Y^q$ be the retained column set, where $q$ is the width of the
interlace. Since $u\ge16$ is dyadic, $9u/16$ is an integer. We select a set
$J$ of exactly $p=9u/16$ blocks retaining base row weight $\ge ux$ and project
$K$ onto them. The weight bound on $K$ and \cref{cl:projection} give
$\Delta_p(\pi_JK)\le\Delta_q(K)\le u/1024$, and the projected interlace is a
restriction of the original one.
Apply \cref{cl:defect-induction} with
\[
a=\log u,\quad \theta=x,\quad \delta=1/64,\quad C=H,\quad e=0.
\]
Here $x\le1/16$, $p>2^{a-1}=u/2$, and $2^ax=ux\le1$.
The hypothesis $\RP{M}{x}{H}$ implies the weaker bounds $\RP{M}{x}{H-1}$
and, since $63/64\ge3/5$, cost $\ge H$ at column weight $63/64$. These are the
two lower-bound hypotheses of \cref{cl:defect-induction}. Finally,
\[
\Delta_p(\pi_JK)\le u/1024
=(9u/16-u/2)/64=(p-2^{a-1})\delta.
\]
The induction gives cost at least $C+a=H+\log u$.
\end{proof}

We use the same induction to restore both deletion bounds at full cost for
$R^+=(\Int{Q}{W})^{\mathsf T}$.

\begin{proof}[Proof of \cref{cl:bridge}]
\phantomsection\label{proof:cl:bridge}
Put $C=\D(W)$ and $f=\ceil{b/\delta}$. The assumptions imply
\[
f\delta\ge b,\qquad (f-1)/Q<1/80,\qquad f<Q/40,\qquad Q\ge5120.
\]
Indeed, $b/\delta\ge64$, $Q\delta\ge80b$, and
$f\le b/\delta+1<2b/\delta$. Fix $j\in\{0,1\}$ and put
\[
a_j=r-j,\qquad p_j=2^{a_j-1}+f,\qquad t_j=2^{a_j}\theta.
\]
These choices give $a_j\ge0$, $2^{a_j-1}<p_j<Q$, and $t_j<1$.

Consider a restriction of $R^+$ to which $\RP{R^+}{2^{-b}}{C+r}$ applies with
this $j$, and transpose it to a restriction of $\Int{Q}{W}$. By
\cref{def:profile}, the latter retains tuple
weight $\ge2^{-b}$ and total block-row weight $\ge(3/5)2^{-j}$.
If fewer than $p_j$ blocks had base row weight at least $t_j$, total row weight
would be less than
\[
\frac{p_j-1}{Q}+t_j
<2^{-j-1}+\frac1{80}+\frac1{20\cdot2^j}
\le\frac35\,2^{-j},
\]
a contradiction, where the last inequality uses $j\le1$. We select $p_j$ blocks $J$ meeting the row-weight threshold. For the retained tuple set $K$, \cref{cl:projection} gives
\[
\Delta_{p_j}(\pi_JK)\le\Delta_Q(K)\le b
\le f\delta=(p_j-2^{a_j-1})\delta.
\]
The projected restriction therefore satisfies the hypotheses of
\cref{cl:defect-induction} with $e=0$. Its two lower-bound hypotheses are
those assumed for $W$, $\theta<1/4$ follows from $Q\theta\le1/20$, and the
row threshold is $t_j$. The lemma gives cost at least $C+a_j=C+r-j$, so \cref{obs:representatives} and
transposition prove $\RP{R^+}{2^{-b}}{C+r}$.

Sending the copy index and then running a protocol for $W$ costs
$r+\D(W)$ bits. This gives equality on the whole matrix, and it turns any
upper bound on $\D(W)$ into an upper bound larger by $r$ on $\D(R^+)$. By definition the
transpose has $n^Q$ rows and $Qm$ columns.
\end{proof}

\subsection{The initial matrix and template layers}
\label{sec:initial-template-proofs}

For the initial matrix $R_0(z,j)=z_j$ on $\{0,1\}^{p_0}\times[p_0]$, we use
rank. Even after column indices are deleted, the restricted matrix has large
rank, so every protocol for it has many leaves of each output value.

\begin{proof}[Proof of \cref{cl:seed}]
\phantomsection\label{proof:cl:seed}

Since \(p_0\) is a power of two, \(a=\log p_0\) is an integer. Fix
\(j\in\{0,1\}\), a retained row set
\(V\subseteq\{0,1\}^{p_0}\) with \(\card V\ge2^{p_0-L}\), and a retained column
set \(J\subseteq[p_0]\) with \(\card J\ge\ceil{(3/5)2^{-j}p_0}\). Under uniform
weights, these are the restrictions tested by \(\RP{R_0}{2^{-L}}{d_0}\). Let \(v\) be
the dimension of the real span of \(V\). The matrix with row set \(V\) has
rank \(v\), so by \citet[Lemma~4.1(1)]{mackenzie2026condensation} it has at most
\(2^v\) distinct rows, that is, \(\card V\le2^v\). The
lower bound on \(\card V\) therefore implies \(v\ge p_0-L\).

Deleting the \(p_0-\card J\) coordinates outside \(J\) can lower the row rank by
at most \(p_0-\card J\). Thus the restricted initial matrix has rank \(r\)
satisfying
\[
r\ge \card J-L \ge \frac35\,2^{-j}p_0-L >2^{a-j-1}.
\]
The strict inequality follows from
\[
\frac35\,2^{-j}p_0-2^{a-j-1} =\frac{p_0}{10\cdot2^j}\ge\frac{p_0}{40}>L,
\]
where \(p_0/40-L>0\) follows from \(p_0\ge64(L+1)\).

By \citet[Lemma~4.1(2)]{mackenzie2026condensation}, a matrix of
depth at most \(a-j\ge1\) has rank at most \(2^{a-j-1}\). Hence the restriction has
depth at least \(a+1-j\). At \(j=0\), this gives \(\D(R_0)\ge a+1\). Conversely,
Bob sends the column index in \(a\) bits and Alice sends the corresponding bit, so
\(\D(R_0)\le a+1\), and hence \(\D(R_0)=a+1=d_0\). The unit-vector rows form the
identity submatrix, so \(\rank(R_0)=p_0=2^{d_0-1}\).

By minimality of the dyadic choice, \(p_0<128(L+1)\). If \(L=\ceil{\log n}\) and
\(n\ge4\), then \(L+1\le\log n+2\le2\log n\), so \(p_0<256\log n\) and
\(2^{p_0}<n^{256}\). Enumerating the \(2^{p_0}\) bit strings and the \(p_0\) column
indices therefore prints the table deterministically in polynomial time.
\end{proof}

\begin{proof}[Proof of \cref{src:profiles}]
\phantomsection\label{proof:src:profiles}
We treat both cases together, except for the step that turns many blocks
above the threshold into a lower bound. Fix $j\in\{0,1\}$, put
$\beta_j=(3/5)2^{-j}$, and let $\lambda=255/256$ be the total weight of the
template rows (\cref{def:source-matrix}). After transposition, each deletion
bound concerns an interlace with the data below. Here $q$ is the interlace
width, $u$ is the block parameter to which we apply
\cref{lem:many-product-blocks} or \cref{det:rareboost}, and the last column is
a lower bound on the average block retention.
\[
\begin{array}{c|c|c|c|c}
\text{target}&\text{base}&q&u&\text{average block retention}\\ \hline
A&R&q_1=s-2&s2^{-j}&\beta_j\\
R_{\rm temp}&A&q_2&q_2 2^{-j}&\beta_j\\
N_{\rm ref}&R_{\rm temp}&s&s2^{-j}&(\beta_j-(1-\lambda))/\lambda
\end{array}
\]
In each case the retained tuple-label event $Y$ has weight $\ge x$. For
$N_{\rm ref}$, we discard the non-template rows. They have total weight
$1-\lambda$, which gives the displayed average retention. We then map each retained $\Omega$
label to its bin part, the tuple $(j_1,\ldots,j_s)$ in the full-product case
and its $\mathcal J_{\rm g}$ label in the explicit case, forgetting the tail
$\xi$, which template entries ignore. This map preserves the original
weights, so its image still has weight $\ge x$, and choosing one retained
label for each image label gives the template entries (\cref{obs:representatives}).

Let $m$ be the number of blocks retaining base row weight at least the
threshold $t=u\rho$, $ux$, and $ux$, respectively. The block-counting bound
$q\beta\le m+qt$ of \cref{sec:models} gives the following lower bounds for $m/u$:
\[
\frac35(1-2/s)-q_1\rho,\qquad
\frac35-q_2x,\qquad
\frac{2^j(\beta_j-(1-\lambda))}{\lambda}-sx.
\]
All three exceed $9/16$ by the numerical bounds of \cref{lem:common-schedule}, which
both parameter schedules satisfy. They also give the row thresholds
$u\rho\le s\rho\le1$ and $ux\le\max\{q_2x,sx\}\le1$, and $u\ge s/2\ge16$
since $s\ge256$.

\emph{Full-product case.} Here $x=2^{-8L}\ge2^{-u/1024}$ by
\cref{lem:full-schedule}, so \cref{lem:many-product-blocks} applies to the
more than $9u/16$ blocks above the threshold.

\emph{Explicit case.} We apply \cref{det:conditioning} to $Y$ with $v=x$ and
$\tau=1$. Outside at most $e(x)$ exceptional coordinates, the conditional
marginals of $\nu\mid Y$ stay close to the product weights, and
$e(x)\le s/128\le u/32$ by \cref{lem:generated-schedule}. Discarding the blocks
at these coordinates leaves more than $9u/16-u/32>u/2$ blocks above the
threshold. Since $u$ is dyadic, at least $u/2+1$ of them remain. These blocks
and $K=Y$, for which $(\nu\mid Y)(K)=1$, meet the hypotheses of
\cref{det:rareboost} with $a=\log u$, which gives the costs computed below.

In both cases the base bound is $\RP{R}{\rho}{d}$ for the first row and, for
the other two rows, the deletion bounds just established for $A$ and
$R_{\rm temp}$. The required costs are therefore
\[
d+\log s-j=H_1-j,\qquad
H_1+\log q_2-j=H_2-j,\qquad
H_2+\log s-j=C-1-j,
\]
which are the three asserted pairs of deletion bounds.

It remains to prove the collision bound for the interlace on $\Gamma$.
The first $q_1$ positions and three distinct tail positions give
$q_1+3=s+1=(2s)/2+1$ blocks. By hypothesis, each retains base row weight
$\ge x_{\rm cap}=2s\rho$, and the retained list labels have weight greater
than $1/2>9/25$. In either case the list $\Gamma$ is exactly eight-wise
independent with the base-column marginal, so \cref{det:exact} with
$a=\log(2s)$ gives $d+\log(2s)=H_1+1$ bits.
\end{proof}

\section{Bounds on transferred inputs}\label{sec:mass-details}
\label{sec:mass_bounds}

We prove \cref{res:incoming,res:foreign}, which bound the $\mathcal D$-weight
arriving at a stopping node (defined in \cref{res:destinations}) and the
weight sent to the stopping node assigned to another copy. Throughout, we
assume the common hypotheses (H1)--(H4) of \cref{sec:internal-interlace} and
drop the subscript $\rm g$ of the explicit construction. In each source copy,
the $\mathcal D$-inputs have total weight $\mu=1/256$, and the template rows
and the $\Omega$-inputs each have total weight $\lambda=255/256$.

Transferred weights below are in units of one complete copy, whereas
$\varepsilon_i$ is normalized within copy $i$'s $\Omega$-inputs.
\Cref{det:conditioning} controls conditional marginals on at most eight
nonexceptional coordinates; the projection inequality following it converts
this control into bounds in the original product weights.
Choosing retained representatives as in \cref{obs:representatives}
realizes these projected tables as actual restrictions.

\subsection{Two template inputs at one stopping node}

The next lemma uses the transposed orientation: its rows
are template-label pairs, and its two column families come from one
copy's $\Omega$-inputs and another copy's $\mathcal D$-inputs.

At the node assigned to copy $i$, an $\Omega$-input of copy $i$ reads the
template label at coordinate $i$, and a $\mathcal D$-input of another copy $j$
reads the template label at coordinate $j$. By (H3), these two labels are
independent under the original weights. The next lemma shows that if both
families are retained, the protocol needs one bit more than the cost $H_2$
given by $\RP{R_{\rm temp}}{x}{H_2}$, even if the two families have the
same column labels in $J_2$.

\begin{lemma}[Two template inputs]\label{res:paired}
Let $A:\Gamma\times J\to\{0,1\}$ satisfy $\RP{A}{x}{H_1}$ with $H_1\ge2$,
and let $q_2$ be dyadic with $2q_2x\le1/8$. Take two weighted lists $S_o,S_f$ of words in $J^{q_2}$, each
exactly eight-wise independent with the base-column marginal, and use
product weights on $S_o\times S_f$. Put $H_2=H_1+\log q_2$ and give
$J_2=[q_2]\times\Gamma$ the product of uniform index weights and the row weights of $A$ on $\Gamma$.
Rows are pairs of template labels $(z_o,z_f)\in S_o\times S_f$. The columns form two separately tagged copies of $J_2$, with entries
\[
A(\gamma,(z_o)_\alpha)\quad\text{and}\quad
A(\gamma,(z_f)_\alpha),
\]
respectively. The tags remain distinct even when the inner indices coincide.
Retain a set of pairs of original product weight at least $49/64$, columns reading $z_o$
of normalized weight at least $1-\varepsilon$, and columns reading $z_f$ of
normalized weight at least $\delta$, each tagged copy of $J_2$ having total weight one, where
\[
\varepsilon\ge0,\qquad 0<\delta\le1,\qquad
\delta>2q_2x+\varepsilon.
\]
Then the retained matrix requires at least $H_2+1$ bits.
\end{lemma}

\begin{proof}
Put $t=2q_2x$. We group columns by their tag and index $\alpha$.
These $2q_2$ groups have total retained $\Gamma$-weight at least
\[
q_2(1-\varepsilon+\delta)>q_2(1+t).
\]
If at most $q_2$ groups had weight at least $t$, this sum would be at most
$q_2+q_2t$. Thus at least $q_2+1$ groups retain weight $\ge t$.

We transpose the matrix. These groups are now distinct row blocks of an
interlace of $A$. Concatenating $S_o$ and $S_f$ with product weights gives
an exactly eight-wise independent tuple list on $2q_2$ tagged coordinates,
since any set of at most eight coordinates splits between the two independent
lists, and each list has the required marginals on its part. We keep the surviving pair
labels as tuple labels, whose original weight is $\ge49/64>9/25$, and
apply \cref{det:exact} with the integer $a=\log(2q_2)$ and row threshold
$2q_2x\le1/8$. This gives $H_1+\log(2q_2)=H_2+1$ bits, as required.
\end{proof}

In the application (the proof of \cref{res:foreign}), we fix a bin at each
of the two copies, and the pairs $(z_o,z_f)$ occur only in the bin fixed at
copy $i$, which we call the \emph{distinguished bin}. On the other bins we know
nothing about the entries in the other copy's columns, so the next claim
extends \cref{res:paired} to allow arbitrary entries there.

\begin{claim}\label{cl:two-stream}
Use the common source weights and deletion bounds of \cref{sec:internal-interlace}, and put
$\lambda_0=2q_2x$, $\tau_0=2sx$, $C=H_2+\log s+1$, and
\[
a_{\min}=\frac{99/100+\lambda^2-1}{\lambda^2}.
\]
Retain columns of normalized $\Omega$-weight at least
$1-\varepsilon$, and separately tagged other-copy columns whose
weight-preserving map $(\operatorname{tail}(\gamma),j,\ldots,j)\mapsto j$ to $J_2$
has retained image weight at least $\delta$.
Number the bins so that the distinguished bin is bin one. It retains a set
$P_*\subseteq S^2$ of pairs $(z_o,z_f)$ of product weight at least $a_{\min}$, and every other bin retains
a set of labels of $S$-weight at least $\tau_0$. Write coordinates of a column in $\Omega$ as
$Z=(\xi,j_1,\ldots,j_s)$, where $\xi\in Y^{s+1}$ and $j_p=(\alpha_p,\gamma_p)\in J_2$. Require entries
\[
\begin{array}{c|cc}
 & Z\in\Omega&\text{other-copy image }(\alpha,\gamma)\\ \hline
(z_o,z_f)\in P_*&A(\gamma_1,(z_o)_{\alpha_1})&A(\gamma,(z_f)_\alpha)\\
(p,z),\ p\ne1&A(\gamma_p,z_{\alpha_p})&\text{arbitrary}.
\end{array}
\]
If $\varepsilon\ge0$, $0<\delta\le1$, and $\delta>\lambda_0+4\varepsilon$, then every such matrix
requires at least $C$ bits.
\end{claim}

\begin{proof}
Suppose this matrix has a protocol of depth at most $C-1$. In each bin, we
follow the retained set of template labels through Alice's bits until Bob
first speaks, using the tracking rule stated before
\cref{lem:collision-tracking} with $t=\tau_0$. The only marked set is the set
$P_*$ of pairs in the distinguished bin. For the thresholds, we measure its
weight by its projection onto $z_o$, and a
split replaces its two parts by their projections onto $z_o$, which are sets
of labels from $S$. A pair set weighs no more than its projection onto $z_o$,
so an unsplit pair set loses less than $\tau_0$ of its original weight.
The common numerical bounds give
$\tau_0\le10^{-3}$ and $a_{\min}-\tau_0>49/64$. An unsplit set of pairs therefore
retains weight $>49/64$, and every retained set at depth $r$ has original
weight $\ge\tau_0 2^{-r}$.

The hypothesis on $\delta$ implies $\varepsilon<1/4$, so the columns from $\Omega$
retain normalized weight $>3/4$, and so does their projection onto any bin.
Since $H_2>0$, the bounds $\RP{R_{\rm temp}}{x}{H_2}$ rule out a leaf
containing a retained set of weight $\ge x$. A retained set also cannot reach depth $\log s+1$:
it would have weight at least $x$, but only $H_2-1$ bits would remain.
Thus every retained set reaches a Bob node by depth $\log s$.

At each stopping node $v$, let $h_v$ be the number of tracked sets and $z_v$
the number of unsplit pair sets among them. By the tracking rule,
$\sum_v(h_v+z_v)$ equals $s+1$, the $s$ initial sets plus the one initial
mark. Thus the stopping nodes satisfy the hypotheses of
\cref{lem:kraft-count} with $q=s$.
Let $v$, at depth $r$, be the node given by that lemma, and write $\kappa=\kappa_v$,
$h=h_v$, $z=z_v$. The sets at $v$ have distinct bin coordinates, since
descendants in different branches never rejoin, and each has original weight
$\ge\tau_0 2^{-r}=2\kappa x$. By \cref{lem:kraft-count},
$h\ge\kappa/2+1$ if $\kappa\ge2$, while for $\kappa=1$, either $h\ge2$ or $h=z=1$.

We first treat the case $\kappa=1$, $h=z=1$, before restricting columns or projecting the pair.
Since only Alice has spoken, both tagged column families are still unchanged. We project columns from $\Omega$
onto $j_1$ and other-copy columns onto their $J_2$ images, choosing one surviving
column per image label. These maps preserve the original weights, so the projected weights
are at least $1-\varepsilon$ and $\delta$. The set of pairs has weight
$>49/64$. Thus \cref{res:paired} requires $H_2+1$ bits, while the stopping
node has only $C-1-r=H_2$ bits left. The original weights on pairs of template
labels are the product of their weights on $S$, as required by that lemma.

In all other cases, at Bob's bit we choose the child with more weight from
$\Omega$, which retains weight $>3/8>9/25$. We replace any set of pairs by its
projection onto $z_o$, which has at least its weight. The entries now form a
restriction of an interlace of $R_{\rm temp}$ on the tracked bins. Its original
column list is the bin part of $\Omega$, which by (H3) has exactly
$8192$-wise independent bin coordinates under the original weights in either
construction. We forget the tail $\xi$, which template entries ignore, but keep
the whole tuple $(j_1,\ldots,j_s)$ in the
full-product case and its whole list index in the explicit case. The retained
image has original weight $>3/8$, and choosing one surviving row or column
for each projected label gives the required entries (\cref{obs:representatives}). Thus
\cref{det:exact} applies in both constructions. For $\kappa\ge2$ we
use width $\kappa$. The $\kappa/2+1$ blocks have row weight
$\ge2\kappa x\ge\kappa x$, so the lemma requires $H_2+\log\kappa$ bits, but
the child has only $H_2+\log\kappa-1$ bits left. For $\kappa=1$ there are at least two
blocks of weight $\ge2x$, and width two gives the stronger bound $H_2+1$.
By the parameter choices, all these row thresholds are at most one. This
contradiction does not use the other-copy entries outside the distinguished
bin, which proves the claim.
\end{proof}

\subsection{Total weight arriving at a stopping node}
\label{sec:incoming-bound-proof}

We first bound the total $\mathcal D$-weight arriving at one stopping node,
including weight from copies without an assigned node. Too much incoming
weight leaves many blocks of substantial row weight, and \cref{det:rareboost},
applied after conditioning, then requires more communication than remains.

\begin{proof}[Proof of \cref{res:incoming}]
\phantomsection\label{proof:res:incoming}
Put $\kappa=1$ at a node at depth $\log w$. Fix any collection of copies, and suppose
the weight of their inputs from $\mathcal D$ arriving at one stopping node is at least $2s\kappa$.
Let $d_i\in[0,1]$ be the fraction of copy $i$'s inputs from $\mathcal D$ reaching it,
and set $d_i=0$ outside this collection. Since each such family has weight
$\mu=1/256$,
\[
D_0:=\sum_i d_i\ge512s\kappa.
\]
These fractions are fixed by the Alice-only transcript. At a node where Bob first speaks before depth $\log w$,
we choose a child whose column set $Y^*$ has original weight $\ge\theta/2$,
which is possible since the tuple event $Y'$ of \cref{res:destinations} has weight $\ge\theta$.
At a node at depth $\log w$ we put $Y^*=Y'$.

To apply the template lower bounds, we need retained column tuples in which
at least half of the incoming weight $D_0$ comes from copies whose label is a
template row.
Let $T_i$ indicate that the label at coordinate $i$ of a column tuple is a template row. When the type, bin, and item labels are drawn according to their fixed weights, the centered variables $T_i-\lambda$ are $256$-wise independent.
The weighted form of \cref{lem:conditioning-moment}, followed by Markov's
inequality, gives
\[
\Pr\left\{\sum_i d_iT_i<D_0/2\right\}
\le C_{128}D_0^{-128}\le C_{128}s^{-128}\le\theta/4.
\]
Here $C_{128}=(2^{14}/(\lambda-1/2)^2)^{128}$ is the fixed constant
of \eqref{eq:common-budgets}, and the last inequality is a row of \eqref{eq:width-margins}.
Thus $Y^*$ intersects the event $\sum_i d_iT_i\ge D_0/2$ in weight
$\ge\theta/4$.
We average over the type/bin/item sequence and the base-label sequence, and fix a
configuration whose retained set $Z$ of template-label sequences has original weight
$\ge\theta/4$ and whose template positions have total weight
$\sum_i d_iT_i\ge D_0/2$.
By (H3), fixing those two sequences leaves the template-label sequence
with its original weights and exactly $8192$-wise independent
$J$ coordinates $z_{i,\alpha}$. We apply
\cref{det:conditioning} to $Z$ with $\tau=1$, which removes at most
$e(\theta/4)\le sq_2$ coordinates.

Each retained input $(\alpha,\gamma)$ from $\mathcal D$ at a copy $i$ with $T_i=1$ reads
$A(\gamma,z_{i,\alpha})$, regardless of its fixed bin. We group these inputs by $(i,\alpha)$
and measure each group within its fiber, that is, by the fraction of the inputs
from $\mathcal D$ of copy $i$ with index $\alpha$ that it contains. Since the
index $\alpha$ is uniform, the $q_2$ groups of copy $i$ then have total weight $q_2d_i$,
so the groups with $T_i=1$ have total base-row weight $\ge q_2D_0/2$.
Set $K_0=2s\kappa q_2$ and $t=K_0x$. From $w^2q_2x\le1$ and $D_0\le w$,
\[
t\le\frac{2s\kappa}{w^2}\le\frac1{256w}<1.
\]
There are at most $wq_2$ groups, each of weight at most one. If $h$ groups have weight
at least $t$, the total weight of all groups is at most $h+wq_2t$, so
\[
h\ge\frac{q_2D_0}{2}-\frac{K_0}{w}
\ge K_0(128-1/w)>K_0\ge K_0/2+sq_2.
\]
After deleting the exceptional coordinates, more than $K_0/2$ blocks outside that exceptional set
remain, each of base row weight $\ge K_0x$. We retain $K_0/2+1$ of them and
all labels in $Z$. Choosing one incoming row for each retained image
$(i,\alpha,\gamma)$ gives a restriction with the required entries, since
distinct images give distinct rows. Row weights use the original $\Gamma$ weights, and the
retained column set has weight one under the fixed conditional weights on $Z$.

Since $K_0$ is dyadic, we can apply \cref{det:rareboost} with base bounds
$\RP{A}{x}{H_1}$ and integer $a=\log K_0$.
It requires $H_1+\log K_0=C+\log\kappa$ bits.
Since the protocol has depth at most $C+\log w-1$, if Bob first speaks before depth $\log w$, the chosen child has only $C+\log\kappa-2$ bits left, and a node at depth $\log w$
has only $C-1$. This proves the bound for every collection of copies.
\end{proof}

\subsection{Transfer to another copy's assigned node}
\label{sec:foreign-bound-proof}

We next bound what one copy can send to the node assigned to another, which
already retains all but the normalized fraction $\varepsilon_i$ of the
assigned copy's $\Omega$-inputs. A large transfer would yield a restriction
of the form in \cref{cl:two-stream}, which needs more bits than remain at that node.

\begin{proof}[Proof of \cref{res:foreign}]
\phantomsection\label{proof:res:foreign}
Let $\varepsilon_i$ be the normalized $\Omega$-weight missing at the node assigned to copy $i$, and put
$\delta_i=2q_2x+4\varepsilon_i$. If $\delta_i\ge1$, the conclusion is automatic:
the inputs from $\mathcal D$ in one copy have weight only $\mu\le\mu\delta_i$. Otherwise, suppose the transfer
exceeds $\mu\delta_i$. Then the node receives a set of inputs labeled by $\mathcal D$ from some copy $j\ne i$ of actual
weight $\mu\delta$ for some $\delta>\delta_i$.

Both copy indices lie outside the exceptional set for $Y'$, and only Alice has spoken before
reaching the assigned node. By \cref{det:conditioning} with error bound $\eta/2$, the projection of $Y'$ onto
coordinates $i$ and $j$ has original product weight $\ge1-\eta/2>99/100$.
We restrict to the event that both coordinates are template rows. This event
has weight $\lambda^2$, so the
retained relative weight is at least
\[
a_{\min}=\frac{99/100+\lambda^2-1}{\lambda^2}.
\]
We average over the two bin labels and fix a pair of bins attaining this relative
weight. Their retained pairs of template labels form a set $P_*\subseteq S^2$ of original
product weight $\ge a_{\min}$. The chosen bin at $i$ becomes the distinguished bin. The
projection at coordinate $i$, which lies outside the exceptional set, loses at most $\eta/2$ of its row weight
by the same error bound, while each bin has weight $\lambda/s$,
so every other
bin retains weight at least
\[
1-\frac{\eta s}{2\lambda}>2sx.
\]
We choose one surviving completion in $Y'$ for every pair in $P_*$ and
for every retained template label in the other bins. The bin label at $i$ distinguishes
the representatives of the other bins from those of the pairs, and the projected labels distinguish the
representatives within each bin. Thus the chosen completions define a
restriction of the original table with the claimed entries.

After transposition, the columns labeled by $\Omega$ from copy $i$ have normalized weight
$1-\varepsilon_i$. The columns labeled by $\mathcal D$ from $j$ have $J_2$-image weight
$\delta>\delta_i$, since that image map preserves their original weights. A pair
$(z_o,z_f)$ reads $A(\gamma_{p_i},(z_o)_{\alpha_{p_i}})$ on columns labeled by $\Omega$ and
$A(\gamma,(z_f)_\alpha)$ on columns labeled by $\mathcal D$. Other bins have the same
template entries on the first family, and their entries on the second family are immaterial. This is the
restriction in \cref{cl:two-stream}, after renaming the distinguished bin.
The claim requires $C$ bits because $\delta>2q_2x+4\varepsilon_i$, whereas the
assigned node has only $C-1$ bits left. This contradiction proves the non-strict
transfer bound, including the case $\varepsilon_i=0$.
\end{proof}

\section{Conditioning and surviving blocks}
\label{sec:conditioning-details}

We prove \cref{det:conditioning} by bounding the number of exceptional
coordinates, and then show that a short protocol must leave two or three
blocks to which \cref{lem:small-blocks} still applies. For these estimates only,
we view a label drawn according to its fixed weights as a random variable.
Expectations then express weighted sums over the fixed tuple list. The
$1024$th-moment bounds use the original weights $\nu$, under which
independence is exact, whereas the fourth-moment bound under $\nu\mid Y$ in
\cref{lem:pair-projections} uses only the approximate marginals outside the
exceptional coordinates.

\subsection{A moment count and the exceptional coordinates}
\label{sec:conditioning-moments}

We use the following moment bound twice: to bound the number of exceptional
coordinates, and, in the proof of \cref{res:incoming}, to bound the probability
that too few template positions survive. In the second use the coefficients are
the fixed fractions $d_i$ of incoming weight, so the bound applies directly to
the incoming template weight $\sum_id_iT_i$.

\begin{lemma}[Moment bound]\label{lem:conditioning-moment}
Let $r\ge1$ be an integer. Suppose $X_1,\ldots,X_m$ have mean zero,
$|X_i|\le1$, and every subfamily of at most $2r$ variables is independent.
For fixed coefficients $a_i\in[0,1]$ with $W=\sum_i a_i\ge1$,
\[
\mathbb E\left|\sum_i a_iX_i\right|^{2r}\le(r^2W)^r.
\]
Unit coefficients give the same bound with $W=m$, and the empty sum is zero.
In particular, centered indicators with $1024$-wise independence satisfy
\[
\mathbb E\left|\sum_iX_i\right|^{1024}\le(2^{18}m)^{512},
\qquad
\mathbb E\left|\sum_iX_i\right|^{256}\le(2^{14}m)^{128}.
\]
\end{lemma}
\begin{proof}
We expand the even power and group terms by which factor positions have equal
indices. A group with an index appearing once has expectation zero. Every
remaining pattern has $k\le r$ index classes, each of size at least two.
Labeling the classes in order of first appearance shows that there are at
most $r^{2r}$ such patterns. If their sizes are $n_1,\ldots,n_k$, the absolute contribution
from a fixed pattern is at most
\[
\sum_{\substack{i_1,\ldots,i_k\\\text{distinct}}}
a_{i_1}^{n_1}\cdots a_{i_k}^{n_k}\le W^k\le W^r,
\]
because $a_i^{n_j}\le a_i$. Summing over patterns proves the bound.
Taking $r=512$ and $r=128$ gives the displayed special cases.
\end{proof}

We take the exceptional coordinates to be the union of a maximal disjoint
family of small coordinate sets with large marginal error. The moment bound
shows that this family is small.

\begin{proof}[Proof of \cref{det:conditioning}]
\phantomsection\label{proof:det:conditioning}
Put $\varepsilon=\varepsilon(v,\tau)$. We choose a maximal disjoint family
$J_1,\ldots,J_m$ of sets of at most eight coordinates whose conditional
marginals differ from their product marginals by more than $\varepsilon$.
If $m=0$, we delete nothing. Otherwise, since $\nu$ has marginal $\mu^{J_i}$ on $J_i$, the total variation bound
gives an event indicator $f_i$ on each $J_i$ with
$\mathbb E_{\nu\mid Y}f_i-\mathbb E_\nu f_i>\varepsilon$ (take its complement
if necessary). Under the original weights $\nu$, any subfamily of at most $1024$
such indicators is independent, because its disjoint sets use at most
$8192$ coordinates.

For $S=\sum_i(f_i-\mathbb E_\nu f_i)$, conditional Jensen's inequality and
\cref{lem:conditioning-moment} give
\[
v(m\varepsilon)^{1024}
<\nu(Y)\mathbb E_{\nu\mid Y}|S|^{1024}
\le\mathbb E_\nu|S|^{1024}
\le(2^{18}m)^{512}.
\]
Thus $m<2^{18}\varepsilon^{-2}v^{-1/512}$.
We delete $E=\bigcup_iJ_i$, so $|E|\le8m\le e(v,\tau)$. Every remaining set of
at most eight coordinates has the desired marginal: otherwise it could be
added to the disjoint family, contradicting maximality.
\end{proof}

\subsection{Some pair has a large projection}

When many blocks survive, we apply \cref{lem:small-blocks} to a pair of
coordinates on which the retained labels project with product weight at least
$9/25$. Such a pair exists because, if every one of many disjoint pairs had a
smaller projection, the retained set itself would be too small.

\begin{lemma}\label{lem:pair-projections}
Fix disjoint coordinate pairs $J_1,\ldots,J_m$ in a tuple list with marginal
$\mu$ and label weights $\nu$. Suppose every projection of a label set $K$ has product
weight $\mu^2(\pi_{J_i}K)<9/25$.
For an exactly eight-wise independent list,
\[
\nu(K)<\frac2{m^2}<\frac9{50m}\qquad(m\ge16).
\]
For an exactly $8192$-wise independent list, let $K\subseteq Y$ with
$\nu(Y)\ge v>0$, and choose the pairs outside the exceptional set from
\cref{det:conditioning}, with any error bound $0<\tau\le1$. Then
\[
(\nu\mid Y)(K)<\frac1{4m}\qquad(m\ge16).
\]
Thus a retained set of conditional weight at least $1/(4m)$ forces one of
these pairs to have product projection weight at least $9/25$.
\end{lemma}
\begin{proof}
Write $P_i=\pi_{J_i}K$ and put
\[
p_i=\mu^2(P_i)<9/25,\qquad
X_i=\ind[z_{J_i}\in P_i]-p_i,\qquad S=\sum_{i=1}^mX_i.
\]
On $K$ every indicator equals one, so $S>(16/25)m$. Under the product weights,
the centered indicators are independent, and
$\mathbb EX_i^2,\mathbb EX_i^4\le1/4$. Only terms with no
singleton index survive the fourth-power expansion. Hence
\[
\mathbb E_{\rm prod}S^4
=\sum_i\mathbb EX_i^4+6\sum_{i<j}\mathbb EX_i^2\mathbb EX_j^2
\le\frac m4+\frac6{16}\binom m2\le\frac{m^2}4.
\]
Exact eight-wise independence gives the same fourth moment under $\nu$,
since each term uses at most four pairs. Markov's inequality gives
$\nu(K)<2/m^2$, because $1/[4(16/25)^4]<2$.
For $m\ge16$, $2/m^2\le1/(8m)<9/(50m)$.

Now suppose the pairs lie outside the exceptional set for the fixed event $Y$. Every term in $S^4$
lies in $[-1,1]$ and depends on at most eight coordinates outside that set. By
\cref{det:conditioning}, conditioning on $Y$ changes its expectation by at most $2\varepsilon$, where
$\varepsilon=\varepsilon(v,\tau)$. Summing the $m^4$ terms and applying
Markov therefore gives
\[
p:=(\nu\mid Y)(K)
\le\frac2{m^2}+\frac{2\varepsilon}{(16/25)^4}
\le\frac2{m^2}+16\varepsilon.
\]
The error term $16\varepsilon$ does not decrease with the number of pairs,
so for large $m$ we also use the original weights $\nu$. Any subfamily of at most $1024$ pair indicators involves at most $2048$
coordinates and is therefore independent under $\nu$. Since $K\subseteq Y$ and
$\nu(Y)\ge v$, \cref{lem:conditioning-moment} and Markov's inequality give
\[
p=\frac{\nu(K)}{\nu(Y)}
\le\frac{(2^{18}m)^{512}}{v((16/25)m)^{1024}}
\le\frac{(2^{20}/m)^{512}}v.
\]
Combining the bounds and using
$\min\{u,w\}\le u^{511/512}w^{1/512}$ yields
\[
\begin{aligned}
mp&\le\frac2m+
\min\{16\varepsilon m,2^{20\cdot512}v^{-1}m^{-511}\}
\\
&\le\frac2m+(16\varepsilon)^{511/512}2^{20}v^{-1/512}.
\end{aligned}
\]
The last term is at most
\[
2^{20-26\cdot511/512}v^{255/131072}<\frac1{32},
\]
by $\varepsilon\le2^{-30}v^{1/256}$ and $0<v\le1$. For $m\ge16$, this gives
$mp<1/8+1/32<1/4$, as required.
\end{proof}

\subsection{Small projections and short continuations}
\label{sec:surviving-block-proofs}

Under $\RP{M}{x}{d}$ with $d\ge2$, \cref{lem:small-blocks} gives $d+1$ bits
for two blocks of row weight $2x$ with product column weight $9/25$, and for
three blocks of row weight $4x$ with product column weight $9/50$, whenever
the row threshold is at most one. We apply it to projections as in
\cref{obs:representatives}, whose column weight $\mu^J(\pi_JK)$ is a product
weight even when the labels have conditional weights. The protocol walks in
the proofs of \cref{det:exact} and \cref{det:rareboost} both end with the next
lemma, which shows that enough surviving blocks and column weight still force
$d+1$ bits.

\begin{lemma}\label{det:suffix}
Suppose $d\ge2$ and $\RP{M}{x}{d}$, and let $b\ge0$ be an integer with
$2^{b+1}x\le1$. Consider a restriction of an interlace on a tuple list with
the base-column marginal weights that retains at least $2^b+1$ blocks, each
of base row weight at least $2^{b+1}x$, and a label set $K$. Assume one of the
following.
\begin{enumerate}
\item[(E)] The list is exactly eight-wise independent with weights $\nu$, and
$\nu(K)\ge(9/25)2^{-b}$.
\item[(C)] The list is exactly $8192$-wise independent with weights $\nu$,
$\nu(Y)\ge v>0$, the exceptional set is chosen from \cref{det:conditioning}
with error bound $0<\tau\le1$, the retained blocks have indices outside that set,
$K\subseteq Y$, and $(\nu\mid Y)(K)\ge2^{-b-1}$.
\end{enumerate}
Then the restriction requires at least $d+1$ bits.
\end{lemma}
\begin{proof}
Let $J$ be a set of at most eight retained coordinates. Under (E), exact
marginals give $\mu^J(\pi_JK)\ge\nu(K)\ge(9/25)2^{-b}$. Under (C), where
$\varepsilon=\varepsilon(v,\tau)\le2^{-30}$, \cref{det:conditioning} gives
\[
\mu^J(\pi_JK)\ge(\nu\mid Y)(K)-\varepsilon
\ge2^{-b-1}-\varepsilon.
\]
For $b\le4$ this is again $\ge(9/25)2^{-b}$, since the smallest margin
is $(7/50)2^{-4}=7/800>2^{-30}$.

For $b\le4$ we select the blocks directly. If $b=0$, two blocks meet the first
threshold above, and if $b=1$, three meet the second. If $b=2$, the
coordinates of five surviving blocks have joint product weight
$\ge9/100$, so \cref{cl:projection} gives a pair of weight
$\ge(9/100)^{2/5}>9/25$. If $b=3$ or $b=4$, eight coordinates have weight
$\ge9/400$, so some pair has weight $\ge(9/400)^{1/4}>9/25$.

For $b\ge5$, we choose $m=2^{b-1}\ge16$ disjoint pairs from the at least $2m+1$
retained blocks. If every pair projection had product weight less than
$9/25$, \cref{lem:pair-projections} would give $\nu(K)<9/(50m)=(9/25)2^{-b}$
under (E), and $(\nu\mid Y)(K)<1/(4m)=2^{-b-1}$ under (C), a contradiction.

Each selected block has row weight $\ge2^{b+1}x$, hence $\ge2x$,
and $\ge4x$ in the triple at $b=1$. These thresholds are at most one by
$2^{b+1}x\le1$. Thus \cref{lem:small-blocks} gives $d+1$ bits in every case.
\end{proof}

\begin{proof}[Proof of \cref{det:exact,det:rareboost}]
\phantomsection\label{proof:det:exact}\label{proof:det:rareboost}
In \cref{det:exact} we measure column weight under $\nu$, and in
\cref{det:rareboost} under the fixed conditional weights $\nu\mid Y$, where
the event $Y$ and its exceptional set stay fixed. Suppose a protocol for the
restriction has depth at most $d+a-1$. We follow it for $a-1$ messages by a
weight-halving walk. At each Alice message, we assign every current block to
a child retaining at least half its row weight and keep a child assigned at
least half the blocks. At each Bob message, we keep a child retaining at least
half the current column weight. If the protocol ends earlier, we pad the leaf with dummy Alice messages, which leave the
restriction unchanged. If $A$ messages are Alice's and $B$ are Bob's, then
$A+B=a-1$, at most $d$ bits remain, and the surviving restriction, with label
set $K'$, has
\[
\begin{aligned}
\text{number of blocks}
&\ge\ceil{(2^{a-1}+1)/2^A}\ge2^B+1,\\
\text{row weight per block}&\ge2^{a-A}x=2^{B+1}x,\\
\text{column weight of }K'&\ge2^{-B}\cdot\text{initial column weight}.
\end{aligned}
\]
In \cref{det:rareboost} the surviving blocks still have indices outside the
exceptional set. The initial column weight is $\nu(K)\ge9/25$ in
\cref{det:exact} and $(\nu\mid Y)(K)\ge1/2$ in \cref{det:rareboost}, so $K'$
satisfies hypothesis (E), respectively (C), of \cref{det:suffix} with $b=B$.
Its row threshold is at most one because $B+1\le a$. It requires $d+1$
further bits, while only $d$ remain.
\end{proof}

\section{Constructing rational input lists}
\label{sec:support-algorithm}

The deterministic reduction prints its tuple lists together with their
rational weights, so \cref{det:code} must bound the encoding length and not
only the number of tuples. We use the sequential construction of
\citet[Section~3]{koller1994sample}: extend each word by every coordinate
value with its product weight, then reduce the list while preserving the
joint weights on all sets of at most $\mathsf D$ coordinates. We recall the
procedure and bound the bit length of the weights.

\begin{proof}[Proof of \cref{det:code}]
\phantomsection\label{proof:det:code}
Fix the independence order $\mathsf D$ and put $P=(1+qm)^{\mathsf D}$, which
bounds the number of constraints below. We start
with the empty word of weight one. At each step, we append every possible
coordinate value $y$ to every retained word and multiply its weight by $\mu(y)$. On the resulting candidate
set $\mathcal U\subseteq Y^i$, the weights satisfy, by induction on $i$,
\[
\sum_{u\in\mathcal U}w_u\ind[u_J=a]
=\prod_{j\in J}\mu(a_j)
\qquad
(J\subseteq[i],\ |J|\le\mathsf D,\ a\in Y^J).
\tag{*}\label{eq:support-cylinders}
\]
There are at most $P$ equalities in \eqref{eq:support-cylinders}, including
the normalization equation for $J=\varnothing$. The cited procedure removes
linear dependences among the zero--one columns of this system until at most
$P$ words remain, and it preserves the equalities and the nonnegativity of the
weights. Each stage therefore extends at most $P$ words to at most $mP$ candidates.
After $q$ stages the equalities give the required $\mathsf D$-wise independence.

Write $a(u)$ for the constraint column of word $u$. A reduction step takes
an integer dependence $\sum_uc_ua(u)=0$ and sets $w'_u=w_u-tc_u$, where
$t=\min_{c_u>0}w_u/c_u$, which removes at least one word of positive weight.
To find the dependence, we take a basis of $r\le P$ constraint columns, one
further column, and $r$ rows on which the basis columns are independent. The
determinant minors of this $r\times(r+1)$ zero--one matrix give an integer
relation valid in every row, with coefficients at most $\max\{1,r^{r/2}\}$
by Hadamard's bound, and hence with $O(P\log(P+1))$ bits. Exact linear
algebra finds this relation in polynomial time, and we break all ties by fixed orderings.

Let $L_\mu$ be the total binary encoding length of the weights $\mu$. Extending a word
multiplies the weights' common denominator by at most the product of the
denominators in $\mu$, adding at most $L_\mu$ bits. An elimination multiplies
it by the positive coefficient attaining the minimum ratio, adding
$O(P\log(P+1))$ bits. Across $q$ extensions and at most $qmP$ eliminations,
every denominator therefore has bit length
\[
O\bigl(qL_\mu+qmP^2\log(P+1)\bigr).
\]
Numerators obey the same bound because the weights lie in $[0,1]$. Hence
all arithmetic has polynomial cost for fixed $\mathsf D$, and so does the
output of at most $P$ weighted words of length $q$.
\end{proof}

\section{Printing a sample of the implicit table}\label{sec:sampling-algorithm}
\label{app:implementation}

The randomized reduction uses the implicit table only through sampled rows,
columns and their entries, which we generate without printing the table. We
first bound the length and generation cost of complete input labels, then
implement their elementary draws with fixed blocks of fair bits. Evaluating
one entry reads only one base entry per recursion level, but its labels must
be generated in full, including the coordinates that the evaluation never reads.

\begin{proof}[Proof of \cref{lem:sample-access}]
\phantomsection\label{proof:lem:sample-access}
We fix a source graph on $n$ vertices above the size cutoff in \cref{lem:full-schedule}, put $L=\ceil{\log n}$, and use the parameter choices of that lemma at
every level. Write $R:X\times Y\to\{0,1\}$ for the current base matrix.

\emph{Labels and entries.}
The two interlaces $U=\Int{w}{N}$, $W=U^{\mathsf T}$ and
$R^+=(\Int{Q}{W})^{\mathsf T}$ give the row and column labels of $R^+$ (the successor labels)
\[
\mathbf a=((i_t,Z_t))_{t\in[Q]},\qquad
\mathbf b=(j,(k_i)_{i\in[w]}),
\]
where $i_t\in[w]$ and $j\in[Q]$ are uniform, $Z_t$ are tagged columns of $B_{\rm raw}$,
and $k_i\in\mathcal K$ are rows of $B_{\rm raw}$. We draw all these coordinates independently, each according to its weight
in \cref{def:interlace,def:source-matrix}. The entry rule is
\[
R^+(\mathbf a,\mathbf b)=B_{\rm raw}(k_{i_j},Z_j).
\]
The entry rules in \cref{def:source-matrix} select exactly one entry of
$R$. To generate a row $k_i$ or column $Z_t$ of $B_{\rm raw}$, we sample according to the row-type and column-family weights from that
definition and draw every listed coordinate, including those not visited by
the eventual entry evaluation. The family $\mathcal D$ copies one tuple
and its tail, so it does not require extra draws for the copied positions.
We generate only actual input coordinates, not the auxiliary sequences used
only in the lower-bound proof. The recursive draw counts are
\[
\begin{array}{c|c|c}
\text{label}&\text{base-row draws}&\text{base-column draws}\\ \hline
\text{template row}&q_2&0\\
\text{row for an item}&1&0\\
\Omega\text{ column}&0&(s+1)+s(2s-1)=2s^2+1\\
\mathcal D\text{ column}&0&2s-1\\
\text{successor row}&0&\le Q(2s^2+1)\\
\text{successor column}&\le wq_2&0
\end{array}
\]
Apart from these recursive draws, generating indices and scanning activity
lists cost polynomial time in $n+s+q_2+w+Q+L+1$, copying a label costs time
linear in its length, and no coloring is required. Every width $s,q_1,q_2,w,Q$ is $2^{O(L)}$ by
\cref{lem:full-schedule}, and so is every draw count in the table. Hence we can choose common computable integer bounds $J_i$ for
the label length, generation time, elementary draw count and entry-evaluation
time at level $i$ such that
\[
J_0=O(L),\qquad J_{i+1}\le2^{c_0L}J_i,\qquad
J_k=2^{O((k+1)L)},
\]
where $c_0$ is fixed and effective. Here each elementary draw counts as one
operation, and we implement it with fair bits below. At the initial matrix $R_0$ of \cref{cl:seed},
a row consists of $p_0=O(L)$ independent uniform bits, a column is a uniform index in
$[p_0]$, and an entry can be read in $O(L)$ time. This gives $J_0=O(L)$.

Since $h=O(L)$ and $d_0=O(L)$, we have $B_k=O((k+1)L)$. The direct sample
size in \cref{thm:compression} then gives $m=2^{O((k+1)L)}$.
We generate and store the $m$ row labels and $m$ column labels, then evaluate all
$m^2$ row--column entries.
Including the stored tuples and printed entries, this costs
\[
O((2m+m^2)J_k)=2^{O((k+1)L)}=n^{O(k+1)}.
\]
All constants are independent of $k$.

\emph{Fixed-bit implementation.}
Let $N_{\rm draw}\ge1$ be a computable bound on the elementary draws for
all $2m$ labels, so $N_{\rm draw}=2^{O((k+1)L)}$. On every tape, each draw
has at most
\[
A=\max\{512,4s(n+s),q_2,w,Q,p_0\}=2^{O(L)}
\]
outcomes, and each rational mixture has a common denominator $\le A$.
This holds for type draws (whose weights in \cref{def:source-matrix} have
denominators $512$ and $256$), item indices, bins, activity indices,
interlace-copy indices and draws from the initial matrix. Base rows and columns are never
drawn directly. We generate them recursively down to the
initial matrix. We set $b=\ceil{\log(200A N_{\rm draw})}$ and reserve
$N_{\rm draw}$ blocks of $b$ fair bits. A draw from $[a]$ reads the next
block as a uniform integer $U\in\{0,\ldots,2^b-1\}$ and returns
\begin{equation}\label{eq:rounded-draw}
\operatorname{Draw}(a)=1+\floor{aU/2^b}.
\end{equation}
A rational mixture with common denominator $a\le A$ maps consecutive
intervals of values of $\operatorname{Draw}(a)$, whose lengths are the
numerators of its weights, to its types. Unused blocks are ignored.

Each value in $[a]$ has $\floor{2^b/a}$ or $\ceil{2^b/a}$ preimages, so a
draw is within total variation $a/2^b\le1/(200N_{\rm draw})$ of the ideal
one, and mapping values to types cannot increase this distance. Because every
draw reads a fresh block, the bound holds even when the alphabet depends on
earlier outcomes.
We couple the implemented and ideal draws one at a time. By a union bound
over the $N_{\rm draw}$ draws, some pair disagrees with probability at most
$N_{\rm draw}/(200N_{\rm draw})=1/200<1/100$.

Every tape thus produces exactly $m$ legal row labels and $m$ legal column
labels, and all $m^2$ entries are evaluated exactly. The tape has fixed length
$bN_{\rm draw}$, where $b=O((k+1)L)$. Arithmetic on $O(b+\log A)$-bit integers adds only
polynomial overhead. Hence the running time and the number of random bits are
$n^{O(k+1)}$ on every tape.
\end{proof}

\section{One parameter choice for both constructions}
\label{sec:parameters}

We verify the numerical requirements of the recipe in
\cref{def:amplification-parameters}: first those shared by both
constructions, then, for each choice of tuple weights, the source-matrix
deletion bounds and the final interlace.

Recall $n=|V|$, $\lambda=255/256$, $\mu=1/256$, and $\zeta=\gamma/24$.
The source promise gives
\[
 n/2\le |E|\le\Delta n/2,\qquad
 \delta_v=\frac{\gamma|E|}{2\Delta n}\ge\frac{\gamma}{4\Delta}.
\]
From \cref{res:destinations} we use the bound $E_{\rm exc}$ on the sum of
$\kappa$ over stopping nodes with no assigned copy, and from the proof of
\cref{res:incoming} the weighted moment constant $C_{128}$:
\begin{equation}\label{eq:common-budgets}
E_{\rm exc}=2e(\theta,\eta/2),\qquad
C_{128}=\left(\frac{2^{14}}{(\lambda-1/2)^2}\right)^{128}.
\end{equation}
The definition
of $e(v,\tau)$ before \cref{det:conditioning} gives the estimate
\begin{equation}\label{eq:conditioning-size-bound}
\begin{aligned}
e(v,\tau)
&\le8+2^{23}\tau^{-2}v^{-1/512}+2^{83}v^{-5/512}\\
&\le2^{84}\tau^{-2}v^{-1/64}
\qquad(0<v,\tau\le1).
\end{aligned}
\end{equation}

\begin{lemma}[Common numerical bounds]\label{lem:common-schedule}
There is an effective dyadic constant $s_0$, depending only on the fixed
graph promise, such that the recipe in \cref{def:amplification-parameters}
with $n\le\rho^{-1}$, $\rho\le s^{-4}$, and $s\ge s_0$ satisfies all source
and internal-interlace bounds listed in \eqref{eq:width-margins} below. The choice of
$s_0$ is independent of the base matrix and the number of rounds.
\end{lemma}

\begin{proof}
The middle column of the table below lists upper bounds obtained by
substituting the recipe. Each tends to zero as $s$ increases. The last column
lists the bound that the proof requires for each quantity.
\begin{equation}\label{eq:width-margins}
\begin{array}{c|c|c}
\text{quantity}&\text{upper bound}&\text{required allowance}\\ \hline
s\eta&s^{-1}&2^{-16}\min\{1,\delta_v,\mu\zeta\}\\
4s|E|/q_2&2\Delta/s&1\\
s^2\rho&s^{-2}&1/4\\
sq_1/q_2&s^{-4}&\zeta/128\\
\max\{q_2x,sx,wx\}&s^{-24}&10^{-4}\\
w^2x/\eta&s^{-14}&1/(32\lambda)\\
w^2q_2x&s^{-10}&1\\
wq_2x/\eta&s^{-16}&1/(128\mu)\\
e(\theta/4)/(sq_2)&2^{85}s^{-53/8}&1\\
C_{128}s^{-128}/\theta&512C_{128}s^{-104}&1/4\\
2sE_{\rm exc}/(\eta w)&2^{89}s^{-5/8}&1/16
\end{array}
\end{equation}
For the two rows involving $e$, \eqref{eq:conditioning-size-bound} with
$\theta=2^{-9}s^{-24}$ and $\eta=s^{-2}$ gives
$e(\theta/4)\le2^{85}s^{3/8}$ and
$E_{\rm exc}\le2^{88}s^{35/8}$, while $sq_2=s^3/\rho\ge s^7$.
The other rows follow by substituting the widths of
\cref{def:amplification-parameters} and using $|E|\le\Delta n/2$, $n\le\rho^{-1}$,
and $\rho\le s^{-4}$.
We choose a dyadic $s_0\ge256$ so that every middle entry is below its allowance,
using $\gamma/(4\Delta)$ in place of $\delta_v$. Clearing the fixed rational
exponents makes this a finite set of integer-power comparisons. Their
left sides decrease with $s$, so the same choice works for every $s\ge s_0$.
None of these comparisons involves the base matrix or the number of rounds.

The table's second row gives $q_2\ge4s|E|$. Since $q_2$ is a power of two,
$q_2\ge q_0=\operatorname{pow}_2(4s|E|)$, as \cref{def:source-matrix} requires. In particular,
$q_2\ge2sn\ge2(n+s-3)$, so all item labels fit: $|\mathcal I|=4s(n+s-3)\le2sq_2$, as
\cref{lem:amplifier-rank} requires. The table's last bound
also implies the two bounds used to count assigned nodes:
\[
E_{\rm exc}\le\frac{\eta w}{32s}<\frac w4,\qquad
e(\theta)\le e(\theta,\eta/2)=\frac{E_{\rm exc}}2
\le\frac{\eta w}{64s}<\frac w{128}.
\]
Here $e(\theta)=e(\theta,1)\le e(\theta,\eta/2)$ because $e(v,\tau)$ is
nonincreasing in the error bound $\tau$.

The remaining fixed-fraction checks appear in the proofs that use them.
For each, we list the table rows it uses and derive here only the bounds that
the proof does not display.
\begin{itemize}
\item \emph{Block counts} (proof of \cref{src:profiles}). Since $s\ge256$
and $s\rho<10^{-3}$, the three lower bounds for $m/u$ satisfy
\[
\frac35(1-2/s)-q_1\rho>\frac9{16},\qquad
\frac35-q_2x>\frac9{16},\qquad
\frac{44}{75}-sx>\frac9{16},
\]
the last since $2^j(\beta_j-(1-\lambda))/\lambda\ge((3/5)-4(1-\lambda))/\lambda=44/75$
for $j\in\{0,1\}$. The bound on $m/\kappa$ in the
proof of \cref{res:alignment} uses $wx\le1/256$ and $4e(\theta)/w<1/32$
from the display above.
\item \emph{Item deletion and collisions} (proof of \cref{src:near-full}).
The row $s\eta$ gives $4096\eta sn\le\delta_v n/16$ and
$4096\eta s(s-3)\le(s-3)/16<s/2$. The rows $s\eta$ and $sq_1/q_2$ give
$8sq_1+4s\eta q_2/\mu\le(1/16+2^{-14})\zeta q_2<\zeta q_2/4$.
\item \emph{Row and template thresholds.} The rows $s^2\rho$, $s\eta$
and $\max\{q_2x,sx,wx\}$ give $2s\rho\le1/8$, $2wx+\eta/(16\lambda)<1$,
$2sx\le10^{-3}$, and hence $999/1000-2(s-1)x>998/1000$ in the proof of
\cref{src:near-full} and
$(99/100+\lambda^2-1)/\lambda^2-10^{-3}>49/64$ in \cref{cl:two-stream}.
\item \emph{Conservation} (proof of \cref{res:alignment}). That proof uses
the rows $w^2x/\eta$, $wq_2x/\eta$ and $2sE_{\rm exc}/(\eta w)$ and displays
the resulting bounds.
\item \emph{Moment bound} (proof of \cref{res:incoming}). Since $D_0\ge512s$,
the probability bound $C_{128}D_0^{-128}$ from \cref{lem:conditioning-moment}
satisfies $C_{128}D_0^{-128}\le C_{128}s^{-128}\le\theta/4$ by the row
$C_{128}s^{-128}/\theta$.
\end{itemize}
\end{proof}

\subsection{Full products}

In the randomized reduction, the graph-encoding width $q_2=s^22^L$ is the only
width that contributes a term linear in $L$ to the per-round overhead $h$. The widths
$s$, $w=s^8$, and $Q=s^{24}$ are polynomial in $\log n$. A smaller overhead allows more
amplification rounds before the final sample becomes too large.

\begin{lemma}[Parameters for full products]\label{lem:full-schedule}
Put $L=\ceil{\log n}$, $\rho=2^{-L}$, and
$s=\operatorname{pow}_2(2^{16}(L+1))$, and use the recipe of \cref{def:amplification-parameters}.
For every sufficiently large promised source, these choices satisfy the hypotheses of
\cref{lem:common-schedule}, the numerical conditions for the full-product deletion bounds for the source matrices, and
the numerical hypotheses of \cref{cl:bridge} with $b=L$ and $\delta=\delta_U$. The per-round overhead is
\[
h=2\log s+\log q_2+\log w+\log Q=L+36\log s=L+O(\log L).
\]
There is an effective fixed cutoff beyond which also
$L\ge16$, $d_0\le L$, and $h\le7L$, where $d_0$ is the cost of the initial matrix in
\cref{cl:seed}. Computing the widths has polynomial bit complexity and
uses neither the base matrix nor the number of rounds.
\end{lemma}

\begin{proof}
Since $s=\operatorname{pow}_2(2^{16}(L+1))$, we have $2^{16}(L+1)\le s<2^{17}(L+1)$. Hence $s\ge s_0$ and
$2^{-L}\le s^{-4}$ beyond an effective fixed cutoff, so the common lemma
applies. In each application of the deletion bounds for the source matrices, the block parameter
satisfies $u\ge s/2$. Since $\log(1/x)=8L$,
\[
u\ge2^{15}(L+1)>4096\log(1/x),
\]
which gives the full-product block threshold
$x\ge2^{-u/1024}$. The common numerical bounds give the required block counts
and row-weight bounds. For the final interlace the remaining checks are
\[
0<\delta_U=\frac1{16s^{10}}<\frac1{64},\qquad
Q\theta=\frac1{512}<\frac1{20},\qquad
Q\delta_U=\frac{s^{14}}{16}\ge80L.
\]
Thus \cref{cl:bridge} applies with $b=L$.

Substituting the widths gives the displayed value of $h$. Writing
$\ell=\ceil{\log(L+1)}$, the initial matrix of \cref{cl:seed} has $d_0=\log p_0+1=\ell+7$ and
$\log s=16+\ell$, so $d_0\le L$ and $h\le7L$ eventually.
Since $s=\Theta(L+1)$ with explicit constants, the ratios $s^42^{-L}$,
$d_0/L$ and $(h-L)/L$ are at most fixed multiples of $(L+1)^42^{-L}$ or
$(1+\log(L+1))/L$. These bounds decrease to zero beyond a computable point.
Hence the cutoff is effective and valid at every larger size. All widths except $q_1=s-2$ are powers of two
with exponents $O(L)$; forming those powers and subtracting two from $s$ gives the stated bit complexity.
\end{proof}

\subsection{Explicit lists}

For the deterministic reduction, we choose $s$ with $s^4\ge n$
and set $\rho=s^{-4}$. Then every width is polynomial in $s$,
including the final-interlace width $Q=s^{24}$ required by \cref{det:restore}.

\begin{lemma}[Parameters for explicit lists]\label{lem:generated-schedule}
After increasing the effective dyadic constant $s_0$ if necessary, every
dyadic $s\ge\max\{n^{1/4},s_0\}$ satisfies the bounds below. Choose the least such $s$ and put
$\mathsf N=s^4$, $L_{\rm g}=\log\mathsf N$, and $\rho=\mathsf N^{-1}$.
With these choices, the common recipe satisfies all numerical hypotheses of the explicit-list
construction and of the final interlace in \cref{det:restore}. Its overhead is
\[
h_{\rm g}=2\log s+\log q_2+\log w+\log Q
=40\log s=10L_{\rm g}.
\]
With $\mathsf N_0=s_0^4$, the padding satisfies
$n\le\mathsf N<16\max\{n,\mathsf N_0\}$. Every width is polynomial in
$\mathsf N$. One round has a fixed polynomial cost in $\mathsf N$ and the
complete base encoding, with exponent independent of the round number.
\end{lemma}

\begin{proof}
Here $\rho=s^{-4}$, so the common lemma applies. Only the deletion bounds for the source matrices
at tuple weight $x$ and the final interlace need additional checks.
By \eqref{eq:conditioning-size-bound},
\[
e(x)\le2^{84}s^{1/2}\le s/128,\qquad
e(\rho,\delta_U/2)\le2^{94}s^{20+1/16}.
\]
We increase $s_0$ so that $2^{84}s^{1/2}\le s/128$ and
$128(2^{94}s^{20+1/16}+1)\le s^{24}=Q$. Dividing each condition by its right
side leaves fixed multiples of negative powers of $s$, so this enlargement is
effective and valid for every larger $s$.

For the deletion bounds for the source matrices, $u\ge s/2$, and by the three
block-count bounds in the proof of \cref{src:profiles}, more than $9u/16$
blocks meet the required row-weight threshold. Removing at most $e(x)\le s/128\le u/32$ exceptional
coordinates leaves more than $9u/16-u/32>u/2$ blocks. The row thresholds obey
$u\rho\le s\rho\le1$ for the first matrix and
$ux\le\max\{q_2x,sx\}\le1$ in the other two.
For the final interlace, $0<\rho\le1$, $0<\delta_U<1/64$,
$Q\theta=1/512<1/20$, and the bound $Q\ge128(e(\rho,\delta_U/2)+1)$ above are
the numerical hypotheses of \cref{det:restore} with $\phi=\rho$ and
$\delta=\delta_U$.

The least dyadic choice satisfies
$s<2\max\{n^{1/4},s_0\}$, proving the factor-$16$ size bound. It can be
found by comparing integer fourth powers, without approximating a root.
The widths are $q_2=s^6$, $w=s^8$, and $Q=s^{24}$, giving the stated
overhead. The independence order stays fixed at $8192$, so by
\cref{det:code} the fixed number of tuple-list constructions have cost
polynomial in these widths and in the complete rational encoding of the base matrix.
Neither the parameter choices nor the exponent of the per-round polynomial
cost depends on the number of rounds.
\end{proof}

\section{Coloring promises from SAT}
\label{sec:sources}

The two communication reductions use different coloring problems: bounded-degree
graph coloring for interlacing (\cref{src:gap-colouring}), and balanced
hypergraph coloring for strict alternation (\cref{yao:source}). This appendix derives both hardness results from the literature and checks
the input conventions that the communication constructions need.

\subsection{Bounded-degree graph coloring of almost-linear size}
\label{app:graph-source}

The graph source of \cref{src:gap-colouring} needs a constant coloring gap
and must keep the almost-linear size of sparse SAT. Composing the reductions of
Dinur and of Bogdanov, Obata and Trevisan gives both. We record this composition
and our graph conventions because the size bound determines the final additive-error bound.

\begin{proof}[Proof of \cref{src:gap-colouring}]
\phantomsection\label{proof:src:gap-colouring}
For a SAT instance of encoding length $L$, Dinur's almost-linear PCP
construction gives a constant-alphabet binary constraint graph of size
$L(\log L)^{O(1)}$, with perfect completeness and a constant
unsatisfied-constraint fraction in the NO case
\cite[Lemma~8.3 and proof of Theorem~8.1]{dinur2007pcp}.
Her degree reduction makes the degree constant, preserves perfect
completeness, and changes the size and the NO-case unsatisfied fraction by
only constant factors~\cite[Lemma~1.7]{dinur2007pcp}.
A constant-size Boolean encoding, followed by
the reduction of \citet[Lemma~11]{bogdanov2002lower}, gives bounded-occurrence $3$-SAT
with perfect completeness, a constant NO-case gap and almost-linear size. The gap-preserving reduction of
\citet[Theorem~14 and its proof, pp.~98--99]{bogdanov2002lower}
then gives a bounded-degree graph of linear size.

It remains to match the conventions of \cref{src:gap-colouring}. We remove isolated vertices and add a
disjoint edge. This preserves completeness and, when the graph already has
an edge, decreases the NO-case gap by at most a factor of two. For trivial inputs we use fixed promise graphs.
For a sparse formula on $N$ variables, $L=O(N\log(N+2))$, so the resulting
graph $H$ has $N(\log(N+2))^{O(1)}$ vertices and edges. Taking
$\max\{1,\ceil{N/|E(H)|}\}$ disjoint copies gives
$N\le |E(G)|\le N+|E(H)|$ without changing the degree or gap.
The remaining edge bounds follow from the absence of isolated vertices
and the constant maximum degree. Parallel edges are retained, and edge
occurrences are counted with multiplicity. We choose a rational $\gamma>0$
below the resulting gap and enlarge the fixed exponent $p$ to cover the
constant factors. The cited reductions and these adjustments are
deterministic and polynomial-time.
\end{proof}

\subsection{Balanced hypergraph coloring with distinct vertices}
\label{app:hypergraph-source}

The hypergraph promise of \cref{yao:source} comes from the hardness result of
Guruswami and Lee~\cite{guruswami2018balanced}, whose edges may contain repeated
vertices. We replace repeated occurrences by distinct labeled clones and
check that the YES and NO conditions of \cref{yao:source} are preserved.

\begin{proof}[Proof of \cref{yao:source}]
\phantomsection\label{proof:yao:source}
Recall the fixed parameters of \cref{yao:source}
\[
s_{\rm a}=2^t,\qquad K_{\rm a}=s_{\rm a}^2,\qquad
S_{\rm a}=2^{s_{\rm a}}.
\]
By Theorem~1.1 of \citet{guruswami2018balanced}, it is NP-hard to distinguish
$K_{\rm a}$-uniform occurrence hypergraphs admitting an $s_{\rm a}$-coloring
with at least $s_{\rm a}-1$ occurrences of each color in every edge from those
with no independent set containing at least a $1/(2S_{\rm a})$ fraction of the
vertices. Here, as in their reduction, an edge of an occurrence hypergraph is a
family of $K_{\rm a}$ vertex occurrences indexed by $[s_{\rm a}]\times[s_{\rm a}]$,
listed in a fixed order. A vertex may occur more than once in an edge, and
colors in an edge are counted with multiplicity, as in their completeness
proof. Hypergraphs of the second kind admit no proper $S_{\rm a}$-coloring,
since some color class would contain at least a $1/S_{\rm a}$ fraction of the
vertices and so would not be independent.

We replace each vertex $u$ by $K_{\rm a}S_{\rm a}$ labeled clones. For each
ordered occurrence edge $(u_1,\ldots,u_{K_{\rm a}})$, we include every tuple
\[
(u_1^{(j_1)},\ldots,u_{K_{\rm a}}^{(j_{K_{\rm a}})})
\]
of pairwise distinct clones, in the displayed order. Retain repeated edge
occurrences as separately labeled edges. Each original edge produces at least
one such tuple, since each vertex has $K_{\rm a}S_{\rm a}\ge K_{\rm a}$ clones. Giving clones
their source vertex's color preserves the color multiset of every edge,
and hence the YES promise.

Conversely, suppose the cloned hypergraph has a proper $S_{\rm a}$-coloring.
We give each source vertex a color appearing on at least $K_{\rm a}$ of its
clones. Suppose some occurrence edge were monochromatic in this induced coloring.
Each source vertex occurs at most $K_{\rm a}$ times in that edge, so we can choose
different clones of that color for its positions. This gives a monochromatic
edge of distinct clones, contradicting properness. Hence the induced coloring
is proper, contradicting the NO promise.

Each vertex produces $K_{\rm a}S_{\rm a}$ clones and each edge produces at
most $(K_{\rm a}S_{\rm a})^{K_{\rm a}}$ cloned edges. These are constants for fixed
$t$, so enumeration is deterministic and polynomial-time, although the degree
may depend on $t$. If the vertex set is empty, we add one isolated labeled
vertex, since \cref{yao:source} requires a nonempty vertex set. Because each
ordered edge has distinct vertices, the position $\iota_\alpha(v)$ in
\cref{def:alternating-table} is well defined.
\end{proof}

\section{Finite-field arrays}\label{sec:arrays}
\label{sec:alternating_implementation}

The alternating lower bound needs a short, explicitly generated list of
arrays on which no prescribed bit pattern occurs much more often than it
would for independent fair bits. We construct that list here, completing
the algorithmic step promised in \cref{yao:array-list}.

We use the polynomial-root construction of Alon, Goldreich,
H{\aa}stad and Peralta~\cite[Construction~3 and Proposition~3]{alon1992smallbias}.
The conversion from parity bounds to pattern frequencies is also recorded
there, with credit to Vazirani~\cite[Lemma~1 and Corollary~1]{alon1992smallbias}.
We use the following direct consequence and keep its seed labels explicit
so that the reduction can print the full array list.

\begin{lemma}[Explicit nearly uniform arrays]\label{yao:finite-field-list}
For integers $n,r\ge1$, a deterministic algorithm prints a nonempty labeled
list $(w^\ell)_{\ell\in\mathcal L}$ of $n$-bit arrays satisfying
\[
\frac{|\{\ell\in\mathcal L:w^\ell|_I=z\}|}{|\mathcal L|}
\le2^{-|I|}+2^{-r}
\qquad(I\subseteq\{0,\ldots,n-1\},\ z\in\{0,1\}^I).
\]
Its running time and complete output length are polynomial in $n2^r$.
\end{lemma}

To instantiate the construction, let $R$ be the least power of two with
$R\ge n2^r$. The construction lists $R^2<4n^24^r$ arrays, indexed by
$\mathbb F_2^{\log R}\times\mathbb F_R$, whose nonempty parities have bias at most
$(n-1)/R\le2^{-r}$, and the bias-to-pattern conversion gives the displayed bound.
A field representation of $\mathbb F_R$ can be found by exhaustive search in time
polynomial in $R$, and we keep repeated arrays as distinct labels.

Finally, we apply \cref{yao:finite-field-list} on the positions
$\Omega_{\rm a}$ of \cref{def:alternating-table} with error $\epsilon_{\rm a}$,
and check that the list has size polynomial in $H_{\rm a}$ for each fixed $t$.

\begin{proof}[Proof of \cref{yao:array-list}]
\phantomsection\label{proof:yao:array-list}
We number $\Omega_{\rm a}$ by $0,\ldots,n-1$ and apply
\cref{yao:finite-field-list} with
\[
n=K_{\rm a}H_{\rm a},\qquad r=4(b+3t).
\]
Here $2^{-r}=\epsilon_{\rm a}$ and
$n2^r=K_{\rm a}2^{12t}H_{\rm a}^5$, which is polynomial in $H_{\rm a}$
for fixed $t$. The integers $n$ and $r$ have $O_t(b+1)$ bits, and the
polynomial bound also covers $H_{\rm a}=1$.
\end{proof}

\section*{AI methodology}
\label{sec:ai-methodology}

The authors' initial goal was to sparsify interlacing while retaining the
hardness needed for the reductions. That research direction preceded the
AI-assisted development of the detailed constructions and proofs.

The intermediate combined development already sparsified interlaces using
balanced families with controlled projections. The present work distinguishes
three constructions: implicit full products with one final random compression,
explicit input lists that preserve the conditions for repeated amplification,
and a separate small-bias construction for strict alternation.

ChatGPT contributed substantive proof proposals for full-product amplification
followed by a single final compression in the randomized construction. ChatGPT later
supplied proposals for the bounded-independence and conditioning arguments in
the unrestricted deterministic construction. ChatGPT also proposed the separate
small-bias argument for strictly alternating protocols.

Codex assisted with reconstructing and integrating the proposals, and with
mathematical review of the arguments and their interfaces. Its work included
checking implementation details and parameter accounting, running finite
verification scripts, and drafting and revising the exposition. The affected
material included the randomized amplification and final compression, the
unrestricted deterministic construction, and the strict-alternation construction
and its supporting appendices, together with exposition throughout the
manuscript. Codex also assisted with the encoding-first proof overview and
the terminology revision in the present version. Claude supplied editorial review.

The authors retain responsibility for the mathematical claims and their
attribution. AI was used in proof development and review, not solely for
copy-editing. AI-assisted checks, agreement between model reviews, and finite
experiments are neither formal verification nor independent human review.





%% ---- appendix: interlacing_details ----

\bibliographystyle{plainnat}
\bibliography{references}
\end{document}
